\documentclass[french]{book}

\usepackage[utf8]{inputenc}
\usepackage[T1]{fontenc}
\usepackage{lmodern}
\usepackage{geometry}
\usepackage{graphicx}
\usepackage{babel}
\usepackage{amsmath}
\usepackage{amsthm}
\usepackage{amssymb}
\usepackage{hyperref}
\usepackage{stmaryrd}
\usepackage{enumitem}
\usepackage{tcolorbox}
\usepackage{dsfont}
\usepackage{multirow}
\usepackage{etoc}
\usepackage{titlesec}
\usepackage{esint}
\usepackage{cancel}
\usepackage{mathdots}

\setlist[itemize]{label=$\bullet$}

\theoremstyle{definition}
\newtheorem*{dfn}{Definition}
\newtheorem*{refe}{Référence}
\newtheorem*{rep}{Repère philosophique}
\newtheorem*{bio}{Biographie}

\begin{document}

\author{EPL}
\title{Philosophie}


\frontmatter
\maketitle
\tableofcontents

\mainmatter

\chapter{Introduction à la philosophie}
\section{Qu'est ce que la philosophie ?}

Une tentative de décrire le monde de manière objective. Un mode de réflexion et mode de vie. Un moyen de réponse (ou tentative de réponse) aux questions existentielles ("Qui suis-je ? Pourquoi suis-je ? Etc.").

\section{Étymologie et signification}

De \textit{philo} en grec, qui veut dire j'aime (de \textit{philein} = aimer) et de \textit{sophia}, la sagesse en grec, la philosophie est étymologiquement l'amour de la sagesse.

Cependant, l'amour revêt 3 façades : \begin{itemize}
	\item la passion pure (\textit{philein}) qui va nous intéresser ; 
	\item la passion érotique (\textit{Eros}) ;
	\item la charité (\textit{Agape}).
\end{itemize}

$\implies$ Ainsi, le philosophe aime la sagesse, et par conséquent, il ne la possède pas et veut tendre vers elle.

Aussi : sagesse = mode de vie (pratique) + connaissance/savoir (théorique)
On se concentrera donc bien évidemment sur ce second aspect en cours de philosophie.

Ainsi, la philosophie revêt deux dimensions : dans sa dimension théorique, elle est la recherche de la Connaissance, du Vrai, et du Bien. Et dans sa dimension pratique, elle est la recherche du mieux ou bien vivre, des principaux moraux d'action et du comportement en vue de bien agir.

\begin{dfn}[Philosophie]
	Manière spécifique de réfléchir en utilisant sa raison dans le but de connaître, d'atteindre la Vérité et de bien vivre
\end{dfn}


$\implies$ Les réponses à ces questions sont donc complexes (d'où la dissertation) et variées.
$\implies$ Le plus important n'est pas le résultat auquel on aboutira (la réponse à la question posée) mais la démarche, la manière, le processus de réflexion mis en œuvre pour arriver à une réponse (ou non d'ailleurs ).

IL NE FAUDRA DONC NÉGLIGER AUCUNE PISTE DE RÉFLEXION !!!

\begin{dfn}[Philosopher]
	Se questionner, refuser la réponse facile et argumenter
\end{dfn}



Repères historiques :

La philosophie naît en Grèce au VIIe siècle avant JC.
On peut dès lors découper l'histoire de la philosophie antique en 3 grandes parties :

La période présocratique :
Les premiers sont des philosophes dits "de la nature". En effet, ces derniers se consacrent particulièrement aux questions concernant le monde et son origine : comment est-il organisé ? comment a-t-il été créé ? … Ce sont les premiers philosophes à raisonner différemment du mythe.

\begin{dfn}[Mythe] 
Récit de fiction ayant pour but d'expliquer le monde, et plus particulièrement son origine. Les protagonistes sont généralement des dieux, des demi-dieux ou des êtres surnaturels.
\end{dfn}


Un de ces philosophes est Thalès de Milet, qui fonda l'École de Milet. Tout comme d'autres philosophes de la nature, Thalès pensait qu'un principe unique (pour lui l'eau, ce pouvaient être autre chose pour les autres, ou alors plusieurs choses) était à l'origine du monde. C'est ce qu'on appelle la philosophie ionienne.

À noter qu'on nomme aussi ces philosophes "présocratiques" car il sont situés temporellement avant Socrate, le premier philosophe qui se consacre à autre chose que le monde et son origine.

La période grecque classique


La période hellénistique 

\chapter{La vérité est-elle libératrice ?}
(Vérité et Liberté)

\begin{dfn}[Vérité]
	Adéquation entre la chose et l'intellect (concordance entre ce que je pense et ce qui est) (thèse de Saint-Augustin)
\end{dfn}


Pour l'instant, face à cette définition complexe, on définit la Vérité simplement comme ce qui est soit une évidence, soit démontrable.


\begin{dfn}[Libératrice]
	Qui libère, qui donne ou rend la liberté
\end{dfn}


Libérer vient de “liberare”, en latin, qui signifiait rendre sa liberté à un esclave, et signifie donc étymologiquement rendre à quelqu'un la libre disposition de son corps et de ses mouvements. Au sens figuré, cela signifie ôter d'un poids, d'une gêne, d'un souci ou d'un remord.

Ainsi, la Liberté (de libertas,atis en latin) a 2 sens : 
La libération matérielle et physique
La libération immatérielle et morale/psychique/psychologique

Cependant, attention, on ne peut pas définir la vérité comme le fait de dire la vérité (ce qui serait tourner en rond). En plus, cela s'appelle déjà la véracité, qui s'oppose au mensonge.

Pour en revenir à notre PBQ de départ/nom de séquence, on soulève immédiatement la contradiction : la vérité libère de l'ignorance, mais nous englobe et nous enferme, nous privant dans une certaine mesure de liberté vis-à-vis de la vérité elle-même.

PBQ :
Ou bien l'expérience de la Vérité me libère de mes erreurs et illusion, et alors la vérité est libératrice, ou bien la Vérité n'a aucune force libératrice face aux erreurs et illusions, et alors elle est contraignante.

\section{La Vérité est libératrice : elle me libère de mes erreurs et illusions}

\subsection{La Vérité me libère de mes illusions et mes croyances}

\begin{refe}
	Platon, \textit{La République}, Livre VII, 514a, "L'Allégorie de la Caverne"
\end{refe}

Cette allégorie dit (grossièrement) que la plupart des hommes sont dans d'illusion : ils prennent des ombres de vérités (des "projections") pour la vérité :

"Aux yeux de ces gens-là, la réalité ne saurait être autre chose que les ombres des objets"

\begin{dfn}[Croyance] 
Ce que l'on tient pour vrai, sans savoir pourquoi ou comment l'expliquer.
\end{dfn}

\begin{rep}[Croire/Savoir]
	A COMPLETER
\end{rep}

\begin{dfn}[Illusion]
Interprétation erronée de la réalité, quand bien la perception de celle-ci serait juste (peu importe si la perception est bonne ou mauvaise, donc).
\end{dfn}


On pourra dire que l'illusion est une croyance qui, généralement, nous conforte dans nos idées reçues.

"La montée de l'âme dans le monde intelligible" 
C'est ainsi que Platon décrit l'homme qui s'échappe de la caverne. Ainsi, le monde de l'ignorance, dit monde sensible, s'oppose au monde de la vérité, le monde intelligible.

L'allégorie de la Caverne illustre donc l'élévation progressive (et potentielle) des hommes : d'abord prisonniers du monde des apparences et des illusions, les hommes progressent peu à peu vers la contemplation de la vérité en s'échappant de la caverne (ce qui est toutefois rare et difficile).
Ainsi, dans ce cas de figure, la vérité est synonyme de liberté, et est donc libératrice.

\subsection{La Vérité me libère de ceux qui voudraient m'asservir}

Revenons-en à Platon : on avait noté sur l’image la présence de personnes/personnages qui créaient les ombres de la Caverne. On peut les associer aux sophistes, contemporains de Platon.
Les sophistes étaient des sortes de professeurs de philosophie qui faisaient payer leurs cours (=la sophistique). Dans ces derniers, ils enseignaient des vérités dogmatiques à leurs élèves. On soulève immédiatement leur différence avec les philosophes, qui enseignent avant tout la remise en question et le doute.
Ainsi, la différence entre philosophe et sophiste est la même que la différence entre vrai et vraisemblable.
Platon voyait cependant un certain intérêt à la sophistique : elle pouvait servir d’initiation à la philosophie.
Ainsi, on voit bien l’intérêt financier qu’avaient les sophistes à, plus ou moins volontairement, maintenir les hommes dans la doxa.
Quelques exemples de sophistes célèbres : Protagoras, Hippias et Gorgias (qui était plutôt un rhéteur).
Les sophistes produisaient des sophismes : des raisonnements faux mais vraisemblables.
=> Ainsi, dans ce cas, la philosophie, qui équivaut à la connaissance de la vérité, est libératrice ; la Vérité nous libère de ceux qui voudraient nous asservir.


Passons désormais à une autre référence philosophique :
\begin{refe}
	Kant, \textit{Réponse à la question : "Qu'est-ce que les Lumières ?"} 
\end{refe}

\begin{bio}[Kant]
	Emmanuel Kant était un philosophe des Lumières allemand (né en Prusse) qui a lutté pour le prix grès des connaissances et de la raison, et qui s'est opposé à l'obscurantisme.
\end{bio}

La devise des Lumières était : "Sapere aude", qui en latin signifie "Ose savoir".

Réponse à la question : "Qu'est-ce que les Lumières ?" est un texte court et engagé où Kant prend la défense des Lumières. Il y analyse aussi leurs conditions de propagation (aux Lumières).

Dans ce texte, la minorité est une sorte de confort ou de paresse dans lesquels le seul besoin est de payer les services des autres. C'est donc la non-nécessité de se donner de la peine et de penser.

La majorité, quant à elle, se déduit en creux de la minorité : il s'agit de la confrontation à la difficulté et du refus de rester dans la paresse.

Définitions utiles à l'explication du texte : 

\begin{dfn}[Indépendance]
Fait de jouir d'une autonomie complète, de ne pas être soumis à un pouvoir ou une force extérieurs et de pouvoir compter sur soi-même
\end{dfn}


\begin{dfn}[Autonomie] 
Fait d'être gouverné par ses propres lois, de faire ses propres choix
\end{dfn}


\begin{dfn}[Mensonge]
Dire volontairement quelque chose de faux dans l'intention de tromper autrui
\end{dfn}


Explication du texte (par paragraphes, non-présents dans le texte d'origine) :
Les hommes restent dans la minorité volontairement et par plaisir
Les hommes
L

Schéma récapitulatif :

\subsection{La Vérité est utile pour mieux agir, et donc pour être libre}

\begin{refe}
	William James, \textit{Le pragmatisme}, “Leçon 6” (1907)
\end{refe}

\begin{bio}[William James]
	William James est un philosophe américain  de la fin du XIXe et du début du XXe siècles.
\end{bio}

Le Pragmatisme, que James théorise dans Le pragmatisme,  c’est considérer que la Vérité est utile pour l’action. Ainsi, le critère de la vérité d’une idée est l’utilité de son action.
=> Une idée est vraie quand elle est utile.
=> Une proposition est vraie si son application est utile.
=> Toutes les vérités sont donc utiles ; il n’y a pas de vérité inutile.

Il énonce notamment que : 
“Le vrai consiste simplement dans ce qui est avantageux pour la pensée.”

Il explique dans la Leçon 6, pour résumer, que posséder (=connaître) des vérités, c’est avoir des “instruments pour l’action”. La vérité n’est donc qu’un “moyen” en vue d’une fin qui est l’action (et possède elle-même une autre fin). Il l’énonce ainsi :

“La possession de la vérité, au lieu d’être à elle-même sa propre fin, n’est qu’un moyen préalable à employer pour obtenir d’autres satisfactions vitales.”

\section{La Vérité est incapable de me libérer : elle est impuissante voire contraignante}

\subsection{La Vérité est contraignante}

\begin{refe}
	Nietzsche, \textit{Par-delà Bien et Mal}, Partie 1 (“Des préjugés des philosophes”, paragraphe 1)
\end{refe}

\begin{bio}[Nietzsche]
	Philosophe allemand contestataire du XIXe siècle ; style d’écriture obscur et littéraire ; philosophie critique ; s’attaque aux idées reçues de son temps et de la tradition philosophique ; a sombré dans la folie et le mutisme 11 ans avant sa mort. Une des citations les plus célèbres de Nietzsche est “Dieu est mort”.
\end{bio}

Explication du texte :
Nietzsche part du constat que la “valeur de la vérité” n’a jamais été questionnée. Il questionne sur la nécessité de la recherche de la vérité, et les causes de la non-recherche de l’ignorance, de l’incertitude, du non-vrai…
ATTENTION : Il n’en conclut pas qu’il ne faut pas rechercher la vérité . Il dit juste qu’il faut questionner ce but et potentiellement en chercher d’autres. Et c’est seulement à ce moment que l’on peut éventuellement en revenir à la quête de la vérité, si toutes les autres options ne nous conviennent pas.
Il explique que les philosophes parlent de la vérité sans jamais se questionner. Ainsi, Nietzsche pense que nous nous enfermons dans la vérité, il s’agit d’une sorte d’aliénation.

\begin{dfn}[Servitude]
Etat de non-liberté où mes choix sont dictés par une autre instance que moi, je suis dépendant de cette instance.
\end{dfn}


\begin{dfn}[Aliénation]
Sorte de servitude dans laquelle on devient étranger à soi-même. L’individu est dépossédé de ce qui le constitue au profit de l’instance qui l’asservit.
\end{dfn}


Pour Nietzsche, la véritable liberté serait de se détacher de la valeur de la vérité et de ce qui est communément admis. Il faut questionner la valeur de la vérité.
Il incite à explorer les pistes de l’illusion, de l’ignorance etc. pour ensuite en revenir à la vérité si ces autres options ne sont pas satisfaisantes :
“Il faut oser les choses”

\begin{dfn}[Libre-Arbitre]
Faculté de l’Homme COMPLETER
\end{dfn}


\begin{dfn}[Appétit]
Ce qui me fait tendre vers telle ou telle chose.
\end{dfn}


\subsection{La Vérité est avilissante}

\begin{refe}
	Rousseau, \textit{Discours sur les Sciences et les Arts}, 1750
\end{refe}

Le progrès des connaissances nous rend plus vicieux et la connaissance nous asservit. A force de connaître, on ne se satisfait plus de l’essentiel et on est obsédé par la connaissance ( différend de la liberté). En cela, la Vérité serait contraignante et avilissante : elle nous maintient dans le désir.

Rousseau va, en partie, à contre-courant des Lumières : il dénonce les points négatifs du progrès et de l'accessibilité des connaissances.

La grande idée de ce discours est résumée par la citation :
“nos âmes se sont corrompues à mesure que nos sciences et nos arts se sont avancés à la perfection”
Selon Rousseau, les Hommes sont donc dépravés par le progrès des connaissances.
Rousseau parle de “vaine curiosité” : nous sommes de plus en plus curieux pour rien. => cercle vicieux
“On a vu la vertu s’enfuir à mesure que la lumière s’élevait sur l’horizon”

Cette soif de connaissance mène selon lui à la démesure (négatif selon Rousseau) : c’est la pléonexie (= non- satisfaction de ce qu’on a, volonté du toujours plus)

Il faut tout de même nuancer que Rousseau a lui-même affirmé de son temps que ce texte était mal interprété et qu’on le diabolisait.

“Le progrès des arts, la dissolution des moeurs et le joug du Macédonien se suivirent de près”
Macédonien = Alexandre le Grand ici à mon avis

Rousseau explique qu’on inonde les gens de connaissance sans que celles-ci soient réellement approfondies. il dénonce cette érudition superficielle. Il faudrait donc selon lui être vertueux avant d'avoir des connaissance :
“Science et vertu sont incompatibles”

Ainsi, Rousseau imagine l’état de nature, où l’Homme n’est pas encore civilisé. Il énonce que l’Homme naît naturellement bon mais que c’est le passage à l’état social (lié à la connaissance) qui le pervertit, c’est le mythe du “Bon Sauvage" :

Etat de Nature
Etat Social/Civilisé
Innocence
Vice
Liberté
Servitude
Bonheur fondamental
Malheur
Ignorance
Connaissance


“La probité est fille de l’ignorance” (probité = honnêteté, innocence)
“Contagion des vaines connaissances"

Pour conclure, selon Rousseau, la Vérité ne libère pas. Au contraire, le progrès des connaissances aliène les Hommes alors que l’ignorance les rend libres et insouciants. Ainsi, selon lui, la Vérité n’est pas libératrice.

\subsection{La Vérité est blessante}
La Vérité me blesse au plus profond de mon être. Mais avant d’avoir la connaissance, on est innocent (même si on est ignorant) : on est donc heureux pour un temps (enfance par exemple).
La Vérité nous rend malheureux et l’ignorance nous maintient dans le bonheur.
On parle de "terrible" ou de "cruelle" Vérité.

\begin{refe}
	"La Genèse" \textit{La Bible} 
\end{refe}

Avant le péché originel, Adam et Eve sont insouciants  et heureux au Paradis, dans les Jardins d’Eden.

Le fruit de la connaissance est dit “Précieux pour ouvrir la connaissance”. Une fois le fruit mangé : “Les yeux de l’un et de l’autre s’ouvrirent et ils connurent qu’ils étaient nus”
La connaissance de la Vérité implique donc la pudeur, qui est liée à une sorte de honte, à quelque chose de blessant.
En parallèle de ce début de connaissance commencent le malheur de l’humanité et le joug de la pudeur et de la connaissance.
En effet, Dieu chasse Adam et Eve du Paradis, et leur vie de malheur commence :travail, douleur physique et morale, joug/esclavage des passions.

Pour conclure, la connaissance de la Vérité nous fait ressentir la honte et nous rend esclave de nos passions : elle ne nous libère donc pas.

TRANSITION :
Pour pouvoir être libre et heureux, il faut donc avoir une certaine connaissance de ce qui nous rendrait libres (par exemple, le jeune enfant n’est pas conscient de sa propre liberté, on peut donc se demander s’il est vraiment libre).


\section{C’est comprendre ce qu’est la Liberté qui nous rend libres}

\subsection{Les Hommes ne sont pas libres mais prendre conscience de cela les rend libres}

\begin{refe}
	Spinoza, \textit{Lettre 58 à Schuller}
\end{refe}

\begin{bio}[Baruch Spinoza]
	Philosophe hollandais du XVIIe siècle ; attaqué par l’Eglise et exclu de la société à cause de ses idées hors du temps telles que le fait que Dieu est toute la nature et que toute la nature est Dieu ; philosophie du Déterminisme
\end{bio}

Spinoza distingue les êtres entièrement libres qui agissent par eux-même comme Dieu (ils sont auto-déterminés et les êtres non-libres dits déterminés ou contraints, c’est-à-dire les "choses créées”.

\begin{dfn}[Déterminisme]
Doctrine selon laquelle l’ensemble du réel est un système de causes et d’effets. Le monde est régi par la loi de causalité (les mêmes causes produisent les mêmes effets)
\end{dfn}

Il y a donc possibilité d’agir sur les causes pour changer les effets.

\begin{dfn}[Fatalisme]
Doctrine selon laquelle tout est écrit d’avance. Croyance au destin.
\end{dfn}

Selon la doctrine fataliste, nous n’avons donc aucune liberté.

\subsection{La connaissance de notre liberté nous rend libres}

\begin{refe}
	Sartre, \textit{L’Existentialisme est un humanisme }(1945)
\end{refe}

“L’existence précède l’essence”
Nous naissons libres, comme une page vierge, libres de toute détermination (il n’existe pas de destin pour Sartre). Ce que je fais dans cette vie va définir qui je suis, c’est-à-dire mon essence : l’existence précède donc l’essence. Notre vie précède le sens qu’on veut lui donner et la question de savoir qui on est.
Cependant, pour les objets, c’est l'inverse : l’essence précède l’existence.

\begin{dfn}[Mauvaise foi]
	Concept sartrien, volonté de se dédouaner d’une responsabilité en invoquant une autre instance (“c’est pas moi c’est ……”)
\end{dfn}

Sartre pense donc que nous sommes toujours responsables de nos actes, puisque personne ne peut décider à notre place.

\begin{dfn}[Responsabilité]
Obligation faite à une personne de répondre de ses actes et de leurs conséquences. Obligation pour une personne de reconnaître ses actes comme étant les siens. Chez Sartre, c’est donc l'obligation de se reconnaître comme totalement libre.
\end{dfn}

Par exemple, un procès détermine la responsabilité de chacun.

Pour Sartre, l’Homme ne peut rejeter sa responsabilité. Nous devons donc selon lui assumer TOUS nos choix.

Dans la conférence L’existentialisme est un humanisme, Sartre s’oppose aux critiques aux critiques et tente de prouver que son Existentialisme est optimiste. Mais une de ses célèbres formule est ambiguë :
“Nous sommes condamnés à être libres.”
Sartre exprime ici le fait que nous ne pouvons pas nous défaire de notre liberté. Il concède donc bien que la liberté peut être perçue comme un fardeau, mais ne partage pas cet avis.

Chez Sartre, la notion centrale de la liberté est le choix. Il pense qu’on ne peut pas ne pas choisir : en effet, selon lui (et d’un point de vue stricto sensu logique), refuser de choisir est déjà choisir.

"L'Homme existe d’abord, se rencontre, surgit dans le monde et se définit après.”
Pour Sartre, l’Homme n’est “rien” naturellement. Il n’y a donc pas de nature humaine chez Sartre. Toutefois, l’Homme doit créer sa propre nature, son essence.

Toutefois, Sartre considère qu’il existe une condition humaine (ou situation humaine) : l’Homme est au monde, et est au milieu d’autres Hommes.

C’est dans le choix de savoir si notre condition est un obstacle pour nous ou pour notre liberté que réside la liberté sartrienne : c’est la liberté en situation.

Pour Sartre, la liberté c’est pouvoir choisir parmi l’univers des possibles déterminé par ma condition.

Cependant, cette liberté totale peut provoquer une angoisse :

\begin{dfn}[Angoisse]
Peur profonde d’un objet non-déterminé (la liberté par ex, la vie…). Inquiétude existentielle pouvant se manifester physiquement. Chez Sartre, sentiment que l’on ressent quand on comprend qu’on ne peut se défaire de notre liberté.
\end{dfn}

Mais c’est justement cette angoisse (passagère, sans quoi l’Existentialisme serait un pessimisme), cette prise de conscience qui nous rend libres.

Pour conclure, chez Sartre, c’est la prise de conscience de notre totale liberté qui nous rend libres.

\subsection{Les quêtes de la Vérité et de la Liberté vont de pair}

\begin{refe}
	Peter Weir, \textit{The Truman Show} (1988)
\end{refe}

Truman Burbank est, à son insu et depuis sa naissance, la vedette d’un show de téléréalité. Son monde est un plateau de tournage géant dont il est impossible pour lui de sortir. Tout n’est qu’illusion. On peut réaliser avec les mineurs du texte de Kant ou bien les prisonniers de la caverne de Platon/Socrate. Cependant, il se rend peu à peu compte que des choses anormales lui arrivent. S’ensuit alors une quête de vérité (pour la liberté), comme pour “sortir de la caverne”.

Les acteurs hypocrites peuvent quant à eux être comparés aux tuteurs ou aux sophistes.

La vie du protagoniste semble parfaite, idyllique (cf. La Genèse), mais ce n’est qu’une illusion. C’est une lente prise de conscience de la part du protagoniste qui s’effectue tout au long du film.

“On accepte la réalité du monde telle qu’elle nous est présentée”
“Truman préfère sa cellule”

\chapter{La foi s’oppose-t-elle à la raison ?}
(Foi et Raison)
\begin{dfn}[Foi]
Croyance d’ordre religieux (mais pas forcément) dans laquelle on est certain de ce en quoi on croit. Certitude, intime conviction. Adhésion ferme et entière de l’esprit. Forme de confiance (du latin fides qui signifie confiance). Assurance profonde de la véracité d’une chose, quand bien même la raison nous dirait le contraire. Sentiment intérieur et personnel que chaque croyant ressent, il y a une communion avec Dieu dans la foi. L'Être supérieur auquel est relié notre foi est transcendant.
\end{dfn}

Ainsi, d’après cette première définition, Foi et Raison s’opposeraient.

\begin{dfn}[Raison]
Faculté proprement humaine qui permet de produire des jugements et des idées. Elle permet à l'Homme de produire des raisonnements sur le monde qui  l’entoure.
\end{dfn}



\section{Foi et Raison sont irréconciliables : soit il faut faire taire la Raison pour entrer dans la Foi, soit il faut valoriser la Raison pour ne pas tomber dans la Foi}

\subsection{L'entrée dans la Foi est un saut dans l’absurde qui fait taire la Raison}

\begin{refe}
	Kierkegaard, \textit{Crainte et tremblement}
\end{refe}

\begin{bio}[Søren Kierkegaard]
	Philosophe danois du XIXe siècle. Il a étudié la théologie et publié la majeure partie de son œuvre anonymement.
\end{bio}

Dans \textit{Crainte et tremblement}, il  parle de “saut dans la foi”. Il assume donc entièrement l’irrationalité de la foi, et l’opposition entre foi et raison. Il va même plus loin en affirmant que la foi doit s’opposer à la raison, sans quoi elle n’est plus de la foi. Ainsi, pour lui, s'engager dans la foi implique de renier sa raison en toute connaissance de cause (pleine conscience), ce n’est donc pas de la folie, mais l’aboutissement d’un long processus de réflexion. Ce saut est décrit comme “dans l’absurde” : Kierkegaard définit donc la foi comme irrationnelle par nature.

\begin{dfn}[Absurde]
Ce qui s’oppose à la raison, ce qui est irrationnel ou illogique. Ce qui ne fait pas sens ou qui n’a aucun sens. 
\end{dfn}

“se jeter tête baissée, plein de confiance, dans l’absurde” 
C’est ainsi que Kierkegaard décrit la foi. Cette foi est donc une sorte de pari, ou de sacrifice (attention à ne pas confondre avec le pari pascalien).

Dans son œuvre, Kierkegaard fait référence à la légende d’Abraham en disant que ce dernier est incompris et perçu comme fou alors que c’est en réalité la preuve de sa foi inébranlable en Dieu.

\begin{dfn}[Folie]
Perte de raison qui entraîne un comportement inhabituel ou anormal, et qui est considéré comme irrationnel ou illogique. De multiples manifestations de la folie sont possibles.
\end{dfn}

Ainsi, dans cette première sous-partie, la foi se rapproche de la folie et se retrouve donc inévitablement opposée à la raison.

\begin{rep}[Transcendant/Immanent]
Ce qui est transcendant dépasse notre réalité, est au-delà de notre réalité ou d’un autre ordre différent du nôtre. Par exemple, Dieu mais aussi l'immortalité ou l’âme sont des concepts difficiles à aborder pour nous, ils sont transcendants (comme de nombreux concepts métaphysiques).
Ce qui est immanent appartient au même ordre que nous ou provient de l’intérieur de notre réalité. On parle de justice immanente par exemple.
\end{rep}


\subsection{Opposer Raison à Foi pour éviter les dérives de la religion}

Ce qui est paradoxal dans cette sous-partie est que nous irons dans le même sens général que Kierkegaard (foi et raison s’opposent naturellement, ce qui sert notre développement), tout en défendant la thèse inverse en préférant la raison à la foi.

Car, en effet, on peut aussi voir dans la légende d’Abraham chère à Kierkegaard une dérive de la religion.

\begin{refe}
	Marx, \textit{Critique de la Philosophie du droit de Hegel}
\end{refe}

\begin{bio}[Karl Marx]
	Philosophe allemand du XIXe siècle. Connu pour avoir théorisé le marxisme, dont l’application concrète deviendra le communisme, il a notamment critiqué du point de vue philosophique l’Etat et la religion.
\end{bio}

Topo sur la philosophie de Marx :

Pour Marx, l’Histoire est depuis la nuit des temps une lutte des classes sans fin pour s’emparer des moyens de production. La société contemporaine de Marx (toujours la nôtre) est capitaliste : le capitalisme est l'exploitation d’une partie de la population par l’autre (bourgeoisie-classe dominante et prolétariat-classe asservie), mais aussi un mode de production des richesses dans lequel la bourgeoisie détient le capital et le fait fructifier. L’Etat se met au service de cette classe dominante (Le Capital). Marx revendique une révolution prolétarienne dans laquelle le marxisme (qui deviendra plus tard en pratique le communisme) supplantera le capitalisme. Marx pensent que les formes juridiques, philosophiques et religieuses différentes correspondent à la situation économique de la société. Ainsi, la forme de la religion est spécifique au modèle économique.

Pour Marx, Dieu est une production/création sociale mais ce n’est pas pour autant qu’il rejette l'étude de la religion. Au contraire, il étudie le phénomène religieux pour en apprendre plus sur la société capitaliste dans laquelle il vit : “C’est l’Homme qui fait la religion”. Selon lui, l’invention de Dieu est une compensation de la souffrance humaine et une consolation face au constat de la réalité : “La religion est l’universel motif de consolation et de justification du monde”. Cette souffrance est plus précisément celle du prolétariat/peuple. Le prolétariat serait en réalité dans la survie, qui implique une vie peu plaisante et de souffrance. Il y a alors l’obligation d’inventer Dieu pour s’y soustraire. Cela résulte en une société qui perd la raison, et Dieu en est sa production. Toutefois, Marx ne blâme pas le prolétariat pour cela, c’est seulement la bourgeoisie et les capitalistes qu’il juge responsables.

Explication du texte :

Lignes 1-2 / Phrase 1 : La religion est à la fois l’expression de la souffrance, un moyen de lutter contre la souffrance et une justification de la souffrance. Nous souffrons car Dieu l’a voulu ; la religion nous soulage car elle fournit dès lors une explication de notre souffrance.

Lignes 2-4 / Phrase 2 : Pamphlet, violente critique de la religion

Ligne 4 / Phrase 3 : L’opium est une drogue qui donne accès à un bonheur illusoire et temporaire. L’opium soulage la douleur et la souffrance mais ne les soigne pas, cela masque le cœur du problème et ne s'attaque pas à la racine du problème. L’opium possède un caractère addictif : les Hommes ont besoin de la religion et ne peuvent s’en passer.

Lignes 4-9 / Phrases 4 à 6 : Marx exige la suppression de la religion pour que le peuple souffrant puisse accéder à un réel bonheur. Il incite à refouler la foi pour se retourner vers la raison. L’opposition Foi/Raison s’illustre par l’opposition Bonheur illusoire/Bonheur réel. L’auteur incite à supprimer les “antidouleurs” (la religion) pour pouvoir s'attaquer à la racine du mal, le capitalisme lui-même. La religion est entretenue par les classes dominantes (bourgeoisie) pour maintenir ce système, cette illusion en place et se garantir le contrôle. La bourgeoisie va dans le sens de la religion pour empêcher le prolétariat de s’émanciper (faire le rapprochement avec les tuteurs de Kant et les sophistes de Platon dans la séquence 1).

Lignes 9-11 / Phrase 7  : Métaphore cosmique qui permet de résumer l’idée de tout le texte. La notion de “gravitation” rappelle l’idée d’attraction, d’addiction.

Conclusion : Il faudrait détruire la religion pour ne plus vivre dans l’illusion, accéder à un réel bonheur et retrouver la raison. La foi s’oppose alors à la raison.

Note personnelle : ne pas oublier l’étymologie du mot “religion” qui provient du latin “religio”, qui signifie “ce qui unit”, “lien moral” ou “relier” → le religion est bien créée pour “relier” les Hommes

\subsection{Foi et Raison sont irréconciliables mais elles sont en cela complémentaires}

\begin{refe}
	Pascal, \textit{Préface au traité du vide }(1647), paragraphes 1 à 30 
\end{refe}

\begin{bio}[Pascal]
	Philosophe français du XVIIè siècle (1623-1662)
\end{bio}

Préface au traité du vide : le traité n’a jamais vu le jour
→ Plusieurs notions : foi, raison, liberté, nature, science
→ Dénonce 2 torts/écueils : ceux qui prétendent établir l’autorité (des Anciens) comme principe démonstratif en science ; et ceux qui prétendent se servir de la raison comme d’une preuve dans des domaines qui échappent à la raison
→ 2 types de preuves/connaissances : la raison et l’autorité
→ Objectif du texte : circonscrire les domaines où raison et autorité font loi, il y a donc le besoin de délimiter deux types de disciplines
→ Dans ce texte, Pascal se concentre sur le premier écueil
→ But : déterminer quand l’autorité est une preuve et quand elle ne l’est pas

Étude linéaire du texte :

Paragraphe 1 :

Pascal énonce ici un problème contemporain de son époque., et plus particulièrement celui qu’est l’écueil 1.

Les Anciens sont ceux qui nous ont précédé et ont trouvé des vérités avant nous ; le plus souvent, ce sont des scientifiques et des philosophes qui proviennent de l’Antiquité gréco-romaine.

Le respect des Anciens est une sorte de soumission à ce qu’ont énoncé les Anciens sans poser de question.

Pascal incite (déjà) à limiter ce respect.

Paragraphes 2 et 3 :

Il ne faut pas supprimer ce respect mais le borner. Cela permet de se prémunir des critiques “je ne prétends pas bannir leur autorité”. 

Pascal présente ici en creux ce qu’il va faire en énonçant ce qu’il ne va pas faire.

Pascal définit ici les deux arguments, types de preuve qui s’opposent : l’autorité et la raison / le raisonnement.

Paragraphe 4 :

“Les unes” que nous appellerons matières 1  et “les autres” que nous appellerons matières 2 sont ici clairement distinguées et même mises en opposition.

Les matières 1 sont les matières historiques où il n’y a pas de découverte à faire, ce sont les disciplines du témoignage et de la mémoire ; tandis que les matières 2 sont les matières du raisonnement où des découvertes sont possibles et souhaitables.

Paragraphe 5 :

Les matières 1 sont bornées car tout a déjà été trouvé ou écrit alors qu’on peut supposer que les connaissances des matières 2 sont infinies.

Paragraphe 6 :

Le respect est la soumission face à ce qui a déjà été trouvé. Il contient une idée de vénération.

La solution au problème du respect (qui est qu’il n’est pas applicable à toutes les matières) réside dans la distinction entre ces matières.

Paragraphes 7 et 8 : Enumération et exemple des matières 1

Histoire : “savoir qui fut le premier roi des Français”

Géographie : “savoir en quel lieu les géographes placent le premier méridien”

Jurisprudence (modification de la loi dans un cas particulier non prévu par la loi) : pour Pascal, le droit est donc un simple enregistrement de faits

Langues : impossible de connaître une langue sans l’écouter ou la lire → “savoir quels mots sont usités”

Théologie : matière emblématique des matières 1 pour Pascal (“surtout”), nous y reviendrons plus tard

Il s’agit là d’une liste non exhaustive bien évidemment.

Critères pour les matières 1 : “Celles qui ont pour principe ou le simple fait ou l’institution divine ou humaine” → le fait est incontestable tandis que l’institution repose sur une création 

La “connaissance entière” de ces matières est (en théorie) possible.

Paragraphe 9 :

L’autorité est la réponse au problème des paragraphes 7 et 8, mais il faut “vérifier” l’autorité de la personne citée. Il faut avoir confiance dans la figure d’autorité → foi.

\begin{dfn}[Autorité]
Pouvoir ou force de s’imposer comme étant la Vérité. C'est un pouvoir qui vaut pour lui-même sans nécessité de justification, incontestable.
\end{dfn}


Selon Pascal, la théologie serait la discipline emblématique des matières 1.

\begin{dfn}[Théologie]
Discours sur Dieu ou science de Dieu. Étude des textes sacrés pour rendre compte des principes et fondements d'une religion.
\end{dfn}


Attention, deux types de théologie existent :
Théologie scolastique//spéculative
Théologie positive/dogmatique
La raison est capable de faire des déductions, de produire des nouveautés sur Dieu. On utilise la raison dans la foi et la théologie → Thomas d’Aquin
Se contente d’accumuler les faits connus dans les Ecritures sans faire de déduction. Recensement des faits divers → Pascal


La théologie provient d’une autorité divine supérieure.

Les vérités théologiques peuvent s’opposer à la raison  et donc seule l'autorité permet de comprendre la théologie.

Les vérités de la foi sont “incompréhensibles à la raison”.

Selon Pascal, la raison doit se soumettre à ce qui la dépasse de très loin (Dieu par exemple). La raison est faible devant Dieu.

Les vérités de la foi échappent à toute discussion, on ne peut rien en tirer de nouveau.

Paragraphe 10 :

Les matières 2 sont “soumises à l’expérience et au raisonnement. Il y a une forte opposition entre les matières 1 et 2. 

Pascal ne dévalue ni la raison ni l’autorité, il exprime le fait que nous avons besoin des deux mais qu’il faut bien délimiter où chacune fait loi. Aucune preuve n'est supérieure à l’autre : “L’une avait tantôt l’avantage ; ici l’autre règne à son tour.” → On a bien ici l’idée de partage, cette dualité entre les deux types de preuve.

Dans les matières 2, la raison humaine peut accéder à ces sujets et y apporter quelque chose de nouveau. Les connaissances sont alors infinies et il y a de la place pour la nouveauté.

Paragraphe 11 : Énumération des matières 2 

Géométrie : branche des mathématiques qui étudie l’espace en utilisant les nombres, les mesures, les longueurs, les figures → on peut découvrir de nouveaux théorèmes

Arithmétique :  branche des mathématiques qui étudie les nombres entiers et les propriétés

Musique : étude de la combinaison (harmonieuse) des sons et des lois harmoniques qui en découlent

Physique : étude des entités du monde et de leurs interactions pour découvrir des lois

Médecine : art de soigner le corps et l’esprit et de prévenir les maladies → le corps humain est proportionné à l’esprit et la raison

Architecture : étude de la conception (et de la construction) des bâtiments, de l’aménagement des édifices, en prenant compte de l'environnement → structure, bâtiments

Critère des matières 2 : “sciences soumises à l’expérience et au raisonnement” 

Dans ces domaines, le perfectionnement et le progrès sont possibles. La perfection doit être visée même si elle n’est pas atteignable.

Chez les Anciens, ces connaissances étaient à l’état “d'ébauches”.

Paragraphe 12 :

Le progrès des connaissances repose sur le “temps” et la “peine” : il faut produire un effort pour progresser dans ces matières ; il y a une progression commune (du point de vue de l’humanité) et un progrès continu (par rapport aux Anciens) ;  le progrès se manifeste par une accumulation ds connaissances.

Les deux sources de la connaissance :
Sources
AUTORITÉ
RAISON
Disciplines/Matières
Histoire, Géographie, Jurisprudence, Langues, Théologie
Géométrie, Arithmétique, Musique, Physique, Médecine, Architecture
Portée de ce savoir
Bornée
“Autant que les livres dans lesquels elles sont contenues”
Infinie
“Fécondité inépuisable”
But de la recherche
Apprendre un maximum de vérités déjà établies
“Savoir ce que les auteurs ont écrit”
Découvrir des vérités encore inconnues ou cachées, inventer
“Chercher et découvrir les vérités cachées”
Type d’objets étudié
Objets infinis
Objets “disproportionnés” à l’esprit humain
Objets finis
Objets “proportionnés” à l’esprit humain
Mode d’acquisition
Lecture et écoute
Expérience et raisonnement
Progrès ?
Non
“Il n’est pas possible d’y rien ajouter”
Oui
“Les sciences soumises à l’expérience et au raisonnement doivent être augmentées pour devenir parfaites”
Attitude psychologique
Crédulité et Foi
Défiance



Paragraphe 13 : Distinction des écueils et “inversion de l’ordre des sciences”
Ecueil 1
Ecueil 2
Accorder trop d’autorité aux Anciens dans les matières 2
Ne pas accorder assez d’autorité aux Anciens dans les matières 1
Ne pas inventer de nouvelles choses dans les matières 2 : ”courage de ces timides” → ironique
Apporter des nouveautés en théologie et dans les matières 1 : “insolence de ces téméraires”
Manque de “courage” = timidité
Témérité, insolence
“N’osent pas inventer”
Osent trop, inventent top
“Aveuglement” → Erreur
“Malice” → Jugement de valeur → Faute


“Profanation”


Perversion / Inversion des ordres :
inversion des ordres qui a des conséquences désastreuses
corruption et injustice

\begin{dfn}[Ordre]
Disposition conforme à ce que l'on attend. Bonne organisation des choses qui est logique, juste et qui répond à une certaine justice.
\end{dfn}

\begin{dfn}[Profanation]
Souillure, violation du sacré. Prendre du sacré pour du profane. Porter atteinte à ce qui est sacré en lui ôtant son caractère sacré.
\end{dfn}

\begin{dfn}[Sacré]
Ce à quoi doit être voué un respect sans borne. Tout ce qui appartient à un domaine séparé (opposé à celui du profane), celui de la religion. Ce qui a un contact avec la divinité. Inspire la crainte et le respect.
\end{dfn}

\begin{dfn}[Profane]
Ce qui est dépourvu de sacré. Sans valeur ni signification religieuse. Se définit en creux du sacré.
\end{dfn}

Paragraphes 14 à 16 :

Argument de la réciprocité : il ne faut pas vouer un respect sans bornes aux Anciens car eux-mêmes ne l’ont pas fait pour leurs Anciens (explicité au paragraphe 16). Il faut borner le respect.

L’hardiesse est ici bienvenue : elle est une sorte d’équilibre entre la timidité et la témérité.

Il faut agir envers les Anciens de la même manière qu’eux ont agi envers leurs Anciens : “borner [l]e respect”que nous avons envers les Anciens comme eux l’ont fait. Il s’agit là d’une incitation à les imiter.

Remarques :
La physique représente pour Pascal les matières 2 alors que la théologie représente pour lui les matières 1.

Le thème de la distinction et de la classification est important pour Pascal car ce dernier veut s’opposer à toute forme de confusion ou de guerre civile.

Paragraphe 17 :

La connaissance de la nature est un processus progressif, qui va toujours de l’avant. Aucun retour en arrière n’est possible dans la science selon Pascal.

Il y a une opposition entre l’ordre de la réalité (ce qui est) et l’ordre de la connaissance (ce que nous avons) : le décalage temporel entre la connaissance de la nature et la nature elle-même. “Le temps les révèle d’âge en âge” ; “quoique toujours égale à elle-même, elle n’est pas toujours également connue”.

Cela s’oppose aux mystères de la foi qui sont toujours “égaux en eux-mêmes et toujours “également connus”.

Paragraphe 18 :

La nature est connue grâce à des expériences.
Ici, “intelligence” tient lieu de “connaissance”.

\begin{dfn}[Expérience] 
Vécu, expériences vécues, mes confrontations avec la vie. Ensemble des données sensibles auxquelles notre esprit est confronté dans la vie de tous les jours.
Expérience vécue répétée, et plus précisément la connaissance qu'on tire de la répétition de ces actions. Ensemble des connaissances concrètes acquises par l'habitude et le contact répété avec la réalité.
Expérimentation, procédure ou test par lesquels on teste la véracité d'une théorie ou d'une hypothèse. Observation provoquée, les évènements sont provoqués dans un cadre précis , réglementé et contrôlé.
\end{dfn}


Paragraphes 19-20 :

Les Modernes font plus d'expériences que les Anciens : ce n’est pas ne pas respecter les Anciens que de découvrir de nouvelles vérités en physique, et ils ne sont pas blâmables pour ne pas avoir trouvé ces vérités.

Il y a une idée de progression et d'accumulation des connaissances.

Nous pouvons “monter plus haut” et il y a “moins de peine” : il y a une possibilité de progresser dans les deux matières grâce au travail des Anciens et à la nouveauté.

Paragraphes 21-22 :

Pascal fait ici la distinction entre les Homme doués de raison et les animaux uniquement doués d’instinct.

La raison permet le progrès des connaissances alors que l’instinct ne le permet pas “Les effets du raisonnement augmentent sans cesse” mais “au lieu que l'instinct demeure dans un état égal”.

Ne pas faire usage de sa raison en sciences serait donc une animalisation : “traiter indignement la raison” → animalisation de l’Homme et honte

L’Homme est dans un progrès historique continuel.

Exemple des abeilles qui construisent des ruches toujours similaires par instinct : “les ruches des abeilles étaient aussi bien mesurées il y a mille ans”

Pas de progrès chez les animaux car il n’y a pas de continuité et de progrès au niveau des autres espèces animales

Cependant, chez les Hommes, une continuité est permise par la transmission des connaissances et l’apprentissage de génération en génération (grâce aux livres notamment, ou à la mémoire et à la tradition orale).

Les animaux sont inscrits dans une “perfection bornée” (un potentiel maximal, une capacité de survie optimale). Les Hommes seraient dans le même cas pour les matières relevant de l’autorité. 

“Il n’en est pas de même pour l’Homme, qui n’est produit que pour l’infinité”;

“Chacun des Hommes s'avance de jour en jour mais [...] tous les Hommes ensemble y font un continuel progrès à mesure que l’univers vieillit” : le progrès individuel et le progrès spécifique (relatif à l‘espèce) coexistent

L’Homme Universel de Pascal se compose alors de la Jeunesse, que sont les Anciens (car ce sont les premiers dans l’ordre chronologique), et de la vieillesse, que nous sommes, les Modernes (car nous arrivons les derniers chronologiquement).

Paragraphes 23-24 :

Exemple de la “Voie de lait” (Voie lactée) : les Anciens pensaient que c’était une zone de l’espace plus dense qui renvoyait mieux la lumière. Mais l’utilisation de nouveaux instruments comme la lunette d’approche ici permet de faire l’expérience de la Voie lactée et observer qu’elle se compose bien d’une “infinité de petites étoiles”.

\begin{dfn}[Instrument]
Ce qui permet d'observer (pas que visuellement) les phénomènes et faits naturels.
\end{dfn}


Ainsi, ce n’est pas parce que les choses sont imperceptibles qu’elles sont inexistantes.

Paragraphes 25-26 :

Référence à la Physique d’Aristote :
Monde supralunaire
Monde sublunaire
Monde des astres
Notre monde, la Terre
Perfection
Corruption
Immobilité
Changement


Toutefois, on observe avec la lunette astronomique des comètes s’enflammer dans l’espace. Il y donc du mouvement dans le monde supralunaire. Cet exemple illustre très bien le fait que les expériences sont importantes pour le progrès des connaissances scientifiques.

Paragraphes 27-29 :

Pascal est ce qu’on appelle un vacuiste : il a l’intuition que le vide existe. Il s’oppose en cela aux plénistes dont font partie Descartes et le Père Noel (oui oui !) : “Mon sentiment est que cet espace est vide”.

Problème historique : les fontainiers de Florence remarquent qu’ils n’arrivent pas à faire monter l’eau au-delà de 10.33m. Torricelli reproduit en 1643 cette expérience avec du mercure (76 cm).

Le savant français Mersenne ramène cette expérience, qui arrive aux oreilles de Pascal. Pascal fait alors faire par son demi-frère l’expérience du Puy de Dôme (en hauteur) pour réaliser l’expérience de Torricelli dans différentes conditions :  c’est “l'expérience d’équilibre des liqueurs” (1648). Cela permet l’invention du baromètre.

Selon les Anciens, la nature a horreur du vide. Mais c’est parce que les expériences ne montraient pas le vide.

\begin{dfn}[Réfutation/Infirmation]
Dire qu'une théorie est abandonnée et remplacée par une autre à la suite du constat d'un cas particulier. Contre-exemple. Montrer comme étant fausse ou incomplète, démentir une théorie, hypothèse ou proposition.
\end{dfn}


“Ce seul [cas] suffirait pour empêcher la définition générale” : Pascal  défend le principe de réfutation

La théorie doit être générale ou universelle.

\begin{dfn}[Théorie]
Proposition cohérente qui rend compte d’un ensemble de lois empiriques ou de certains phénomènes en les expliquant, en les classant, en les interprétant et éventuellement en en déduisant de nouveaux phénomènes.
\end{dfn}


On peut décliner les sciences en 3 catégories : les sciences formelles (mathématiques), les sciences physiques et les sciences du vivant.

La science est vouée à l’instabilité car elle est en progrès continu; Les théories sont incertaines puisqu’elles sont toujours réfutables La vérité absolue ne peut donc être trouvée en sciences.

Toutefois, c’est l’adéquation entre l’esprit et la nature, permise grâce à l'expérience, qui fait en sorte que les théories scientifiques progressent.

Paragraphe 30 :

Il y a un décalage entre l’être de la nature et la connaissance qu’on en a. Le temps et l’effort ont donc une grande importance en sciences. Il faut avoir de la patience en sciences.

Transition : Une telle attitude visant toujours à opposer foi et raison ne serait-elle pas contradictoire ? Ne pourrait-on pas en effet envisager une réconciliation entre foi et raison ?

\section{Raison et foi ne s'opposent pas nécessairement : une foi rationnelle est possible}

\subsection{La vérité de la foi n'est pas contraire à la raison}

\begin{refe}
	Thomas d'Aquin, \textit{Somme contre les Gentils}
\end{refe}

Gentils = païens ; somme = résumé

Thomas d’Aquin refuse de faire un choix entre foi et raison, et soutient même, au contraire, que la raison soutient la foi et l’”éclaire”. En effet, il serait absurde de renier sa raison puisque c’est Dieu qui nous en a dotés.

La raison ne peut jamais produire des vérités à l’encontre de la foi puisque notre raison provient de Dieu. Ce que Dieu a “inscrit” en nous est la foi, mais aussi la raison. Or, Dieu ne va pas à sa propre encontre, donc la raison et la foi sont liées.

La raison est absolue ou elle n’est pas. L’Homme ne peut pas renier en pleine conscience sa raison. Il est impossible de dire que ce que me dit ma raison est faux puisque la raison irait à l’encontre d’elle-même, ce qui est impossible (problème de langage, comme quand on dit “ce que je dis est faux”) : “il n’est même pas possible de penser que c’est faux”. Il est dans la définition même de la raison de ne pas s’opposer à elle-même.

Thomas d'Aquin, Somme contre les Gentils


Trouver le syllogisme du texte:

Si on suppose que la raison ne va pas à l'encontre d’elle-même, on a deux solutions :
alors la foi est forcément fausse puisque la raison l’emporte sur tout (Marx)
ou alors la foi n’est pas contraire à la raison et une foi rationnelle est possible.
Par conséquent, il n’y a pas d’opposition entre foi et raison selon Thomas d’Aquin.

On a donc ici une démonstration logique (syllogisme) pour montrer que la foi est rationnelle.

\subsection{La raison peut même en arriver à prouver l'existence de Dieu}

\begin{refe}
	Descartes, \textit{Méditations métaphysiques}, méditations 3 et 5
\end{refe}

L’idée de cette sous-partie est de se dire que si l’on arrive à prouver par un raisonnement logique que Dieu existe, alors la foi est nécessairement rationnelle.

\begin{bio}[René Descartes]
	Philosophe et scientifique français du début du XVIIè siècle, à l’origine de ce que l’on nomme la philosophie moderne, veut concilier philosophie et théologie.
\end{bio}

Plan des Méditations métaphysiques : 
1 : Descartes entreprend de douter de tout (cf. corpus DST1), c’est le doute hyperbolique
2 : Descartes trouve la première vérité indubitable, le “cogito” (sa propre existence)
3 : Descartes prouve l'existence de Dieu pour la première fois
4 : Descartes s’interroge sur l’erreur
5 : Descartes prouve d’une autre manière l’existence de Dieu
6 : Descartes prouve enfin l’existence des choses matérielles (le reste du monde).

On a donc ici une démarche très rigoureuse et rationnelle, quasi-scientifique.

Dans la “Lettre dédicatoire” (=préface), Descartes s’adresse à des théologiens pour expliquer ce qu’il s’apprête à faire.

"Tout ce qui se peut savoir de Dieu peut être montré par des raisons qu'il n'est pas besoin de chercher ailleurs que dans nous-mêmes, et que notre esprit seul est capable de fournir." On a ici une différence avec Pascal : en effet, au contraire de ce que dit Pascal, la raison pourrait appréhender le divin. Selon Descartes, l'esprit humain est proportionné à Dieu et c'est le SEUL moyen pour comprendre les vérités de la foi.

Dans la méditation 3, on a la preuve dite "par l'infini" de l'existence de Dieu : les hommes sont finis mais sont capables d'appréhender l'infini, c'est donc qu'un être infini (Dieu) a placé dans notre esprit l'idée de l'infini, c'est la marque de notre création par Dieu : "Je n'aurais pas néanmoins l'idée d'une substance infinie, moi qui suis un être fini".

"De cela seul que j'existe, et que l'idée d'un être souverainement parfait (c'est-à-dire Dieu) est en moi, l'existence de Dieu est très évidemment prouvée." (Méditation 2)

Dans la méditation 5, on a la preuve dite "ontologique" (ce qui a un rapport avec l'existence ou la réalité) de l'existence de Dieu : Descartes part du principe qu'il y a une distinction entre l'essence (ce qui est au plus profond de nous-mêmes, nos caractéristiques) et l'existence (avoir une vie concrète, être là). Selon lui, chez Dieu, l'existence et l'essence ne peuvent être séparées puisque Dieu est parfait, donc Dieu existe : "Il n'est pas en ma liberté de concevoir un Dieu dans existence".


\subsection{La foi doit se fonder sur la raison sinon elle risque d'être faible, voire dangereuse}

\begin{refe}
	Kant, \textit{La religion dans les limites de la simple raison}
\end{refe}

Ce livre a pour but de démystifier la religion.

Pour Kant, la religion peut être appréhendée de manière rationnelle, c'est la religion naturelle, qui s'oppose à la religion révélée.

Plan de La religion dans les limites de la simple raison :
1ere partie : problème du mal et de ce que ce principe implique pour la liberté humaine
2eme partie : figure du Christ (= bonté et bien absolu) et la lutte contre le mal
3eme partie : caractéristiques de la religion et distinction entre la religion naturelle et la religion révélée

La différence majeure entre les religions révélée et naturelle réside dans la manière de connaître Dieu.

\begin{dfn}[Révélation]
Événement particulier, historique, précis qui fait que Dieu se dévoile, se révèle aux yeux des Hommes alors qu'il était caché
\end{dfn}


Dans la religion révélée, Dieu est connu par la révélation.

Dans la religion naturelle, Dieu est connu par la raison. Ma simple observation du monde, de la nature me fait déduire l'existence d'un être qui aurait créé ce monde, en l'occurrence Dieu. Cette capacité d'observation est innée, ce qui explique les termes de religion naturelle.

C'est parce que je m'aperçois qu'il y a un certain ordre et une certaine beauté dans la nature, dans le monde, qui ne pourraient venir que d'un créateur divin, que j'en déduis que Dieu existe.

Si Dieu était connu par la révélation, la croyance en Dieu ne serait qu'une croyance historique qui serait contingente, donc seule la religion naturelle est nécessaire. Ainsi, selon Kant, la religion doit être naturelle mais peut aussi être naturelle et révélée.
Dans tous les cas, pour toutes les religions, les Hommes "auraient pu ou dû d'aviser de cette religion par eux-mêmes".

Ainsi, la religion n'est pas du tout obscure et la foi ne s'oppose pas pas à la raison.

"Chacun peut dorénavant se convaincre de la vérité de cette religion par lui-même et sa propre raison"

Il y a plusieurs risques à une religion uniquement révélée :
Elle risque de tomber dans l'oubli : "elle disparaît du monde", en effet la révélation peut être oubliée ou perdue
Il faut renouveler cette révélation : "il faut renouveler de temps en temps publiquement cette révélation"
Il faut consigner cette religion dans un livre puis la diffuser en publiant le livre

\begin{dfn}[Superstition]
Croyance totalement irrationnelle selon laquelle le monde est constitué de signes (dont la valeur est heureuse ou funeste). La superstition vient d'une forme d'ignorance mais aussi et surtout d'une crainte.
\end{dfn}


Par exemple, croire qu'une action est nécessaire pour empêcher un grand malheur, quand bien même cette action n'aurait rien à voir avec le malheur.


Transition : Mais cette conception selon laquelle la raison est toujours fiable est discutable et relève d'une foi, d'une croyance en la toute-puissance de la raison.

\section{La confiance en la raison est elle-même une forme de croyance}

\subsection{La science, qui assure ne compter que sur la raison, relève toujours d'une forme de croyance}

\begin{refe}
	Nietzsche, \textit{Le Gai savoir} (1882), Livre V, §343-347 (344 en particulier)
\end{refe}

La science possède un paradoxe : on présuppose, on a une croyance dans le fait que notre raison est infaillible et mène au progrès, ce qui est en réalité une forme de foi.

Dans \textit{Ecce homo} (retrospective de Nietzsche suir sa vie et son oeuvre), Nietzsche décrit Le Gai savoir comme étant celui de la joie, où il est plein de vie et célèbre la connaissance après avoir vaincu une maladie.

Nietzsche, Le Gai savoir, Livre V, §344

Dans Le Gai savoir, Nietzsche explique que la croyance en Dieu décline, avec la célèbre formule “Dieu est mort” (§125, III). Dieu n’existe que dans notre esprit, mais comme cette foi décline, alors l’existence de Dieu elle-même décline (Dieu meurt), c’est en tout cas ce que constate Nietzsche.
Il constate alors que la science remplace la religion en tant que foi : “La croyance au Dieu chrétien est tombée en discrédit” (§343). En effet, le dieu chrétien est devenu “indigne de croyance”.

Ainsi, pour Nietzsche, la croyance s’est transformée. Avec l’avènement des sciences, la croyance a pris une autre forme : ce n’est plus la croyance en Dieu mais en la raison et la capacité de la science à trouver la vérité. Nietzsche s’oppose donc au postulat selon lequel les scientifiques abandonnent leurs croyances pour chercher la vérité.

\begin{dfn}[Piété]
Extrême croyance, croyance exagérée
\end{dfn}


La science reste attachée à une vieille tradition : la vérité est incontestable. La science serait débarrassée de toute croyance.
Mais il existe toujours une croyance en l’existence de la vérité, sur laquelle repose la science.
“Où donc la science prendrait-elle son absolue croyance, sa conviction sur laquelle elle repose, à savoir que la vérité serait plus importante que toute autre chose.” : cette question rhétorique prouve qu’il n’y a pas de réel fondement à cette conviction.

\begin{dfn}[Conviction]
Certitude de l’esprit en la véracité d’une chose, forte croyance. On parle souvent d’intime conviction, donc la conviction provient du plus profond de nous-mêmes. La conviction vient à la suite d’un jugement clair et lucide (clarté et stabilité).
\end{dfn}


Nietzsche définit ceux qui n’auraient aucune croyance comme des “esprits libres” (parmi lesquels il se classe d’ailleurs, preuve de sa grande modestie).
Nietzsche insiste pour se défaire de cette croyance en la vérité qui serait métaphysique, transcendante. Cette croyance de la science en la vérité est donc aussi une croyance en la raison.

Ainsi, selon Nietzsche, la science, et donc la raison, relèvent toujours de la croyance.

\subsection{L'athéisme et l'agnosticisme relèvent toujours d'une forme de croyance}

\begin{refe}
	Nietzsche, \textit{Généalogie de la morale} (1887), Traité III, §24-25
\end{refe}

Dans la Généalogie de la morale, Nietzsche cherche à comprendre comment la morale est née, pour mieux l'attaquer ensuite.

\begin{dfn}[Athéisme]
Nier fermement l'existence de Dieu, ne pas croire en Dieu. Revendiquer la non-existence de Dieu
\end{dfn}



\begin{dfn}[Agnosticisme] 
Ne pas se prononcer sur l'existence de Dieu. Dire qu'on ne peut pas savoir si Dieu existe ou pas. Ne pas s'engager, ne pas trancher sur l'existence de Dieu puisque celle-ci nous dépasserait de toute façon (transcendance) dans l'hypothèse où elle est réelle.
\end{dfn}


Pour Nietzsche, les athées manifestent leur foi en s'opposant à la foi en Dieu puisqu'ils ont une croyance en la véracité de leur affirmation. Ce ne sont donc pas des esprits libres puisqu'ils veulent rétablir une vérité, celle que Dieu n'existe pas.

Pour Nietzsche, les agnostiques "adorent le point d'interrogation lui-même comme un Dieu". Il y aurait donc chez les agnostiques une foi dans le questionnement. Ils sont qualifiés "[d'] idolâtres de l'inconnu et du mystérieux en soi".

Les athées se revendiquent des "incroyants", des "antichrétiens" et se croient esprits libres mais Nietzsche leur montre le contraire puisqu'ils croient à la valeur de la vérité : "[les athées] sont encore loin d'être des esprits libres car ils croient encore à la vérité". Ce sont des "prétendus" esprits libres : "Ils croient en la vérité plus fermement, plus inconditionnellement que les autres". Nietzsche dit que c'est une forme de "foi inconsciente".

Nietzsche incite donc à questionner ces croyances et la valeur de la vérité (cf. Sqce 1) : "Nous devons une bonne fois et de façon expérimentale mettre en question la valeur de la vérité.".

\subsection{Les productions que l'on attribue à la raison relèvent en réalité d'une croyance}

\begin{refe}
	Hume, \textit{Enquête sur l'entendement humain}
\end{refe}

\begin{bio}[David Hume]
	Philosophe empiriste du début du XVIIIe siècle 
\end{bio}

Empirisme = Toutes nos connaissances viennent de l'expérience, les connaissances sont dites "a posteriori". La connaissance n'est pas fondée sur des "a priori". La connaissance vient précisément de l'expérience, d'une habitude.

Ceux qui pensent que la connaissance est fondée sur des "a priori" pensent que nous avons des schèmes de pensée pour appréhender l'expérience.

Par exemple, la notion de temps (antériorité, simultanéité et postériorité) est perçue par les uns comme innée et acquise pour les autres.

Pour Kant, le principe de causalité est connu a priori : tout événement est a une cause. Alors que pour Hume, à l'inverse, le principe de causalité est a posteriori, c'est-à-dire qu'il tient de l'expérience qu'on en fait.

Le principe de causalité, c'est dire qu'il existe nécessairement un lien entre la cause et l'effet.

Par exemple, lorsque je m'approche du feu, ça brûle. J'en induis que le feu brûle et chauffe. C'est l'habitude qui va permettre de prévoir. Toutefois, Hume explique que cela relève d'une forme de croyance (croire en ce lien). Il veut donc prouver ici qu'une grande partie de notre savoir vient d'une croyance.

La confiance en la raison d'une forme de foi puisque ce que l'on suit provient d'une forme de croyance.

Pour Hume, les généralisations, qui ne peuvent être universelles, sont des inductions.

\begin{dfn}[Induction]
Type de raisonnement qui tire des lois générales à partir de l'observation d'un ou plusieurs faits particuliers. C'est passer de l'effet à la cause.
\end{dfn}

L’induction est l’inverse de la déduction.

"Le contraire d'un fait est toujours possible." Hume prend l'exemple du soleil : il se pourrait que le soleil ne se lève pas demain, et nous ne pouvons pas l'infirmer.

Le contraire d'un fait auquel nous sommes habitués est toujours possible : "Ne puis-je point concevoir que cent événements différents pourraient aussi bien suivre de cette cause" (section 4).

Pour Hume, c'est ma raison qui ajoute cette dimension nécessaire, la nécessité n'est pas dans les choses elles-mêmes, et c'est ajout est inexplicable, c'est donc une croyance.

Tout le travail de la raison reposerait donc que une croyance, la croyance que les mêmes causes produisent toujours les mêmes effets (déterminisme, Spinoza, sqce 1). Cela rejoint la croyance selon laquelle le futur sera semblable au passé et au présent : "Toutes nos conclusions expérimentales procèdent de la supposition que le futur sera conforme au passé." Ce n'est donc qu'une hypothèse, pas quelque chose de certain.

Cette croyance, bien qu'infondée, est néanmoins utile puisqu'elle nous empêche de tomber dans une paranoïa constante. Il n'est donc pas nécessaire de s'en débarrasser mais simplement d'en prendre conscience.

Dans la section 3, Hume fait le constat de cette croyance et dans la section 4, il critique ce lien de causalité et cherche l'origine de cette croyance, il va étudier la tendance humaine qui nous pousse à savoir cela.

Cette connaissance de la relation de cause à effet naît de l'espérance. Qu'est-ce qui, dans notre esprit, nous fait créer un lien entre causes et effets ? Cette connaissance vient plus précisément de l'expérience, de l'habitude, de l'accoutumance.


\begin{dfn}[Habitude]
	Tendance à reproduire ce qu'on a toujours fait.
Tendance à attendre ce qu'on a toujours expérimenté.
\end{dfn}


Hume dit que l'habitude est "le grand guide de la vie humaine".

Selon Hume, l'habitude contractée nous l'idée que la succession de causes et d'effets se reproduira toujours. Or, cette connexion est "dans notre pensée" (section 7), elle est produite par mon esprit, et à force d'expérimenter, mon esprit se met à croire que la cause et l'effet sont liés.

Pour Hume, cette induction n'est donc pas un savoir mais une croyance. Pour lui, cette croyance est naturelle chez nous, c'est une "tendance mécanique".

Toutefois, sous cette croyance, nous serions dans le doute constant et l'incapacité d'anticiper, donc cette croyance est utile dans notre vie pratique. Ainsi, il ne faut pas l'abandonner puisque toute notre connaissance repose sur cette croyance (Hume est un sceptique modéré). D'un point de vue pragmatique, cette croyance est donc une vérité.

Mais il faut proportionner cette croyance : il faut amoindrir les effets peu fréquents et prendre plus en compte les effets fréquents (section 6). "Il faut que nous assignions à chacun de ces effets une autorité et un poids particuliers en proportion de la plus ou moins grande fréquence que nous avons leur avons trouvée." La section 6 parle de probabilités et de pondérer notre croyance selon les probabilités des évènements.

\chapter{Peut-on se connaître soi-même ?}
(Conscience et Inconscient)

Problème dans la question de la connaissance de soi : le sujet connaissant et l'objet connu/connaissable sont la même chose.

"Peut-on ?" interroge la capacité et la légitimité.

Provisoirement, nous dirons que la capacité à se connaître soi-même, c'est la conscience.

Longtemps, le mot conscience a eu une signification seulement morale (avant Descartes).

\begin{dfn}[Conscience morale]
	Capacité que nous avons de sentir et de juger la moralité de nos actions
\end{dfn}

=> Mauvaise ou bonne conscience

Le sens moderne du mot "conscience" renvoie depuis Descartes au domaine de la connaissance de soi. Étymologiquement,le mot renvoyait déjà à la connaissance (implicitement de soi), provenant du latin "cum scientia", qui signifie "accompagné de science, de connaissance".

\begin{dfn}[Conscience]
	Connaissance tournée vers soi-même, vers son intériorité. Capacité de l'esprit humain à percevoir ce qu'il se passe en lui et à savoir ce qu'il pense. Permet d'être présent à soi-même.
\end{dfn}


La connaissance de soi est-elle valide, valable, fiable, exhaustive (existence de potentielles zones d'ombre ?) ? La conscience est-elle totalement transparente à elle-même ?

PBQ :
Ou bien la conscience que j'ai de moi-même est immédiate et infaillible, et alors je peux avoir une connaissance complète et exhaustive de moi-même ; ou bien la conscience que j'ai de moi-même souffre des médiations et est faillible, et alors je ne peux pas avoir une connaissance complète et exhaustive de moi-même. La conscience de soi est-elle infaillible ?


\section{Je peux parvenir à me connaître car je peux avoir une conscience de moi claire et infaillible}

Grâce à ma conscience, je suis la personne la plus apte à me connaître moi-même et me permet d'avoir une connaissance de moi immédiate.

\begin{rep}[Médiat/Immédiat]
Immédiat, ce qui ne souffre, n'implique aucune médiation, c'est-à-dire aucune intermédiaire, ce qui met en relation deux termes sans avoir besoin d'aucun autre terme.
Médiat, ce qui passe par une médiation, un intermédiaire, ce qui met en relation deux termes via un autre terme, sinon la relation serait impossible.
\end{rep}


\subsection{La conscience de soi comme une faculté présente en chaque être humain}

Mais est-ce seulement le privilège des humains ?

\begin{refe}
	Descartes, \textit{Lettre à Morus du 5 février 1649}
\end{refe}

Thèse : La conscience de soi est le propre de l'Homme, ce qui le différencie des animaux

Pour Descartes, la conscience de soi est un retour sur soi, un retour réflexif par lequel je sais ce que je pense.

Pour Descartes (et dans toute cette première partie), la conscience de soi est une saisie immédiate de la pensée.

1/ La conscience de soi permet de lier corps et esprit : elle permet aux hommes de produire des mouvements voulus et pensés par eux-mêmes. Par exemple, je veux marcher, je marche. De mes volontés/volitions découlent des actions. Elle engendre la liberté de l'Homme.

2/ La raison implique le langage : la conscience de soi va de pair avec la raison et cette dernière se manifeste par le langage. Le fait de parler prouve qu'on est doué de raison. Le langage est, chez Descartes, un signe pour marquer quelque chose se rapportant à la seule pensée. "Le langage se rapporte à la pensée"."La parole est l'unique signe et la seule marque assurée de la pensée cachée et renfermée dans les corps".

L'animal, à l'inverse de l'Homme, n'a pas de conscience de soi d'après Descartes mais il n'en a pas vraiment de preuve ( aveu dans la lettre). "On ne saurait prouver qu'il y ait des pensées dans les bêtes". Descartes compare les animaux à des machines : l'animal est comparé à une horloge, il agirait dans une perfection bornée, comme les hirondelles qui reviennent toujours.

Descartes dit que les animaux n'agissent pas selon leur volonté : un mouvement "sans le secours d'aucune pensée". Descartes appelle cela des "mouvements convulsifs" => automates instinct. 

Selon Descartes, les animaux n'ont pas de raison car ils n'ont pas de langage : "On n'a point cependant encore observé qu'aucun animal fut parvenu à ce degré de perfection". Les cris ne sont pas pour Descartes des manifestations de la pensée, les cris ne se rapportent qu'à leurs émotions et pas à leurs pensées. L'argument de Descartes est que si les animaux avaient des pensées, ils nous les communiqueraient.

\subsection{Grâce à la Conscience de soi, je peux parvenir à une connaissance claire et infaillible de moi-même et de ma propre existence}

\begin{refe}
	Descartes, \textit{Méditations métaphysiques}, méditations 1 et 2
\end{refe}

Dans ses \textit{Méditations métaphysiques}, Descartes recherche les “premiers principes de toute chose”.

En quoi la conscience que j'ai moi-même peut-elle me permettre d'être certain de mon existence ?

“Des choses que l’on peut révoquer en doute” : Renvoie à la première méditation.

Depuis son enfance, on a toujours fait croire à Descartes des choses fausses, et il s’en aperçoit. Descartes se dit alors que tout peut potentiellement être faux, il rejette donc toute proposition et se met à douter de tout, c’est le doute hyperbolique.
Le doute est aussi dit ”méthodique” car il est volontaire (il manifeste la liberté).
Il est dit “hyperbolique” car rien n’y échappe.

Descartes doute de tout et en vient même à douter de sa propre existence. Il essaye de dissocier l’évidence de la réalité.

A la fin de la première méditation, Descartes n’a plus aucun repère et est dans un état de torpeur.

Comment Descartes peut-il arriver à la vérité alors qu’il a tout remis en doute ?

Dans le premier paragraphe, Descartes rappelle comment s’est terminée la première méditation

Descartes cherche une vérité indubitable et garde en vue son objectif dans le deuxième paragraphe (“chemin”=méthode).

Dans le troisième paragraphe, Descartes énonce cette première vérité : “je suis, j’existe”. Si on peut me tromper, je reste quelque chose, donc j’existe par le fait qu’on peut me tromper. La condition pour me tromper est que je sois sois, que j'existe : si je doute, on peut me tromper donc j’existe.

Dans le quatrième paragraphe, Descartes se demande ce qu’il est puisqu'il existe. Dans le passage omis, il se demande s’il est un corps : non car . Il se demande alors s’il est une âme et énumère les attributs de l’âme : se nourrir, marcher et sentir (par rapport aux sens) nécessitent un corps. Or, Descartes a douté de l’existence de ce corps. Mais il reste un dernier attribut, la pensée, qui ne nécessite pas le corps.

Dans le cinquième paragraphe, on a la conclusion de cette deuxième étape. Je sais un peu plus précisément ce que je suis, je suis un esprit, une chose qui pense (une “res cogitans”). On n’existe alors que par la pensée. Il y a une identité de la pensée et de l’existence : j’existe, je suis en tant que je pense.

A partir du sixième paragraphe, Descartes tente d’apporter plus de précision. Il procède à une définition de la pensée.
Concevoir (similaire à entendre) se rapporte aux concepts et représentations mentales abstraites et générales.
Imaginer est se faire une image d’une chose (au sens large, pas que visuel).
Descartes réduit les attributs de l’âme aux attributs de la pensée. Penser, pour Descartes, c’est tout ce qui se passe dans notre esprit. La forme des pensées est toujours vraie, mais le contenu peut être faux (cf. imagination). Les sensations peuvent être fausses, je peux avoir une hallucination, mais le fait que je sente quelque chose demeure vrai.

\begin{dfn}[Introspection]
S’examiner soi-même, s’observer soi-même, regarder en soi-même. 2 buts possibles : regarder en soi-même pour se comprendre psychologiquement ; et regarder en soi-même pour comprendre le fonctionnement de l'esprit humain en général.
\end{dfn}


Pour Descartes, la conscience est la présence immédiate à soi par la connaissance immédiate de ses pensées. Toute pensée est alors consciente d’elle-même et cela implique qu’aucune pensée ne serait inconsciente.

Au septième paragraphe, Descartes illustre son propos :de la perception de chapeaux, je juge qu’il y a des Hommes. Mais le contenu de cette pensée pourrait être faux, ce pourraient être des spectres ou des automates. Toutefois, la forme de la pensée reste vraie malgré tout.

Descartes s’imagine voir des Hommes par la fenêtre mais le contenu de cette pensée peut être faux. Mais, indépendamment de son contenu, la forme de la pensée reste vraie, et donc je pense, et par conséquent j’existe.

Le doute de Descartes est donc hyperbolique, méthodique mais aussi provisoire. Il se différencie là du doute sceptique, c’est un doute sceptique feint pour trouver des vérités.

La conscience de soi chez Descartes permet de savoir ce que je suis mais ne permet pas de savoir qui je suis.

\subsection{La conscience de soi permet la continuité de l’identité A COMPLETER}

\begin{rep}[Identité / Égalité / Différence]
L’identité, c’est le caractère de ce qui reste le même à travers le temps, du latin “idem”, “le même”. L’identité personnelle, c’est l’identité de la personne, de la personnalité, du moi, c’est le fait que je suis le même à travers le temps.
L’égalité, c’est le principe de droit selon lequel tous les individus ont la même valeur devant la loi.
La différence, c’est le caractère par lequel une chose se distingue d’une autre. Pour que deux choses puissent être différentes, il faut qu’elles aient un socle commun.
\end{rep}




\section{Il est difficile, voir impossible, de se connaître soi-même car la conscience de soi est faillible}

\subsection{On peut remettre en question l’existence du “moi”}

\begin{refe}
	Pascal, \textit{Pensées}, L688 “Qu’est-ce que le moi ?”
\end{refe}

Comment peut-on réduire un être à une unité que serait le moi, et qu’est-ce qui rend possible cette identification de l’unité et de la singularité de la personne ?

Pour Pascal, cette réduction d’une personne à un “moi” est illusoire. Ce fragment est une critique plus ou moins directe de Descartes et de ses Méditations métaphysiques, et surtout de la conclusion de la deuxième méditation qui est que l’esprit (le “moi”) est plus facile à connaître que le corps. La première vérité énoncée par Descartes n’est alors plus du tout évidente.Dans ce fragment, les thèmes de l’amour et du “moi” sont très liés. Quand j’aime une personne, qu’est-ce que j’aime chez cette personne ? La personne elle-même ? Non, seulement certaines de ces qualités. Mais si le “moi” existe, on ne peut pas le réduire à ces qualités.
Ainsi, on ne peut pas savoir ce qu’est le “moi”, le “moi” est introuvable. Je ne peux pas trouver un “moi” qui serait la substance de toutes mes qualités. Attention, Pascal ne dit pas que le “moi” n’existe pas, mais qu’il est très complexe à trouver.

Pour Pascal, il est impossible de déterminer ce qui forme l’essence singulière de chacun, de chaque moi : il retourne l’exemple de Descartes.

Pourquoi aime-t-on quelqu’un ? Pour sa beauté ? Mais la beauté est éphémère, elle ne peut donc pas seule constituer le “moi” puisqu’elle peut disparaître alors que le “moi” ne disparaît pas.

D’autres raisons de l’amour : le jugement et la mémoire ? Là aussi ce sont des qualités éphémères : on peut les perdre sans perdre son “moi”. 

Il y a une inconstance des qualités corporelles et intellectuelles, elles sont donc accidentelles.

\begin{rep}[Accidentel / Essentiel]
Est essentiel ce qui appartient à l’essence d’une chose, ce qui est nécessaire à l’existence d’une chose : si on supprime ce qui est essentiel, on supprime la chose.
Est accidentel (du latin \textit{accidere}, “arriver par hasard”) ce qui ne fait pas partie de cette essence, ce qui peut être supprimé sans supprimer la chose. Simple qualité qui n’altère pas l’existence.
\end{rep}


Ainsi, le “moi” est introuvable ni dans le corps, ni dans l’esprit. Mais dans l’amour, on ne peut pas aimer l’âme d’une personne : on ne peut donc aimer quelqu’un abstraitement, ou quelle que soit sa qualité. 
Pascal souligne l’importance des conventions sociales et la nécessité de reconnaître.

Pour conclure, le “moi” est inatteignable, il n’est qu’une idée abstraite qui n’existe pas dans la réalité.

\subsection{Des facteurs extérieurs et inconscients déterminent ma conscience}

\begin{refe}
	Marx et Engels, \textit{L'idéologie allemande} (1845)
\end{refe}

\begin{bio}[Engels]
	Engels est un ami philosophe de Marx qui a beaucoup collaboré avec ce dernier.
\end{bio}

La manière dont nous construisons nous goûts et nos envies est le fruit de l'influence de notre famille, de notre milieu social et de notre environnement en général (amis, entourage mais aussi éducation par exemple). Nous sommes dépendants de notre culture.

Comment tout cela peut-il limiter la connaissance que nous avons de nous-mêmes ?

Marx et Engels partent du principe que le sujet humain n'est pas isolé, il reçoit diverses influences de la société : on les nommes des déterminismes, qui peuvent être plus ou moins forts.

Nous ne sommes pas vraiment des sujets libres, nous ne sommes pas transparents envers nous-mêmes. Nous subissons des médiations extérieures de notre environnement (il s'agit surtout de la classe sociale pour Marx et Engels). 

"Ce n'est pas la conscience qui détermine la vie, mais la vie qui détermine la conscience" : la conscience n'est pas une donnée première et médiate, elle n'est pas vierge de toute détermination (cf. Sartre Sqce 1)

Marx s'oppose à l'idée d'une conscience libre : pour Marx, le sujet n'est pas entièrement libre.
La conscience n'est pas vierge, elle est surtout le fruit des influences de notre environnement.

La conscience ne détermine pas notre rapport au monde, mais c'est le monde qui détermine notre conscience (voir citation précédente).

Le sujet est le produit des influences historiques et sociologiques de sa vie.

Ainsi, les consciences du bourgeois et du prolétaire sont différentes du fait de leur naissance avant tout.

Cette détermination de la conscience est la plupart du temps inconsciente. Le sujet a en effet rarement assez de recul pour se rendre compte qu'il s'agit d'une construction sociale.

\begin{dfn}[Sociologie] 
Science des faits sociaux humains et des groupes sociaux. La sociologie étudie comment les sociétés fonctionnent et se transforment. La sociologie pense aussi la capacité qu'ont les hommes à se détacher de ces déterminismes.
\end{dfn}


Tous ces mécanismes décrits par Marx et Engels limitent la connaissance que nous avons de nous-mêmes, mais dans une moindre mesure car nous pourrions connaître ces déterminismes et ainsi nous en détacher.


\subsection{Le psychisme humain est en grande partie inconscient : l'hypothèse de l'inconscient par la psychanalyse}

ATTENTION, l'inconscient tel qu'il sera développé ici est hypothétique, il n'est pas prouvé scientifiquement qu'il existe mais cela est fortement probable car sans lui, de nombreux comportements seraient inexplicables. Nous développerons ici la théorie du point de vue de Freud, qui suppose avérée l'existence de l'inconscient.

Pour commencer, il est impératif de distinguer l'inconscience de l'inconscient :

L'inconscient devient donc une instance à part entière dans le psychisme humain.
L'inconscient n'est plus seulement une sous-conscience (comme l'inconscience) mais plutôt un rival à la conscience ayant une certaine force.

\begin{refe}
	Freud, \textit{Introduction à la psychanalyse}
\end{refe}

Princesse Marie, Benoît Jacquot → Téléfilm sur Freud et l'inconscient

Sigmund Freud : Philosophe mais surtout neurologue autrichien des XIXe et XXe siècles (1856-1939)

Des pulsions irrationnelles et inconscientes déterminent ce que nous pensons et ce que nous pensons d'après Freud.

\begin{dfn}[Pulsions]
	Excitations psychiques d'origine somatique visant un objet et recherchant une satisfaction 
Vient de l'allemand trieb(en) : pousser vers l'avant
Les pulsions sont l'expression de désirs et d'instincts profonds qui peuvent apparaître masqués, déguisés et qui dirigent nos actions sans que nous en soyons conscients. Ces pulsions masquées se situent dans l'inconscient.
\end{dfn}


Freud se fonde sur les travaux de Breuer sur l'hypnose pour traiter l'hystérie. Cette maladie avait des symptômes physiques (tremblements etc.) mais ne s'expliquant pas par une maladie purement physique.

\begin{dfn}[ Symptôme]
	Indice, manifestation d'une maladie. C'est ce que le patient rapporte au médecin et/ou ce que le médecin observe.
Selon Freud, il s'agit d'une "substitution à la place de quelque chose qui n'a pas réussi à se manifester au-dehors".
\end{dfn}


Breuer parvient à traiter ces malades par l'hypnose, en revient l'émotion associée à ces symptômes. Il y a une décharge émotionnelle, c'est la méthode cathartique (Catharsis, Aristote → Tragédie, purge le mal par le mal)

Freud reprend cette méthode et l'appelle la cure psychanalytique : le patient s'allonge dans un fauteuil sans voir le psychanalyste, et le patient dit tout ce qui lui passe par la tête, sans se censurer. Le psychanalyste orienté le patient et le pousse à ne surtout pas s'auto-censurer. La cure psychanalytique repose sur un "échange de paroles". La parole est donc extrêmement importante dans le processus de guérison.

Freud se rend compte que les symptômes sont dûs à une pensée dite "refoulée" (le refoulement), une idée écartée de la conscience du patient car insoutenable, intolérable, immorale, taboue de façon inconsciente. Les pensée/pulsions inconscientes refoulées vont lutter pour s'auto-satisfaire de manière masquée, ce qui crée une tension psychique et se manifeste par des symptômes. Freud observe plus précisément que ces symptômes ont souvent pour origine des traumatismes (souvent sexuels pour Freud), qui ont souvent eu lieu dans l'enfance.

Le traumatisme originaire est souvent lié à des expériences sexuelles réelles ou fantasmées dans l'enfance (le complexe d'Œdipe par exemple).

\begin{dfn}[Traumatisme]
	Évènement tellement brutal, choquant et violent émotionnellement que c'est comme s'il n'avait pas été vécu (en première personne) par le sujet. Le traumatisme a une influence durable (négative) sur la personne → c'est donc subjectif, l'influence dépend des personnes.
\end{dfn}


\begin{refe}
	Freud, \textit{5 leçons sur la psychanalyse}, cas de Elizabeth Von R.
\end{refe}

Elizabeth Von R. avait des “douleurs violentes” dans les jambes. Elizabeth était amoureuse du mari de sa sœur. Sa sœur meurt. L’idée d’épouser son beau-frère surgit en elle sur le lit de mort de sa sœur, mais cette pensée est intolérable au vu des mœurs de l’époque et est donc refoulée. C’est ce refoulement qui a/aurait produit l’hystérie.

\begin{dfn}[Refoulement]
	Processus inconscient par lequel des représentations liées à des pulsions intolérables sont repoussées hors du champ de la conscience.
\end{dfn}


Un fois qu'Elizabeth se souvient de cet évènement et l’avoue à Freud, elle commence à guérir de son “hystérie” : c’est l’abréaction.

\begin{dfn}[Abréaction]
	Action par laquelle, par l’exercice de la parole pratiqué lors de la cure psychanalytique, le patient se libère d’une pulsion refoulée en la faisant surgir à la conscience
\end{dfn}


“Notre théorie agit en transformant l’inconscient en conscient”

\begin{refe}
	Freud, \textit{Le moi et le ça}, 2ème topique (années 1920)
\end{refe}
(Schopenhauer avait déjà fait ce schéma et Freud ne s’en cache pas)

Les manifestations de l’inconscient :

Les maladies mentales/psychique divisées en deux classes : les névroses et les psychoses

\begin{dfn}[Névrose]
	Accumulation de refoulement où le moi est dans l’impasse. Le sujet est conscient (1) de sa maladie et reste ancré dans la réalité (2). Pas de modification de la personnalité (3), pas de perte du moi.
\end{dfn}

\begin{dfn}[Psychose]
	Trop-plein de refoulement, il y a rupture totale entre le moi et la réalité car les souffrances sont trop importantes. Le sujet n’est pas conscient de sa maladie, il y a une perte de contact avec la réalité et une modification profonde de la personnalité. C’est une perte totale du moi et le ça prend toute la place dans le psychisme.
\end{dfn}


Les rêves :
Selon Freud, les rêves sont une sorte de rébus dans lesquels il y a des choses à interpréter.
\begin{refe}
	Freud, \textit{L'interprétation du rêve}
\end{refe}

Contenu manifeste du rêve 
Contenu latent du rêve
Ce que l’on voit dans le rêve, souvent absurde
Sens caché, sens déchiffré du rêve

Le passage du contenu manifeste au contenu latent est appelé interprétation.

“Le rêve est la réalisation d’un désir refoulé” (réalisation=contenu manifeste et désir refoulé=contenu latent)
“Le rêve est la voie royale qui mène à l’inconscient”

Les actes manqués :
\begin{refe}
	Freud, \textit{Psychopathologie de la vie quotidienne}
\end{refe}

Acte manqué = Échec d’une action habituellement réussie

Lapsus = dire quelque chose d’autre que ce que l’on voulait dire => révèle les désirs inconscients selon Freud

Selon Freud, les actes manqués ne peuvent être associés au hasard et s’expliquent par un désir inconscient.

Erreurs de lecture, d’écriture et oublis (de réveil avant examen par exemple) sont des actes manqués.

Le résultat attendu d’un acte manqué est manqué et le résultat effectif est inadéquat.

Pour Freud, rien n’est insignifiant dans les conduites humaines (théorie des actes manqués)


La psychanalyse a montré qu’il pouvait exister de nombreux phénomènes où il y a des contradictions entre ce que je sais de moi via ma conscience et la signification effective, réelle et inconsciente des pensées et comportements.

On ne peut pas lire en soi comme dans un livre ouvert, on ne peut pas se connaître soi-même. Le psychisme humain ne se réduit pas à la conscience, l’inconscient a même une très grande part dans le psychisme humain, et d’autres instances comme le surmoi et le subconscient existent.

Non, on ne peut pas se connaître soi-même, l’introspection de Descartes est impossible.


En tout cas, on ne peut pas se connaître soi-même entièrement mais une connaissance complète et médiate est possible.On peut se connaître partiellement par des médiations.

\section{Une connaissance de soi médiate est possible grâce à autrui}

\subsection{Une forme de connaissance de soi est possible grâce à la cure psychanalytique}

La psychanalyse a, en plus de prouver l’existence de l’inconscient, pour but de mieux se connaître.

Toutefois, la cure psychanalytique est longue et incertaine : “Nous ne pouvons rien promettre”. Cela distingue la psychanalyse de la médecine (effet placebo et mensonge pour mieux guérir).

Le psychanalyste doit diriger les idées et paroles du patient, donner des explications (interpréter des rêves par exemple) et prendre des notes pour permettre au patient de mieux se connaître.

En thérapie, Saison 2 Episode 12 
=> Docteur Dayan et Robin (8 ans)

\subsection{Je prends conscience de ma propre conscience grâce à autrui}

\begin{refe}
	Sartre, \textit{L'Être et le Néant}, partie III, chapitre 1
\end{refe}

En quoi autrui est-il constitutif de la conscience que j'ai de moi-même.

Pour Sartre, autrui est constitutif de la conscience que j'ai de moi-même : "L'enfer c'est les autres", Huis-Clos. Ce qu'on est pour soi, c'est d'accord ce qu'on est médiatement pour autrui. Cela va à l'encontre de la thèse de Descartes (grand I Sqce 3).

Pour illustrer cela, Sartre prend l’exemple de la honte.

1- Le geste honteux vécu comme un pour soi n’est en fait pas perçu comme un geste honteux
2- La présence d’autrui déclenche l’apparition du sentiment de honte 
3- Grâce au regard d’autrui, je peux porter un regard objectif sur moi-même et ainsi me voir comme honteux
4- La honte est reconnaissance

Le jugement peut être positif ou négatif tandis que la blâme est nécessairement négatif.
La morale présuppose la vie en société, en effet, on ne se juge pas seul.

L’apparition d’autrui et le fait qu’il me voie, suivi du fait que j'aperçois qu'autrui me voit, fait qu’il y a médiatement un sentiment de honte.
=> Le jugement et la présence sont intrinsèquement liés.

Grâce au jugement d’autrui, j’effectue un jugement sur moi-même, une réflexion. Autrui est donc un médiateur, il opère une liaison entre moi et moi-même. Je ne peux pas avoir honte tout seul, mais “j’ai honte de moi tel que j’apparais à autrui”.

Autrui me permet de me voir comme un objet et de porter un regard objectif sur moi-même. Le regard d’autrui rend la connaissance subjective que j’ai de moi-même objective.

La 4ème partie est un raisonnement par l’absurde : si ce n’était que le jugement d’autrui et non la vérité sur moi-même, ce regard ne pourrait pas réellement m’atteindre. Ainsi, je reconnais le jugement d’autrui et donc je le partage, c’est donc le même jugement que je porte sur moi-même. C’est la reconnaissance.

Avoir honte, c’est avoir honte de moi tel que j’apparais aux autres.

La présence d’autrui doit être effective mais pas nécessairement réelle et physique, elle peut être une habitude ou une représentation. En effet, on porte toujours autrui en nous. 
L’existence d’autrui est contenue dans la connaissance intérieure que j'ai de moi-même, on ne peut pas séparer les deux. La connaissance de soi ne s’acquiert pas de manière solipsiste mais grâce aux autres. 


Je ne peux échapper au regard d’autrui, la présence d’autrui fait que je dois assumer mes actes, je ne peux pas fuir. Autrui m’oblige à me regarder moi-même.

Pour Sartre, l’inconscient n’existe pas et est une erreur morale qui relève de la mauvaise foi.


\subsection{La psychanalyse existentielle}

Sartre et Freud ont le point commun qu'ils s'opposent à Descartes : la connaissance de soi a besoin d'un médiateur.

\begin{refe}
	Sartre, \textit{L'Être et le Néant}, partie IV, chapitre 2
\end{refe}

Pour Sartre, un désir inconscient est impossible : "Mon désir est forcément la conscience de ces objets comme désirables".
Ceci est lié au fait que la connaissance du désir est la définition du désir sartrien.

Critique épistémologique de l'inconscient : L'inconscient n'existe pas scientifiquement.

Critique morale de l'inconscient : Croire en l'inconscient est dangereux car cela pourrait permettre d'imputer toute action à notre inconscient.

Toutefois, Sartre admet que la psychanalyse permet de faire des choses utiles et souligné ses points communs avec Freud : le sujet n'est pas conditionné par ses gènes, et chaque geste ou action révèle quelque chose sur nous.

"L'homme s'exprime tout entier dans la plus insignifiante et la plus superficielle de ses conduites. Autrement dit, il n'est pas un bout, un tic, un acte humain qui ne soit révélateur."

La solution proposée par Sartre est la psychanalyse existentielle :
Étudier l'Homme dans sa situation
Rejeter toute logique de symbolisation

\section{Définitions complémentaires}
\subsection{La conscience}

\begin{dfn}[Atome]
“Atome” vient du grec “atomos”, qui signifie “indivisible”. L’atome est une particule infiniment petite, c’est la plus petite unité de mesure. C’est un élément matériel insécable et indivisible, se mouvant dans le vide et dont les combinaisons constituent la matière de toutes les choses du monde. L’atomisme est une doctrine matérialiste soutenue par Démocrite au Ve siècle avant J.C.
\end{dfn}

\begin{dfn}[Contemplation]
Non-action qui consiste à réfléchir sur des objets de pensée en ayant pour vocation de se mêler à ces objets, sans les modifier.
\end{dfn}

\begin{dfn}[Corps]
Enveloppe matérielle dans laquelle nous vivons. Siège des fonctions organiques et physiologiques. L’étude de la conscience interroge la relation entre corps et âme. Lequel des deux constitue l’identité personnelle ?
\end{dfn}

\begin{dfn}[Une personne]
La personne désigne l’Homme en tant qu’il est un sujet moral, conscient, responsable et raisonnable. Voir en l’Homme une personne, c’est lui accorder de la valeur absolue, c’est le considérer comme une fin en soi.
\end{dfn}


\begin{dfn}[Culpabilité]
Sentiment douloureux qu’une personne éprouve en prenant conscience d’une faute dont elle s’accuse
\end{dfn}


\begin{dfn}[Phénoménologie]
Courant de pensée né avec Husserl au XIXe siècle qui prône un “retour aux choses-mêmes” (Husserl). C’est l’étude descriptive des phénomènes tels qu’il apparaissent à la conscience, considérés indépendamment de tout jugement de valeur.
\end{dfn}

\subsection{L’inconscient}

\begin{dfn}[Compulsion]
Tendance irrépressible et d’origine inconsciente à répéter certains actes ou certaines conduites. Ce type de conduite peut mener à des angoisses.
\end{dfn}

\begin{dfn}[Hallucination]
Perception par un sujet d’un objet sensible qui n’existe pas dans la réalité (hallucination positive) ou non-perception d’un objet qui existe dans la réalité (hallucination négative)
\end{dfn}

\begin{dfn}[Phobie]
Peur démesurée et irrationnelle d’un objet ou d’une situation précise. Elle se caractérise par une réaction d’angoisse ou de répulsion devant l’objet ou la situation en question. La phobie n’est pas la peur du danger, mais c’est une peur difficilement explicable (ce qui laisse penser qu’elle ne peut se comprendre qu’à la lumière de l'inconscient)
\end{dfn}

\begin{dfn}[Subconscient]
Il ne faut surtout pas confondre inconscient et subconscient. Le subconscient est littéralement ce qui est en dessous de conscient. C’est ce qui n’est pas saisi par la conscience mais qui est à la limite de l’être, aussitôt que l’attention s’y portera.
\end{dfn}

\begin{dfn}[Sublimation]
Satisfaction d’une pulsion dans des domaines socialement acceptables et spirituellement valorisés (art, philosophie, science, actions altruistes, religion …). L’artiste transforme des pulsions internes névrotiques en des manifestations externes à travers la création d’une oeuvre
\end{dfn}

\begin{dfn}[Transfert]
Processus par lequel les désirs inconscients s’actualisent au cours de la cure psychanalytique et s'extériorisent dans la relation avec le psychanalyste
\end{dfn}

\chapter{La justice est-elle naturelle ?}
(Justice et Nature)
Brainstorming de départ :
Nature / Naturel
Justice
inné (s’oppose à l'acquis, au conventionnel, au culturel et au social)
légitime
sans intervention de l’Homme
appartient à la nature à l’inverse de ce qui est artificiel
sens/sentiment de justice
l’institution
action de juger


\begin{dfn}[Naturel]
Qui appartient à la nature, sans intervention de l’Homme, à l'inverse de l’artificiel
Ce qui est instinctif et s’oppose à la raison
Ce qui est inné
\end{dfn}


\begin{dfn}[Justice]
Sens de la justice, sentiment qui implique le respect d’un idéal à suivre qui serait le bien
L’institution, l’application concrète de la justice dans une cité ou un État. Implique la création de lois, de tribunaux, de métiers, etc…
\end{dfn}

Premières pistes de réponse :
OUI
NON
la nature = l'ordre des choses (divin ?)
base naturelle
morale naturelle
justice peut être naturelle pour l’Homme
pas besoin d’une institution
justice = subjective et relative (aux époques, lieux et religions)
cas des jurisprudences => une modification est possible
la nature ne dicte pas toujours le bien
la nature est plus parfaite que la justice
la justice est uniquement l’institution
la justice va contre l’ordre des choses naturelles
la justice est apprise et a besoin d’être appliquée


\begin{rep}[Légitime / Légal]
Le droit naturel, qui se rapporte à la légitimité, est ce qui est naturellement interdit ou ce qu’on doit faire naturellement. Le droit naturel est inné et indépendant de toute loi écrite.
Le droit positif, tout ce qui est légal, est l’ensemble des règles établies par des conventions. Ce sont les lois (généralement écrites).
\end{rep}


Droit positif = Légal    alors que    Droit naturel = Légitime
Désobéissance civile : faire qqch d’illégal au nom de la légitimité

\begin{refe}
	Kant, \textit{Doctrine du droit}, Introduction
\end{refe}

Le droit naturel repose sur des principes a priori, c’est-à-dire avant toute expérience, alors que le droit positif procède de la volonté d’un législateur.

Le droit naturel “appartient à chacun par nature, indépendamment de tout acte juridique”, alors qu’à l'inverse, pour le droit positif “un acte juridique est requis”.

Pour Kant, le seul droit naturel est la liberté, un droit qui “appartient à tout Homme en vertu de son humanité”.

Un exemple de cette différence entre droit positif et droit naturel est la tragédie Antigone.

Si je juge illégitime ce qui est légal, je me réfère à une idée de la justice supérieure à celle de la légalité ou à un quelconque système juridique.

Dire que la justice est naturelle signifie pouvoir se passer de la légalité et que les Hommes auraient un sens de la justice qui ne justifie pas le besoin des lois.

PBQ :
Ou bien la justice ne peut exister sans la légalité, et alors la justice ne peut pas être naturelle ; ou bien la justice peut se passer de la légalité, et alors la justice est naturelle.

\section{La justice n’est pas naturelle, elle naît avec la légalité}

La justice ne peut exister sans les conventions que sont les lois (droit positif), et a besoin d’être instituée : c’est la légalité.
Le sentiment de justice n’est pas inné chez les Hommes, ils ont besoin de l’apprendre.

\subsection{Le sentiment de justice n’est pas inné mais acquis}

\begin{refe}
	Rousseau, \textit{L’Émile}, Livre II (1762)
\end{refe}
→ traité d’éducation (cocasse sachant qu’il a lui-même abandonné ses enfants)

On suit l’évolution d’Émile, un personnage fictif, de bébé à l’âge de ses 20 ans en 5 volumes.

Pour Rousseau (à la différence de Locke), l’éducation doit être avant tout pratique : l’enfant doit faire des expériences en 1ère personne. Il montre comment l'enfant apprend la justice par l’expérience du sentiment d'injustice. Or si ce sentiment s’acquiert, il n’est donc pas inné.

“Avant l’âge de raison, nous faisons le bien et le mal sans le connaître.”

\begin{dfn}[Propriété]
Possession d’un bien considérée comme juste légalement, le bien nous appartient de droit
\end{dfn}

Rousseau cherche à faire ressentir à Émile l’injustice en passant par la notion de propriété.

Émile prend conscience de la notion d'appartenir et la fait de voir “son” jardin saccagé lui cause de la souffrance. La propriété vient ici du travail, de l’effort, de l'investissement…

Retournement de situation : la victime supposée est en réalité coupable d’usurpation

\begin{dfn}[Usurpation]
S’attribuer illégalement la propriété d’un autre
\end{dfn}

Le marché qu’Émile passe avec Robert constitue la genèse du sens de la justice. Émile possède donc des motivations qui vont le conduire à respecter le principe de propriété privée.

La notion primitive de justice “a été inculquée” à Émile par la propriété,. Le sentiment de justice n’est donc pas naturel.

\subsection{Le droit se réduit aux conventions : les seules lois valables sont celles qui sont instituées}

\begin{refe}
	Pascal, \textit{Pensées}, L60
\end{refe}

La justice se réduit à la légalité et les seules lois valables sont les conventions établies entre Hommes, dans quoi ce serait le chaos.

\begin{dfn}[Loi (sens politique)]
Droit positif. Ensemble de règles et de prescriptions qui régissent la vie d'un État et les rapports des Hommes dans la société en énonçant leurs droits et leurs devoirs. Les lois sont établies par une autorité dite "souveraine" qui impose alors la loi comme une obligation.
\end{dfn}


Pas de justice universelle = Pas de droit naturel

Pascal constate que les lois sont relatives aux lieux et aux époques : les lois ne sont donc pas universelles et absolues.
Il raisonne par l'absurde : "l'éclat de la véritable équité aurait assujetti tous les peuples".

Ce qui est considéré comme juste ou injuste varie considérablement selon là où l'on regarde : la justice est contingente et relative aux temps, lieux et sociétés "On ne voit rien de juste ou d'injuste qui ne change de qualité en changeant de climat".
Par exemple, l'inceste et le meurtre ont parfois été autorisés.

Il n'y a pas de norme pour juger ce qui est bien ou mal dans l'absolu.

La seule loi valable est le droit positif

Il est nécessaire de mettre en place des lois pour éviter l'anarchie et le chaos, puisque toute loi naturelle nous est inaccessible ou bien n'existe tout simplement pas.
Les lois ne sont pas là conséquence d'une loi naturelle, sont illégitimes, mais elles sont nécessaires. C'est cette nécessité qui rend le droit positif légitime.

La légitimité des lois vient uniquement de leur légalité. Selon Pascal, ces lois se font passer pour légitime aux yeux du "peuple".

Une illusion nécessaire

La tromperie devient légitime car elle permet d'éviter le chaos : c'est la raison des effets.
Pour Pascal, il est nécessaire de faire croire au peuple que ces lois sont légitimes et respectent le droit naturel puisque si l'effet est utile, la chose est légitime :
"Pour le bien des Hommes, il faut souvent les piper"

Peuple
"Demi-habiles"
Habiles
Illusion
Révolte
"Pensée de derrière"

Le comportement illusoire en apparence est le plus fondé, le plus légitime. Les habiles ont la "pensée de derrière" : ils ont conscience que la loi est illégitime mais font comme si elle l'était ("conservatisme désabusé").

La valeur la plus importante aux yeux de Pascal est l'ordre.

\subsection{La naissance de la justice va de pair avec la naissance de l’État}

La naissance de l'État repose sur un contrat, une convention : la naissance de l'État n'est pas naturelle. Ce sont les théories du contrat social qui ont notamment été développées par Hobbes et Rousseau.

\begin{dfn}[État]
Autorité politique souveraine à laquelle est soumise un groupe d'humains sur un territoire donné
Limites géographique de cet État
Forme spécifique que prend un gouvernement, nature spécifique de régime politique
\end{dfn}

\begin{refe}
	Hobbes, \textit{Le Léviathan}, 1651
\end{refe}

\begin{bio}[Hobbes]
	Philosophe anglais du XVIIe siècle qui s'intéresse à la philosophie politique (Le Léviathan, Du Citoyen) ; considéré comme l'un des penseurs de l'État moderne (État découlant d'un contrat).
\end{bio}

\textit{Le Léviathan} : Pourquoi l'État est-il né ? Pourquoi les Hommes ont-ils accepté de se soumettre à l'autorité supérieure qu'est l'État ?

"État de nature"
Fiction théorique/expérience de pensée qui imagine l'état dans lequel seraient le Hommes sans État
Hobbes fait une description bien spécifique de l'état de nature : "l'Homme est un loup pour l'Homme" (Plaute, "homo hominis lupus (est)"). Les Hommes sont gouvernés par l'instinct de survie, dit de "conservation", et ils font tout pour survivre. Dans cet état, les Hommes sont tous égaux par nature, ils ont tous les mêmes droits, il n'y a pas d'État ni de justice. Dans l'état de nature, c'est la "guerre de tous contre tous" ("bellum omnium contra omnes").

Les 3 caractéristiques de l'état de nature sont la défiance, le désir de gloire et la compétition. L'objectif est d'acquérir l'estime des autres.

Défiance/Méfiance : peur/risque de trouver quelqu'un de plus fort que soi et peur de la mort => initiatives d'attaques => état de violence permanente sans sentiment de justice

"L'Homme vit dans la crainte et le risque continuels d'une mort violente"

Nécessité de créer un État et des lois

Par ce constat, les Hommes consentent pour pouvoir acquérir la paix, la tranquillité et la sérénité. Ils consentent à transférer leurs droits à un souverain unique nommé "Léviathan" (serpent géant qui doit détruire le monde dans la mythologie) : on se soumet alors à l'État comme on se soumet à un monstre hyper puissant et terrifiant. Il impose automatiquement le respect.

Les Hommes acceptent de se soumettre aux lois et de se séparer de leur liberté totale pour gagner en sérénité.

Le passage à l'État relève toujours d'une forme de survie mais collective. L'abandon de ces droits naturels doit être mutuel (principe de consentement), sauf pour le Léviathan. Les droits sont "transmis" à un autre qui doit assurer notre liberté. Cette transmission, c'est le contrat social.

"C'est comme si chacun disait à chacun «J'autorise cet Homme et je lui abandonne mon droit de me gouverner moi-même à cette condition que tu lui abandonnes ton droit»"

Critique
Le problème chez Hobbes est qu'on perd la valeur fondamentale de la liberté. La recherche de la sérénité absolue n'est pas compatible avec la liberté.


\begin{refe}
	Rousseau, \textit{Du contrat social}
\end{refe}
Chez Rousseau, l'Homme est bon par nature (et c'est la société qui le pervertit). Il n'y a pas de sens inné de la justice chez Rousseau.

Rousseau, Discours sur l’origine et les fondements de l’inégalité parmi les Hommes, I

État de nature

La vertu principale de l'Homme est la pitié. L'Homme est perfectible et capable du meilleur comme du pire.

Pourquoi nécessiter un passage à l'État et à l'état civil via le contrat social ?

Raisons ambiguës et peu claires
Une catastrophe naturelle aurait forcé les Hommes à s'associer à cause du manque de ressources et du besoin de travailler. Il faut alors un État pour encadrer cette association.

Rousseau s'interroge sur les sources de l'autorité de l'État.

Les 4 sources de l'autorité de l'État selon Rousseau :

Force 
Sang
Dieu
Convention 
Tyrannique
Despotique
Totalitaire
=> Hobbes
Monarchie
Oligarchie
=> Repose sur une dynastie

Tribal 
Ethnique 
Clanique
Mafieux
Théocratique 
Démocratique 
=> Accord libre entre les citoyens
=> Libre engagement réciproque

=> Rousseau


Contrairement à Hobbes, chez Rousseau, c'est la volonté générale qui exerce la souveraineté et les lois devraient satisfaire tous les Hommes. L'Homme agit en fonction de l'intérêt public et l'Homme doit avoir le pouvoir législatif. Les Hommes restent donc libres puisqu'ils obéissent à la loi qu'ils ont choisie.
=> Obligation (et plus contrainte comme chez Hobbes)

Passage de l'état de nature à l'état civil :
État de nature
État civil
Liberté naturelle
Liberté politique 
Égalité naturelle
Égalité politique
Amoralité
=> Instinct
Moralité, sentiment de justice
=> Raison
Contraintes
Obligations
L'Homme obéit à des "impulsions physiques"
L'Homme obéit à la "voix du devoir"
Possession
Propriété 


Ce passage à l'état civil produit un changement très remarquable : "l'impulsion du seul appétit est esclavage et l'obéissance à la loi qu'on prescrit est liberté".
L'Homme n'est véritablement libre qu'à l'état civil. Il ne faut pas sacrifier sa liberté à la sécurité, la sécurité sans liberté n'est pas une vie souhaitable (prison haute sécurité par exemple).

"On y vit tranquille aussi dans les cachots mais est-ce assez pour s'y trouver bien ?"

Un peuple qui sacrifié sa liberté pour sa sécurité est qualifié par Rousseau de "peuple de fous".

Régimes et gouvernements :

Nombre de personnes ayant le pouvoir
Pouvoir législatif → Type de régime
Pouvoir exécutif → Type de gouvernement 
Un(e)
Dictatorial
Despotique → Hobbes
Dictatorial
Despotique → Hobbes
Plusieurs
Aristocratique
Oligarchique
Aristocratique
Oligarchique → Rousseau
Tous(tes)
Démocratique → Rousseau
Démocratique 


La création de l'État et de la justice n'est pas naturelle mais humaine. Elle doit être réglée pour respecter le principe de la liberté.

\section{La justice est naturelle : elle nous dicte ce qui est juste}

\begin{refe}
	\textit{Déclaration des Droits de l'Homme et du Citoyen} (1789)
\end{refe}

Dans le préambule, il est cité : "des droits naturels, inaliénables et sacrés" ; "Les hommes naissent et demeurent libres et égaux en droit" (article 2) et "Droits naturels et imprescriptibles de l'homme" (article 2).
=> Il y aurait des droits naturels pour l'Homme : liberté et égalité

Il y aurait une naturalité (possédés dès la naissance, innés) et une inaliénabilité (enlever ces droits revient à enlever ce qui nous constitue) des droits de l'Homme.

\subsection{Il existe une justice naturelle et absolue qui vient de Dieu}

\begin{refe}
	Malebranche, \textit{Entretiens sur la métaphysique et la religion} (1688), 8ème entretien
\end{refe}

\begin{bio}[Malebranche]
	Philosophe français (fin XVIIe - début XVIIIe) rationaliste
\end{bio}

\textit{Entretiens} : La justice est un ordre de la nature qui vient de Dieu. Malebranche s'oppose ici aux théories du contrat social. Dieu indique ce qui est juste, il est la mesure du juste.

Critique des théories du contrat social

Les théoriciens de l'état de nature et du contrat social sont dits "insensés".

"Tout n'est naturellement permis à tous les Hommes" => Critique directe de Rousseau. Ce serait "peindre la société humaine comme une assemblée de bêtes" et animalier l'Homme.
Pourquoi donc ?
Dieu interdit naturellement que tout soit permis à tous les Hommes dans l'état de nature.

Il y a un ordre naturel dicté par Dieu

Cet ordre est dit "immuable", il a toujours été le même.

"Je conviens que l'ordre immuable de la justice est une loi dont Dieu même ne se dispense jamais et sur laquelle il me semble que tous les esprits doivent régler leur conduite"

Il faut suivre le droit naturel dicté par Dieu et pas le droit positif.

Problème : Comment avoir connaissance de ce droit ?
Il est inscrit en nous par Dieu d'après Malebranche.

Second problème : Pourquoi le mal existe-t-il alors ?

Les notions de juste et d'injuste ne sont pas des inventions de l'esprit humain mais sont innées, données par Dieu et inscrites en nous.
Selon Malebranche, la perfection de l'ordre de la nature règne en nous.

La philosophie qui consiste à défendre la perfection de l'ordre de la nature ainsi que l'existence de Dieu, bien que le mal existe, se nomme théodicée.
En effet, selon les auteurs de théodicées on ne peut pas juger Dieu ni comprendre son dessein caché.

\subsection{La véritable justice est celle de la nature que les lois ont corrompue}

Le droit positif serait une corruption.

\begin{refe}
	Platon, \textit{Le Gorgias} (483d-484b), Discours de Calliclès
\end{refe}

Discours de Calliclès : Socrate, Gorgias (un rhéteur), Polos et Calliclès sont présents. Il s'agit d'un discours sur la rhétorique.

Pour Socrate, l'art oratoire n'est qu'un "simulacre" de l'art politique.

La conception de la justice énoncée ici est similaire à la thèse de Thrasymaque dans La République, Livre I

En quoi pour Calliclès la justice est-elle naturelle ?

Les nombreuses marques dans la nature qui montrent que le plus fort domine le plus faible

Les Hommes n'ont pas la même valeur (s'oppose au principe d'Égalité). Le critère de cette valeur est la "capacité" (physique, morale ou intellectuelle). Il s'agit d'une description de la loi du plus fort.

Pour Calliclès, il est juste que les plus forts dominent les plus faibles.

Tandis que chez Hobbes l'égalité mène à la guerre de tous contre tous, chez Calliclès, la loi du plus fort instaure un certain ordre. Pour Calliclès, il y a bien de la justice à l'état de nature et c'est même le seul moment où il y en a.

4 exemples de domination :
Règne animal (lion par exemple)
Cités des Hommes = Assemblées d'Hommes
Famille (le père est traditionnellement le chef de famille)
Xerxès, roi des Perses, fils de Darius 1er => Ont mené les guerres médiques d'invasion de la Grève par la Perse
Chaque exemple va contre une supposée égalité naturelle des Hommes. La justice naturelle renie l'égalité entre les Hommes.

Mais un exemple n'est qu'un constat et sûrement pas un argument. Ce sont des cas particuliers qui n'ont aucune valeur de preuve.

La justice instituée est antinaturelle et asservit les Hommes forts

Ici, la définition de "nature" renvoie à l'ensemble des qualités qui définissent un être, son essence, ses caractéristiques premières.

La nature du juste est ici un principe normatif qui regroupe les normes établies. C'est l'essence du juste.
La conduite de Xerxès (et autres exemples) est conforme à la nature du juste, elle respecte ce qui doit être d'après les normes de justice de la nature.

Il y aurait une mesure du juste indiquée par la nature à laquelle il faudrait se conformer.

Par conséquent, ce qui est juste n'a pas besoin d'être établi puisqu'il est déjà donné par la nature elle-même : le droit positif se trompe, nous nous trompons. Le droit positif n'est pas conforme à la nature du juste.

Le droit positif affaiblit les Hommes forts, c'est injuste. L'égalité devient alors une injustice.

Le tour de force de Calliclès est donc de dire que l'égalité est injustice et je l'inégalité est justice

"Incantations, sortilèges, formules, sorcelleries" pour décrire les lois instaurées
=> Côté péjoratif, mystique, surnaturel et illusoire de ces lois

Les plus forts sont muselés, maîtrisés et contrôlés par la masse des Hommes faibles via les lois : il y a un renversement des forces.

La tentative d'un Homme fort pour rétablir le droit naturel 

La solution imaginée par Calliclès pour inverser ce premier renversement est que le plus fort des Hommes forts lutte contre ces lois, les foule aux pieds. C'est une insurrection, une révolte (différent de la désobéissance civile car d'intérêt privé et pas d'effet de groupes sauf selon Calliclès bien sûr).

Calliclès souligne ici la nécessité de recourir à la justice naturelle qui avait été souillée et remplacée par le droit positif. Ce retour implique la révolte d'un Homme fort.

La loi positive est faillible car il suffit d'un seul Homme fort pour la renverser. Elle n'est donc pas absolue.

Cette thèse se rapproche de celle de Nietzsche (Généalogie de la morale, Traité I) où Nietzsche oppose les Hommes forts, les "Nobles", aux Hommes faibles assemblés, le "Troupeau", qui fait la "révolution des faibles".

La nature nous indique ce qu'il est juste de faire, elle montre que l'inégalité entre les Hommes est juste et que le droit positif est fait par et pour les plus faibles afin d'humilier les plus forts et de les asservir.

\subsection{L’Homme est naturellement fait pour s’associer et vivre en société}

\begin{refe}
	Aristote, \textit{La Politique}, I, 2
\end{refe}

\begin{bio}[Aristote]
	Philosophe grec du IVè siècle avant J.C, disciple de Platon, a fondé le Lycée (son école)
\end{bio}

“L’Homme est par nature un animal politique”

\begin{dfn}[Cité]
Communauté politique organisée, une société avec des lois et valeurs communes
\end{dfn}

\begin{dfn}[Politique]
Tout ce qui concerne la vie de la cité.
Art de gouverner la cité ou l'État.
\end{dfn}

“Toute cité est dans la nature” : selon Aristote, le regroupement en cité est le dernier stade de la perfectibilité des Hommes

Lorsque l’Homme se regroupe en société, il atteint son “développement complet”. En effet, atteindre ce développement complet, c’est pouvoir se suffire à soi-même pour survivre et vivre. L’Homme ne peut atteindre ce stade qu’en s’associant en cité avec d’autres Hommes.

Si un Homme n’appartient à aucune cité, il est qualifié de “créature dégradée” par Aristote.

\section{La justice est ce qui vient corriger la nature, même si elle doit toujours s’appuyer sur le droit naturel}

La justice est certes instaurée, elle commence avec la législation, mais elle s’appuie sur le droit naturel et doit rester légitime.

La justice serait alors une perfection du droit naturel et permet de rendre ce dernier applicable.

\begin{rep}[En fait/En droit]
Est dit en fait ce qui est dans la réalité (à rapprocher de “en pratique”).
Est dit en droit ce qui doit ou devrait être, indépendamment de la réalité (à rapprocher de “en théorie, en principe”).
Le passage du “en droit” au “en fait” est permis par la justice.
\end{rep}

=> Calliclès confond "en fait" et "en droit"

\subsection{La justice comme institution s’oppose à la loi naturelle que serait la vengeance}

\begin{refe}
	Hegel, \textit{Principe de la philosophie du droit}, Partie I, section 3, 1820
\end{refe}

\textit{Principes de la philosophie du droit} : Hegel fait une théorie juridique, économique et sociale. La section 3 se nomme "l'injustice".

Hegel parle ici de la justice corrective (i.e. celle des tribunaux) qui, gérée par les magistrats, établit la justice dans l’Etat. Elle s’oppose à la loi du talion (œil pour œil, dent pour dent).

\begin{dfn}[Punition]
Peine infligée à quelqu’un pour une infraction, un délit ou un crime dont il est jugé responsable
\end{dfn}

\begin{dfn}[ Vengeance/Vendetta]
Action par laquelle une personne lésée ou un proche inflige en retour un mal à l’offenseur pour le punir
\end{dfn}

La justice instituée est créée pour contrer la pseudo-loi naturelle qu’est la loi du talion.

La loi du talion serait un “droit immédiat”, naturel, qui n’a pas besoin de médiateur.

La personne lésée réagit avec ses passions et non avec sa raison : la vengeance n’est pas objective ni impartiale. Hegel parle de volonté particulière.

Au contraire, pour que le jugement soit impartial, il faut l’action d’une volonté objective, un tiers qui n’a rien à voir, c’est le droit médiat”.

Dans le cas de la vengeance, la peine délivrée est contingente et arbitraire, dans le sens où elle aurait pu être autre. 

Au contraire, dans le cas de la justice, la punition décrétée est censée être nécessaire et objective dans le sens où la même infraction doit toujours correspondre à la même peine.

La justice n’est pas naturelle, elle est instituée et s’oppose à la pseudo-justice de la vengeance. La justice doit être impartiale et délivrer des peines nécessaires.

De plus, la vengeance instaure un cercle vicieux qui ne se termine pas : “un processus de l’infini qui se transmet de génération en génération”. C’est ça la vengeance selon Hegel.

\subsection{La justice est ce qui vient corriger la nature}

\begin{refe}
	Rawls, \textit{Théorie de la justice} (1971)
\end{refe}

\begin{bio}[Rawls]
	Philosophe américain du XXe siècle (mort en 2002) ; il met en place une nouvelle conception de la justice ; il s’intéresse à savoir comment et sur quels principes fonder une justice ; il va faire un nouveau contrat social
\end{bio}

\begin{rep}[Origine/ Fondement]
L’origine, c’est le commencement historique, temporel, le début de l’existence d’une chose.
Le fondement implique une légitimité, c’est ce qui va fonder, justifier l’existence de la chose. C’est ce qui soutient l’existence d’une chose, le pilier.
Ainsi, ce qui est fondé est légitime mais ce qui a une origine ne l’est pas nécessairement.
Par exemple, l’origine de la monarchie est l’instauration historique, temporelle de la monarchie tandis que son fondement est le droit divin, la monarchie tire(rait) son pouvoir de Dieu.
\end{rep}


Pour Rawls, il y a à l’origine de l’injustice. Cette injustice doit être réparée à posteriori par un principe de redistribution et de compensation (1er principe).

Les inégalités sociales et économiques doivent être agencées de sorte qu'elles soient au bénéfice des plus démunis.

“Des inégalités de richesses et d’autorité sont justes si et seulement si elles produisent en compensation des avantages pour chacun, et en particulier pour les membres les plus désavantagés de la société.”

Pour Rawls, on peut établir une justice “satisfaisante” en veillant à ce que les plus démunis tirent un certain bénéfice des inégalités.

Le deuxième principe pour Rawls est que certains droits doivent être fondamentalement égaux, il ne doit pas y avoir d’inégalité pour ces droits-là (vote, culte, pensée …).

Comment appliquer le premier principe ? => La méthode du voile d’ignorance

On essaye de se mettre un voile devant les yeux pour ignorer quelle sera sa situation et, à ce moment-là, mettre en place  les principes de la justice.

En effet, chacun aura tendance à favoriser sa propre situation.

C’est une méthode hypothétique qui ne peut fonctionner dans la réalité, mais Rawls nous invite tout de même à réaliser cette expérience de pensée.

Avec le voile d’ignorance, les Hommes oublient leur point de vue personnel et leur situation. Ils ne pourraient pas alors favoriser leur situation puisqu’ils ne la connaissent pas, et c’est donc là qu’il faudrait fonder la justice. Cela permettrait de garantir une impartialité dans la redistribution.

\subsection{Ouverture : C’est à la justice de protéger la nature}

\begin{dfn}[Écologie]
	Théorie qui souhaite une meilleur adaptation de l’Homme à son environnement naturel dans le but d’atteindre une sorte d’harmonie avec la nature
\end{dfn}


\begin{refe}
	Hans Jonas, \textit{Le Principe responsabilité} (1979)
\end{refe}

\begin{bio}[Hans Jonas]
	Philosophe allemand du XXe siècle (mort en 1993) ; élève de Heidegger ; s’intéresse à la vie de l’humanité future
\end{bio}

Jonas part du principe que la technique et la théorie de la croissance sont devenues une menace pour la nature et l’humanité. Il faut alors créer de nouvelles règles morales et éthiques pour repenser l’éthique et contrer ces menaces.

\begin{dfn}[Éthique]
	L'éthique, c’est une réflexion morale appliquée à une situation concrète. L’éthique est une application de la morale. Une éthique est une conduite morale particulière.
\end{dfn}



Finalisme
Mécanisme
Doctrine d'Aristote selon laquelle tous les mouvements naturels répondent à une fin. Aristote tente d’expliquer la nature par ses “causes finales”.
Doctrine selon laquelle la nature obéit à des lois régulières, comme une machine. Grâce à la raison, nous pouvons connaître ces lois (cf. Descartes et le rationalisme) 
Doctrine dite téléologique : tout obéit à une certaine fin
Bergson s’oppose à ce mécanisme


Note : On parle plus de morale kantienne que d’éthique kantienne car on reproche à Kant l’absence d’exemple d’applications concrètes de se amorale.

Jonas va proposer un nouvel impératif pour remplacer l’impératif catégorique de Kant : 
“Agis de façon que les effets de ton action soient compatibles avec la permanence d’une vie authentiquement humaine sur Terre”

L’Homme a une responsabilité envers les générations futures et la planète.

\section{Définitions complémentaires}
\subsection{La justice}

\begin{dfn}[Coup d’Etat]
Appropriation du pouvoir politique par la force. C’est une mise en suspension des lois et de la légalité, une action contre le droit commun.
\end{dfn}

\begin{dfn}[Équité]
Modification de la justice écrite pour essayer de mieux correspondre à l’esprit de la justice. La loi, en tant qu’elle est générale, peut parfois devenir injuste au regard de certains cas particuliers. C’est là qu’intervient l’équité, qui modifie la loi pour l'appliquer aux cas particuliers. Ce principe est théorisé par Aristote dans le Livre V de l’Ethique à Nicomaque.
\end{dfn}

\begin{dfn}[Fascisme]
Système politique autoritaire érigé par Mussolini et caractérisé par la tout-puissance de l’Etat et par l’exaltation de la nation (ce qui passe par le culte, un parti unique, l’importance du militaire …)
\end{dfn}

\begin{dfn}[Libéralisme]
Doctrine politique et économique prônant la non-intervention ou la limitation de l'intervention de l’Etat et voulant ainsi donner le plus de liberté possible aux individus. Le libéralisme favorise la libre concurrence et ouvre la porte au non-respect des droits fondamentaux.
\end{dfn}


\begin{dfn}[Raison d’État] 
Principe justifiant des pratiques de l’Etat au nom de l’intérêt de l’Etat et de la nation. C’est ce qui permet à l’Etat de se conserver : la raison d’Etat assure la pérennité du gouvernement et permet à l’Etat d’avoir l’avantage sur la morale et le droit.
\end{dfn}

\begin{dfn}[Socialisme]
Doctrine politique qui affirme la priorité du bien général de la société, qui doit être contrôlé et assuré par l’Etat. Elle incline vers une conception plus égalitaire de la société et a pour fondement des sentiments humanitaires.
\end{dfn}

\subsection{La nature}

\begin{dfn}[Biologie]
Science qui étudie le vivant. Terme apparu au début du XIXè siècle, dans le livre Théorie de l’évolution des espèces de Lamarck, notamment avec le début de l’étude de la génétique (mise en évidence des lois qui règlent la reproduction et la transmission héréditaire) et l’essor des théories de l’évolution. La biologie peut poser des problèmes éthiques : peut-on étudier le vivant si cela implique des expérimentations sur le vivant
\end{dfn}

\begin{dfn}[Cosmos]
Univers en tant qu’il est considéré comme un ensemble ordonné
\end{dfn}

\begin{dfn}[Environnement] 
Ensemble des éléments et des phénomènes physiques qui entourent un organisme vivant et peuvent interagir avec lui.
\end{dfn}

\begin{dfn}[Mouvement]
Au sens large, il s’agit de tout changement de substance : le changement spatial, le changement qualitatif, le changement quantitatif, la génération ou la corruption. Au sens restreint, il s’agit uniquement du changement de position dans l’espace, le déplacement d’un corps par rapport à un point fixe, et à un moment déterminé. Tout mouvement s’observe grâce à des repères.
\end{dfn}

\begin{dfn}[Taxinomie]
Science des lois de la classification. La perspective biologique aristotélicienne de vouloir répertorier toutes les espèces existantes en indiquant leurs caractéristiques spécifiques et communes est une entreprise de taxinomie.
\end{dfn}

\begin{dfn}[Le vivant]
Ensemble de ce qui vit, c’est-à-dire de ce qui naît, croît, se nourrit, se meut, décroît et meurt. Les êtres vivants sont des organismes, c’est-à-dire des êtres dont toutes les parties sont interdépendantes.
\end{dfn}

\chapter{L’art se réduit-il à la technique ?}
(Art et Technique)
\section{L’art se réduit à la technique : l’artiste est un genre d’artisan}

\subsection{Les deux utilisent des outils et instruments}


\begin{dfn}[Outil]
	Objet technique conçu et fabriqué pour un usage manuel. La particularité de l'outil est qu’il permet à l’Homme d’agir sur la matière, impliquant l’idée que l’Homme agirait différemment sur la matière sans cet outil. L’outil a une forme et une matière qui sont adaptés à la matière qu’il travaille (cf. texte Aristote corpus Travail). L’outil est le prolongement de la main et permet d’armer le corps humain.
\end{dfn}



\begin{dfn}[Instrument]
	Objet technique qui sert à mesurer, observer quelque chose. Sans l’instrument, il serait impossible d’obtenir ces informations (lunette astronomique par exemple). L’instrument délivre l'information sans agir sur la matière.
\end{dfn}


L’artiste et l'artisan utilisent tous les deux des outils pour modifier la matière. Ils opèrent donc une modification de la matière. Par exemple : le burin pour le sculpteur, le pinceau pour le peintre, etc…

\begin{refe}
	Bergson, \textit{L’évolution créatrice}, 1907
\end{refe}

Pour Bergson, ce qui caractérise l’Homme est l’utilisation d’outils, mais surtout la capacité de fabriquer des outils. En effet, s’il est communément admis que l’Homme est un Homo sapiens (au sens de l’Homme qui pense, et non pas de l’espèce), Bergson ajoute que l’Homme est un “Homo faber”.

L’Homme se caractérise par sa “faculté de fabriquer des objets artificiels, en particulier des outils à faire des outils et d’en varier indéfiniment la fabrication”. Non seulement l’Homme crée des outils, mais la création “d’outils à faire des outils", la méta-fabrication, témoigne de sa place particulière dans le règne animal.

Non seulement l’Homme est capable de fabriquer des outils, mais surtout il a la capacité de fabriquer des outils pour fabriquer des outils (méta-utilisation des outils), qui n’est pas présente chez les animaux.

\subsection{L’art nécessite un savoir-faire et une application des règles}

L’artiste doit connaître des règles de production, et surtout doit savoir les appliquer.
Le talent artistique serait alors acquis : on apprendrait à devenir un artiste.

Quelques règles en art :
Danse : les positions
Peinture et dessin : la perspective
Photographie : le ratio
Canons dans la Grèce Antique, principalement liés à la sculpture :
“khanon” = ensemble de règles qui expliquent comment doit être telle ou telle oeuvre pour être réussie
Musique : gammes, intervalle, accord et orchestration


Peinture : bonne connaissance de la géométrie, de l’anatomie et de l’harmonie des couleurs


Littérature : vocabulaire riche




Il y a deux types de règles en art :
celles de maîtrise technique de base qui sont les mêmes pour tout le monde (ou presque) => colonne 1
celles plus particulières qui inscrivent l’oeuvre dans un courant artistique => colonne 2
Pour appliquer ces règles, l’artiste doit faire preuve d’habileté :

\begin{dfn}[Habileté]
Qualité de celui qui maîtrise les techniques nécessaires (adresse, finesse, intelligence, savoir agir à propos …). Est habile celui qui comprend les règles et sait les appliquer.
\end{dfn}

Dans la Grèce Antique, on appelait canon un critère de réussite d’une œuvre d’art. Notamment, au -Vè siècle, Polyclète (Le Doryphore) établit des rapports de proportion mathématiques pour le corps humain qui doivent être respectés en sculpture afin de créer une sculpture parfaite.

Ceci suppose que la réussite d’une œuvre d’art pourrait être mesurée objectivement.

On retrouve ce respect des règles poussé à l’extrême durant la période classique en France : les jardins de Versailles (Le Nôtre, 1661), le théâtre codifié respectant les trois unités et les deux règles (Molière, Racine, Corneille), la poésie versifiée … 

\begin{dfn}[Perspective]
Façon de représenter la réalité en donnant une impression de profondeur pour faire coïncider l’oeuvre avec le regard du spectateur
Il existe plusieurs types de perspective : linéaire, cavalière, curviligne, atmosphérique …
\end{dfn}


\begin{refe}
	Raphaël, \textit{L’Ecole d’Athènes}, 1508, Musées du Vatican
\end{refe}

Note historique : la perspective linéaire a été inventée par l’architecte italien Brunelleschi au début du XVè siècle.

Ainsi, la perspective se rapproche de l’imitation, et on a longtemps considéré que l’art devait imiter la réalité (du moins pour les arts picturaux). L'œuvre la plus réussie serait alors celle qui imite le mieux la réalité.

\begin{dfn}[Imitation]
Représentation d’un modèle existant dans la réalité
Essayer de faire passer la copie pour l’original
Simple inspiration (sens plus faible)
\end{dfn}

Pour la première définition, on a tous l’image d’un peintre impressionniste avec son chevalet qui tente de reproduire un paysage.
Pour la deuxième définition, la copie serait le représentant et l’original le représenté (à rapprocher du signifiant et signifié pour le Langage).

La copie est un plagiat, une contrefaçon, ou encore, au sens péjoratif, un simulacre. Il s’agit de quelque chose de vraisemblable à rapprocher de la tromperie.

L'œuvre la plus réussie serait alors un trompe-l’œil ?

\begin{dfn}[Trompe-l’œil]
Procédé de représentation visant à créer l’illusion de réalité
\end{dfn}

Pour conclure, comme dans la production technique, l’artiste doit se plier à un certain nombre de règles, qu’elles soient techniques ou plus particulières à une époque. Les meilleurs artistes sont alors ceux qui ont la maîtrise la plus parfaite, au point de savoir imiter à la perfection.

\subsection{L’artiste, comme l’artisan, doit fournir un travail}

\begin{refe}
	Nietzsche, \textit{Humain, trop humain}, IV, paragraphe 155 “Croyance à l’inspiration”
\end{refe}

L’art peut être considéré comme un travail non pas au sens d’une action rémunérée mais au sens d’une action qui nécessite de fournir un effort.
L’artiste doit fournir des efforts et utiliser des qualités comme l'adresse, la finesse etc… Il travaille.

Quelles sont les deux thèses que Nietzsche oppose ici ?

Nietzsche oppose les thèses selon lesquelles l’artiste tire son art de l’inspiration (Platon, Ion ; et le génie selon Kant dans Critique de la faculté de juger) ; et la thèse selon laquelle l’artiste est un travailleur (Nietzsche ici).

Énumérez les actions qui caractérisent le travail de l’artiste (verbes). Quelle est la faculté à l'œuvre ici ? Définissez-là.

Actions de l’artiste : produire, rejeter, choisir, combiner

La faculté principalement utilisée ici est l’imagination (la raison est aussi utilisée).

\begin{dfn}[Imagination]
Faculté de créer des images qui n’existent pas dans la réalité
Faculté de se projeter vers le futur
Faculté de produire des images (au sens large) et d’en faire quelque chose
\end{dfn}



Analysez l’exemple des Carnets : en quoi cet exemple permet-il d’illustrer le fait que l’artiste est un travailleur ?

Carnets : Composition, recherche, tâtonnement, temps, petit à petit, expérimentations, essais, erreurs, succès, synthèse d’un travail

L'artiste utilise des brouillons (comme nous devrions tous le faire en philo d’après la prof). Ainsi, il produit un effort et donc il travaille.

On a l’exemple du repentir en peinture (revenir sur ses traits et coup de pinceaux, avec notamment :
\begin{refe}
	Vélasquez, Le Bouffon aux calebasses, 1639, Madrid
\end{refe}
(calebasse à droite)


Quelles sont les différences entre la production et la reproduction ?

Production : relève de l'inspiration et de la nouveauté => vers le futur
Reproduction : relève de l’imitation => vers le passé

Conclusion :
Contre l’idée que l’artiste est un génie qui n’aurait pas besoin de travailler (thèse de Kant, voir II-A), Nietzsche montre que l’artiste, comme l’artisan, doit fournir un travail conséquent issu de l’imagination et de la raison.

Transition vers le II :
L’artiste pourrait avoir besoin d’une certaine sensibilité qui serait innée et s’opposerait alors à l’artisan.
L’art n'est pas qu’un moyen mais aussi une fin.
L’artisan a beau maîtriser toutes les techniques, il lui manquera toujours une chose nécessaire à produire de l’art.

\section{L’art ne peut se réduire à la technique : l’artiste est autre chose qu’un artisan}

\subsection{L’artiste va au-delà des critères techniques }

\begin{refe}
	Kant, \textit{Critique de la faculté de juger}, Partie I, Section 1, Livre 2 “L’analytique du sublime”
\end{refe}

Note : Kant a écrit trois critiques : Critique de la raison pure, Critique de la raison pratique et  Critique de la faculté de juger

\begin{dfn}[Génie artistique]
Aptitude innée remarquable qui permet de la création d’oeuvres artistiques jugées exceptionnelles, remarquables et inégalables (chefs-d’oeuvre)
\end{dfn}

“Les beaux-arts sont uniquement les arts du génie” : pour Kant, le génie est une “disposition innée de l’esprit”.

Les 3 critères du génie artistique selon Kant :

L’originalité

Le génie n’imite personne :
“Les œuvres du génie ne suivent aucune règle déterminée.”

\begin{dfn}[Originalité]
Caractère de ce qui porte une originalité en soi, et qui ne s'appuie sur aucun modèle 
\end{dfn}

L’exemplarité

Le génie n’imite personne, mais les œuvres du génie doivent être des modèles, des exemples pour ceux qui apprennent la technique :
“Bien qu’eux-mêmes ne procèdent point d’une imitation, ils doivent cependant servir à d’autres de mesure ou de règle d’appréciation.”

\begin{dfn}[Exemplarité]
Caractère de ce qui, par sa perfection, peut être donné comme un modèle à suivre et/ou à imiter
\end{dfn}

=> Il y a donc deux types d’exemples : l’exemple lambda, le cas particulier pris au hasard ; et l’exemple à suivre, le modèle quasi-parfait.

L’incommunicabilité

Le génie est incapable d’expliquer comment son œuvre naît, il ne le sait pas lui-même :
“Le poète ne peut indiquer comment ses idées poétiquement riches surgissent et s’assemblent dans son cerveau.”

\begin{dfn}[Incommunicabilité]
Caractère de ce qui ne peut être transmis à autrui (et même à l'artiste lui-même), de ce qui n'est pas accessible à autrui
\end{dfn}



L’artiste crée, il réalise une chose qui n’existait pas auparavant, une chose originale. L’artiste part du réel et le modifie, il ne crée pas ex nihilo (ex + ablatif = à partir de, hors de ; et nihilus,i,m = rien, le néant).
La production, ou fabrication, c’est mettre au monde, réaliser un objet technique ayant un usage déterminé, une fin précise, en suivant des règles spécifiques.

Création
Production, fabrication
Artiste
Artisan, technicien
Originale
Copie
Pas de création ex nihilo (seul Dieu crée ex nihilo)





\subsection{Le processus artistique est différent de celui de la technique}

L’artiste, à la différence de l’artisan et du technicien, n’a pas de conception précise de son œuvre avant de la réaliser.
L'œuvre prend forme au fur et à mesure de sa conception. L'œuvre obtenue est généralement différente de la conception initiale qu’en avait l’artiste, voire surprenante.

L’artiste est surpris par son œuvre, l'œuvre d’art échappe à son créateur (différentes interprétations possibles par exemple).

Il n’y a pas de planification dans l’art, ou du moins pas autant que dans la technique.

\begin{refe}
	Alain, \textit{Système des Beaux-Arts}, Livre I, Chapitre 7 “De la matière”
\end{refe}

\begin{bio}[Alain]
	De son vrai nom Émile Chartier, philosophe français de la fin du XIXe et du début du XXe, invente le style des Propos
\end{bio}

Technique
Art
Oeuvre mécanique => pas de réflexion, travail non-libre
Oeuvre libre => Réflexion, action et esprit
Production
Création
Maîtrise de l’oeuvre
Pas de maîtrise de l’oeuvre
Plusieurs exemplaires
Oeuvre unique
Idée bien définie de son oeuvre en tête
Pas d’idée bien définie en tête (ou floue, imprécise)


Conception/Réalisation
Idée
En second dans la technique
En premier dans la technique
En premier dans l’art
En second dans l’art
Dans la réalité
Dans l’esprit
Réalité matérielle
Réalité spirituelle
Pratique
Théorie


Avant de fabriquer son objet, l’artisan ou le technicien a en tête un plan préconçu et préexistant, une sorte de prototype, et toute l’exécution consiste à réaliser ce plan. Plus la réalisation est adéquate, plus l’objet est réussi.

L’artiste est pour Alain “spectateur de son œuvre, mais n’en agit pas moins pour autant. Il est spectateur en tant qu’il voit des choses. Exemple du peintre : le peintre voit quelles couleurs vont bien aller ensemble.

Il n’y a pas de protocole ou de recette pour créer la beauté ou créer une œuvre d’art belle.

\subsection{L’art et la technique ont des fins différentes}

Tandis que les objets techniques sont un moyen en vue d’une fin déterminée et pratique (vitale ou autre), l'œuvre d’art n’est pas un moyen et n’a pas de fin précise : l'œuvre d’art est désintéressée.

Cependant, l’art n’est pas pour autant inutile : il permet l’évasion, l’élévation, la sublimation, la réflexion, l’enrichissement, l’éducation…

Attention à la nuance entre l’utile et l’utilitaire : l’utile (art et technique) relève de la vie, alors que l’utilitaire (technique seulement) relève de la survie.

Heidegger réalise cette distinction à propos du tableau :
\begin{refe}
	Van Gogh, Les Souliers, 1886, Amsterdam
\end{refe}

\begin{dfn}[Utilitaire]
Ce qui répond à des besoins de survie ou à une fonction pratique et/ou économique
\end{dfn}

\begin{dfn}[Utile]
Ce qui a un intérêt pour la vie sans que ce soit nécessairement les fonctions suc-citées et pour répondre à des besoins de survie
\end{dfn}

L’art répond donc à des besoins dépassant la simple survie.



\begin{refe}
	Kant, \textit{Critique de la faculté de juger}, Partie I, Section 1, Livre 1 “L’analytique du beau”
\end{refe}

\begin{dfn}[Beauté]
Caractéristique d’une chose (ici d’une oeuvre d’art) qui par sa perception procure une satisfaction et suscite une émotion
\end{dfn}

\begin{dfn}[Esthétique (nom)]
La science du beau
La philosophie de l’art (master d’esthétique = master de philosophie de l’art)
Appréciation personnelle de la beauté
Du grec “aisthesis” qui signifie “sensation”
\end{dfn}

\begin{dfn}[Goût]
Capacité à juger de la beauté (d’une oeuvre d’art ici)
\end{dfn}

\begin{dfn}[Jugement esthétique / Jugement de goût]
Jugement de la beauté, savoir si une chose est belle ou non
\end{dfn}

Le jugement de goût / le jugement esthétique chez Kant :
Il est dit désintéressé : il procure une satisfaction désintéressée.

Tandis que la perception de la beauté est désintéressée, l’agréable est perçu de façon intéressée. Le beau plaît à l’esprit alors que l'agréable plaît aux sens, ce qui montre que la perception de l’agréable est intéressée car elle satisfait nos sens facilement. L’agréable ne fait pas intervenir l’esprit mais seulement nos sens, alors que le beau fait intervenir l’esprit et les sens. Le jugement esthétique est donc un jugement désintéressé qui plaît à l’esprit, il est par conséquent libre.

Le jugement de goût est subjectif mais il a une “prétention à une universalité subjective” : il n’y a pas d’universalité du jugement de goût dans les faits, mais nous sommes en bon droit d’exiger que ce jugement soit le même pour tous.

On peut rapprocher cette volonté d’universalité de l’universalité de la loi morale kantienne.

Cette prétention à l’universalité n’est qu’une “idée”.

“Est beau ce qui plaît universellement sans concept”

Conclusion du grand 2 : Les critères de l’oeuvre d’art
Originalité → Génie
Exposition dans des musées → Reconnaissance publique de l’artiste
Doit plaire à l’esprit → Être belle
Pas de planification
Libre et désintéressée
Doit faire réfléchir
Intemporalité
Utilité
Unicité
\begin{refe}
	Weber, \textit{Le Savant et le Politique}
\end{refe} “Une oeuvre d’art vraiment “achevée” ne sera jamais surpassée et ne vieillira jamais” (par opposition à la technique et à la science)

L’oeuvre d’art, une notion problématique :
Vers la fin du XVIIIe, et durant une grande partie du XIXe, une définition de l'œuvre d’art se met en place et semble cohérente. Mais à peine établie, cette définition est remise en question par les artistes eux-mêmes. L’art contemporain la conteste.
Durant une grande partie de l’histoire, l’oeuvre d’art n’est pas isolée des autres productions
Spécificité de l'œuvre d’art au sens classique (fin XVIIIe, début XIXe). La notion d'œuvre d’art est la référence absolue de l'œuvre d’art.
Remise en cause du concept d'œuvre d’art dans l’art contemporain. Contestation de la notion d'œuvre au nom de la création artistique.
Pas d'opposition tranchée entre objets esthétiques, objets sacrés, objets ornementaux, outils…
Spécificité de l'oeuvre d'art, par opposition à la production en série des objets artisanaux et industriels
Répétitions (pop art) ; séries ; esthétique industrielle.
Modification permanente de la culture orale ; reprise-modification continuelle des oeuvres ancestrales ; chaque artiste-artisan imite et modifié ses prédécesseurs 
Clôture, achèvement de l'œuvre.
Perfection interne.
Protection juridique de l'œuvre.
Un artiste = une oeuvre = une date
Non finito, esquisse ; “œuvre en progrès” ; l’inachèvement et la répétition comme vitalité interne ; l'œuvre dépendant de ses “lecteurs” ; l'œuvre et ses variantes infinies.
Fonctions religieuse, politique, magiques de l'œuvre. L'œuvre est au service d'un but social. Elle n'existe pas pour elle-même ; elle est outil.
Finalité interne : l'œuvre n'a pas d'autre utilité qu'elle-même.
L'art pour l'art.
La beauté gratuite.
L'œuvre comme arme politique, transgression, subversion de l’ordre social (dadaïsme, surréalisme …).
Horizon restreint de la plupart des cultures : on veut marquer l'identité de son peuple en produisant des objets durables ; mais parler au nom de l'humanité n'a pas de sens.
Permanence, intemporalité de l'œuvre. Elle est faite pour s'installer dans la durée. L'œuvre est marquée par la volonté de laisser une trace immortelle. Elle est censée parler au nom de l'humanité entière.
Oeuvres éphémères : body art, land art, “emballages” de Christo
Production artisanale ou technique ; reproduction d'éléments traditionnels. Le gest technique ne se sépare pas du sens esthétique.
"Création" ; originalité, nouveauté ; non-reproduction
Mise en scène : ready-made (Duchamp), utilisation d’objets courants, détournement du quotidien. Reproduction (pop art).
Anonymat des concepteurs ; intervention des commanditaires
Caractère individuel de l'œuvre :
Une œuvre = un artiste = une date
Art populaire, art brut, expériences de production collectives
Caractère artisanal, technique de la production : l'art n'est pas séparé du monde social.
Sérieux de la production. Sacralisation de l'art : l'art comme nouvelle religion dans un monde sans religion.
Autodérision, gags, provocation. Contestation des musées, des institutions
Art et décoration ; art et tatouages ; art et mime 
Matérialité de l'œuvre ; caractère objectivité : une matière, une toile …
Concept art, land art, happening, productions visuelles …
Les matériaux sont utilisés pour leur valeur présumée intrinsèque : l'or est précieux, donc une statue ouvragée d'or ne peut qu'être belle
Noblesse des matériaux en eux-mêmes (bois, pierre) : le sculpteur respecte la pierre ou le bois comme le chasseur respect telle ou telle race de chiens.
Usage de matériaux de récupération, de matériaux pauvres, de déchets. Idée de transfiguration : toute matière peut être anoblie par le geste qui l’utilise.
Intégration des productions dans la vie sociale, religieuse, politique.
Rapport distant objet/spectateur.
Relations physiques avec les spectateurs, intervention des spectateurs …
Caractère traditionnel de l'inspiration.
Caractère réfléchi, maîtrisé du processus créateur.
Rôle du hasard ; action painting ; écriture automatique


Critique des critères de l’oeuvre d’art 
L’art conceptuel (Kosuth, One and Three chairs, 1965, Paris)
Le pop art
Les “happenings” (performances) → Manzoni “Consommation de l’art dynamique par ses spectateurs mêmes, dévorateurs d’art”
La reproductibilité technique → Pop art, sérigraphies, CD, streaming
\begin{refe}
	Benjamin, \textit{L’œuvre d’art à l’époque de sa reproductibilité technique}
\end{refe}

L’art contemporain déstabilise et remet en question les critères classiques de l'œuvre d’art.

\section{L’art tente de pallier les défauts de la technique}

\subsection{Les dangers de la technique}

\begin{refe}
	Descartes, Discours de la méthode, Partie 6
\end{refe}

Descartes défend ici que le progrès technique et scientifique permet d’améliorer la condition humaine puisqu'il nous permet de maîtriser la nature. Le progrès technique permet de nous rendre “comme maîtres et possesseurs de la nature”.

\begin{refe}
	Louis de Broglie, \textit{Physique et Microphysique}, Conclusion, 1947
\end{refe}

\begin{bio}[Louis de Broglie]
	(Prononcer “de Breuil”) Physicien français ; prix Nobel de physique en 1929 « pour sa découverte de la nature ondulatoire des électrons »
\end{bio}

La science se décline en deux sous-parties :
la science pure et désintéressée qui relève de la théorie
la science pratique et appliquée qui relève de la pratique

La technique n’est pas dangereuse en elle-même, mais c’est l’usage qu’on en fait qui peut être dangereux. Par exemple, la chimie peut sauver des vies (médicaments) mais peut aussi en détruire (gaz mortels durant la Première Guerre mondiale).

L’Homme élève sa sagesse au niveau de sa puissance technique en utilisant sa raison.

La raison possède elle aussi deux aspects :
celui qui permet de distinguer le vrai du faux, la raison rationnelle
celui qui permet de distinguer le bien du mal, la raison raisonnable

Ainsi, selon de Broglie, il faut avoir un usage de la raison qui ne doit pas qu'être rationnel, mais qui doit aussi être raisonnable.

"Toute augmentation de notre pouvoir de d’action augmente nécessairement notre pouvoir de nuire. Plus nous avons de moyens d’aider et de soulager, plus nous avons aussi de moyens pour répandre la souffrance et la destruction." -> Retenir l’idée

Malheureusement, dans les faits, l’aspect raisonnable n’est pas toujours utilisé en sciences et en technique.

\subsection{L’art réenchante le monde que la technique a désenchanté}

\begin{refe}
	Weber, \textit{Le Savant et le Politique}, “Le métier et la vocation du savant”, 1919
\end{refe}

\begin{bio}[Max Weber]
	Sociologue et économiste allemand, il est considéré comme le père de la sociologie moderne (avec Emile Durkheim) ; il a travaillé principalement sur la politique ; fin XIXe - début XXe
\end{bio}
Le désenchantement par la technique :

Weber constate qu’il y a en Occident une rationalisation et une “démystification du monde” (à rapprocher de Nietzsche et de la phrase “Dieu est mort”). Il constate qu’à son époque, le spirituel et le mystique sont dévalués. La religion n’est plus une institution centrale, elle est remplacée par la raison rationnelle. Il y a un déclin des valeurs : c’est le “désenchantement du monde”.

“L’intellectualisation et la rationalisation croissantes signifient qu’il n’y a en principe aucune puissance imprévisible et mystérieuse qui entre en jouet que l’on peut en revanche maîtriser toute chose par le calcul.”


Le réenchantement par l’art :

\begin{refe}
	Jankélévitch, \textit{La musique et l’ineffable}, 1961
\end{refe}

L’art permettrait de s’émerveiller devant le monde qui nous entoure. La poésie nous fait notamment percevoir le langage.

La musique réussit à exprimer ce que les mots n’arrivent pas à transmettre, ce qui est ineffable. L’ineffable est exprimable au travers de la musique :
“Où manque la parole commence la musique”

La musique “agit comme un enchantement” et elle offre une “parenthèse de rêverie”. La musique nous ouvre sur d’autres mondes.

\subsection{La critique de la technique par l’art}

\begin{refe}
	Hayao Miyazaki, \textit{Le vent se lève}, 2014
\end{refe}

Jiro → Concepteur du Chasseur Zéro
Caproni → Ingénieur italien ami de Jiro

Ce film dénonce l’ambiguïté du progrès technique.

“Tout a été mis en pièces” → Même s’il y a de la beauté dans l’avion, il est voué à être détruit.


\chapter{Pourquoi voulons-nous tuer le temps ?}

Notions impliquées : Temps et Bonheur

Introduction :

« Tuer le temps » est une expression métaphorique : il ne s’agit pas réellement de tuer le temps à proprement parler, comme si le temps était une personne. On retrouve par ailleurs cette expression dans d’autres langues (= ce n’est pas une spécificité de la langue française). Il y a de nombreuses autres expressions de ce genre avec le mot « temps » (prendre son temps, gagner du temps, passer son temps à faire quelque chose...). Cette expression veut dire : « faire passer le temps ». C’est aussi une manière de vivre, une manière de se rapporter au temps, et donc de se rapporter à la vie. → Implique une façon de se représenter le temps, et de l’éprouver.

Pourquoi le voulons-nous ? Pour continuer la métaphore, il s’agira de voir quel est le mobile du crime, quelles en sont ses justifications. D’où vient ce ressentiment contre le temps ? Est-il justifié ?

Il s’agira également de nous demander si nous le voulons réellement, et ainsi d’interroger la notion de volonté. Y a-t-il des conditions particulières qui font qu’on voudrait tuer le temps, ou bien s’agit-il d’une disposition fondamentale de l’humanité dans son rapport au temps ? Que voulons-nous dire et que voulons-nous faire quand nous disons que nous voulons « tuer le temps » ?

→ Il s’agira d’éclairer le rapport spécifique de l’être humain au temps.

Le temps :

La notion de « temps » en philosophie a une dimension existentielle, s’interrogeant – entre autres – sur des sujets comme la mort, l’ennui, l’angoisse, l’existence, l’éphémère, le remords, etc. Mais on peut également étudier le temps d’un point de vue épistémologique (exemple : Le temps est-il un objet de la science ? Est-il un objet de la conscience ?). La notion de temps comprend de nombreux thèmes : la relation temps/espace, la perception du
temps, le temps historique, l’action... Pour ce sujet, c’est la dimension existentielle qui est surtout impliquée, et le rapport que nous entretenons au temps.

Le temps est une notion très complexe, voire insaisissable. C’est un terme très difficile à définir (encore plus que les autres notions du programme...). Comme nous allons le voir dans ce sujet, la perception du temps est quelque chose de très subtil.

Le mot « temps » a beaucoup de sens différents en français, variété qui vient notamment de son étymologie. En effet, « temps » vient du latin « tempestas », qui désigne à la fois : un moment, une époque ; la température ; le mauvais temps, la tempête, l’orage ; le trouble, le malheur, la calamité... Mais il y a aussi en latin le mot « tempestus », qui signifie le moment opportun pour la prise des augures, et par extension le moment opportun, le bon moment pour agir (le « kairos » en grec).

Le mot « temps » a par conséquent de multiples significations en français (le temps météorologique, le temps qui passe, le temps en conjugaison...), et les expressions multiples de la langue française avec le mot « temps » corroborent cette diversité de sens. Dans le cadre d’une interrogation philosophique, nous nous intéresserons au temps uniquement dans le sens du temps qui s’écoule.

Le temps est une tension entre continuité et succession : c’est une succession d’instants qui sont différents mais qui en même temps forment un tout et ne peuvent ainsi être complètement hétérogènes, le temps étant un écoulement continu et ininterrompu. Le temps dans un sens plus particulier peut être considéré comme une mesure de cet écoulement, mesure qui s’accompagne généralement d’une unité de mesure et d’une donnée chiffrée. Mais nous remettrons en question dans le sujet cette définition du temps comme mesurable objectivement. En effet, la volonté de tuer le temps découle d’une certaine appréciation du temps, d’un rapport que nous entretenons avec le temps, qui lui est subjectif.

Le temps peut aussi se référer à un instant particulier, un instant précis, repérable dans une succession chronologique fixée par rapport à un avant (le passé) et un après (le futur). Cet instant est alors le moment présent. Le temps présent se différencie alors de ce qui est antérieur à lui – le passé, conçu par la faculté de la mémoire – et ce qui est postérieur à lui – le futur, conçu par la faculté de l’imagination ou de l’anticipation.

Problématisation du sujet :

A priori, notre préoccupation première est de passer le temps. Nous passons le temps grâce à des occupations, destinées à remplir notre emploi du temps. Nous préférons faire quelque chose plutôt que rien.

Quand voulons-nous particulièrement tuer le temps ? Lorsque nous avons l’impression qu’il ne s’écoule pas. Par exemple : quand nous attendons fortement un événement prochain ou quand nous ressentons de l’ennui dans une situation, ou encore quand nous n’avons rien à faire. → Nous voulons tuer le temps quand, dans l’attente ou dans l’ennui, nous avons l’impression qu’il ne passe pas (en tout cas pas comme nous le voudrions), que le temps est figé (mais ce n’est qu’une impression car il s’écoule bien objectivement). Nous cherchons alors à nous occuper pour que le temps passe plus rapidement, et pour ne pas y penser (avec des « passe-temps »). Ainsi, nous voudrions tuer le temps pour être heureux, ou plutôt pour ne pas être malheureux en ressentant de l’ennui.

Au contraire, lorsque nous sommes occupés, que nous ne nous ennuyons pas, que nous sommes pris par une tâche passionnante, le temps passe très rapidement, ou plutôt : nous ne le sentons pas passer. → Quand le temps passe, il passe inaperçu.

→ Cela implique que nous ne ressentons le passage du temps que lorsque le temps nous est long, et lorsque nous cherchons à le tuer. C’est ainsi le passage du temps lui-même que nous voudrions ne pas ressentir.

Pourquoi ne voulons-nous pas ressentir le passage du temps ? Car cela nous rapproche de la mort, nous fait réaliser que nous ne sommes pas là éternellement et que chaque minute qui passe nous rapproche petit à petit de la mort. Cette pensée sur la mort peut être angoissante (et ainsi source de malheur).

Paradoxe : Si nous voulons tuer le temps pour ne pas ressentir le temps qui passe et ainsi ne pas sentir l’approche de la mort, tuer le temps revient aussi à vouloir que le temps passe plus rapidement, et donc cela reviendrait à une volonté de se rapprocher plus rapidement de la mort. → Vouloir tuer le temps relèverait donc de deux volontés contradictoire (vis-à-vis de la mort). Ne s’agit-il pas en dernière instance de la même volonté ?

Problématique : Ou bien nous voulons tuer le temps pour ne pas ressentir le passage du temps et alors pour ne pas sentir l’approche de la mort, ou bien nous voulons tuer le temps pour accélérer son passage et alors se rapprocher de la mort.

→ Quel rapport à la mort et à l’existence la volonté de tuer le temps révèle-t-elle ?

\section{Nous voulons tuer le temps pour ne pas le sentir passer et ainsi ne pas sentir l’approche de la mort}

\subsection{La dimension insaisissable du temps}

Le temps apparaît tout d’abord comme une chose insaisissable, très difficile à définir. En découle une complexité dans notre rapport au temps et dans la manière dont nous l’appréhendons. Le fait que tout discours passe par des métaphores et des expressions figurées montre bien ce côté insaisissable du temps : nous ne pouvons parler du temps qu’à travers des images. Quelle sorte d’existence est celle du temps ?

\begin{refe}
	Saint Augustin, \textit{Confessions}, Livre XI, Chap. 14
\end{refe}

\begin{bio}[Saint Augustin]
	Aussi appelé Augustin d’Hippone, ou simplement Augustin (comme Thomas d’Aquin, qui peut être appelé Thomas, Saint Thomas, etc. : ne pas confondre les deux philosophes !) est un philosophe et théologien chrétien romain médiéval (fin du IVème siècle / début du Vème siècle) et né en Afrique (en Numidie, dans l’actuelle Algérie). Il fût évêque à Hippone (en Numidie).
\end{bio}

Au chapitre 14 du Livre XI de ses Confessions, il s’interroge et s’étonne sur la notion du temps, qu’il présente comme insaisissable. Il se pose la question : « Qu’est-ce que le temps ? » et remarque que la réponse est loin d’être facile : « Si personne ne m’interroge, je le sais ; si je veux répondre à cette demande, je l’ignore ». Nous savons tous intuitivement ce qu’est le temps, mais le définir parait impossible. Pourtant, il constate que le
temps existe bel et bien : « Et pourtant j’affirme hardiment, que si rien ne passait, il n’y aurait point de temps passé ; que si rien n’advenait, il n’y aurait point de temps à venir, et que si rien n’était, il n’y aurait point de temps présent ». Le temps est nécessairement.

Cette difficulté à définir ce qu’est le temps vient du fait que ce qui constitue la dimension du temps, le passé et le futur, ne sont pas réellement. En effet, le passé est littéralement passé et le futur n’est pas encore : « ces deux temps, le passé et l’avenir, comment sont-ils, puisque le passé n’est plus, et que l’avenir n’est pas encore ? ». Quant au présent, il n’est saisi que par le passé et le futur, il n’a pas d’existence seul, sinon « il serait l’éternité ».

\begin{dfn}[Éternité]
L’éternité peut être soit ce qui n’a ni commencement ni fin, soit ce qui a un commencement mais pas de fin. Nous pouvons dire simplement que c’est un présent éternel, et ainsi une dimension où le temps n’existe plus (dimension atemporelle).
\end{dfn}

Le paradoxe du temps est alors qu’il est parce qu’il disparaît, il est « à condition de n’être plus ». Augustin en conclut alors que le temps est « parce qu’il tend à n’être pas ». Pour résoudre ce paradoxe, Augustin explique que le temps a bien une existence, mais dans l’esprit humain, déclinant ainsi trois modes de présence du temps dans l’esprit : « le présent du passé, le présent du présent et le présent de l’avenir ». Le présent du passé = la mémoire. Le présent du présent = « l’attention actuelle » (on pourrait dire maintenant la conscience). Le présent du futur = l’attente.

→ La réalité du temps est en elle-même insaisissable. Le temps n’existe ainsi que dans les opérations de l’âme humaine. Il n’y a qu’elle qui puisse donner une réalité au temps.

\subsection{Sentir le passage du temps est douloureux}

Le temps a donc bel et bien une réalité, une réalité dans l’esprit humain. Or, si nous voulons tuer le temps, c’est qu’il agit aussi sur nous, et nous sentons les effets de cette action sur nous. Les effets de cette action du temps sur nous sont les effets observables du vieillissement (physique et intellectuel). Cette action du temps est ainsi une action de dépérissement, voire de destruction. Le temps est comme une sorte de puissance enveloppante qui nous entoure, nous domine et nous dépossède petit à petit de nous-mêmes.

Cette représentation du temps est ancrée dans l’imaginaire collectif, notamment à travers l’image de Chronos, le Dieu grec du Temps et qui dévore ses enfants. → Le temps engloutit toute chose sans rien épargner. Nous subissons le temps et le passage du temps.

Le temps paraît agir sur nous par une opération qui nous échappe et qui nous surprend chaque fois que nous en prenons conscience. La conscience du passage du temps passe toujours par une prise de conscience, soudaine et momentanée : brusquement on s’aperçoit que le temps a passé et que ce passage s’est fait à notre insu.

→ Prendre conscience que le temps passe, c’est prendre conscience que le temps nous échappe. Le temps est une puissance invisible à laquelle nous ne pouvons échapper mais qui elle nous échappe à chaque instant.

Nous sommes passifs et impuissants face au passage du temps. En effet, le temps est irréversible, et ceci en deux sens : 1. Il est impossible de revivre ce qui a été vécu. 2. Il est impossible de défaire ce qui a été fait (« What is done cannot be undone », Macbeth, Shakespeare). On ne peut arrêter la fuite du présent. Chaque instant, bien ou mal vécu, disparaît, mais continue d’avoir été bien ou mal vécu pour toujours.

\begin{refe}
	Husserl, \textit{Leçons pour une phénoménologie de la conscience intime du temps}
\end{refe}
« Le temps est rigide et pourtant le temps coule. Dans le flux du temps, dans la descente continue dans le passé, se constitue un temps qui ne coule pas, absolument fixe, identique, objectif. Tel est le problème ».

On retrouve beaucoup de traces de cette puissance irréductible du temps chez Baudelaire. Par exemple dans Les Fleurs du mal, le temps est appelé « l’Ennemi » dans le poème du même nom : « Ô douleur ! Ô douleur ! Le temps mange la vie, / Et l'obscur ennemi qui nous ronge le cœur / Du sang que nous perdons croît et se fortifie ! », ou encore dans Petits poèmes en prose (un poème d’ouverture) : « Le Temps règne en souverain maintenant ; et avec le hideux vieillard est revenu tout son démoniaque cortège de Souvenirs, de Regrets, de Spasmes, de Peurs, d’Angoisses, de Cauchemars, de Colères et de Névroses. [...] Oui ! le Temps règne ; il a repris sa brutale dictature ».

Les moments de prise de conscience du passage du temps sont douloureux car nous nous rendons compte que ne pouvons échapper au temps, nous sommes prisonniers de la dimension temporelle de l’existence. Cela nous rappelle que nous sommes des êtres finis, qui sommes condamnés au même sort : la mort. Nous nous sentons impuissants face à la puissance irréductible qu’est le temps. Et cela procure de l’angoisse.

→ Face à cette angoisse, nous faisons alors tout pour ne pas sentir le passage du temps et nous désirons alors tuer le temps.

Nous sentons principalement le temps lorsque nous le trouvons long, quand il ne passe pas, et donc quand nous nous ennuyons.

\begin{dfn}[Ennui]
Sentiment désagréable mêlé de lassitude, fatigue, de découragement, de vide, d'inutilité, causé par l'inaction et le manque d'activité ou le manque d'intérêt dans une action.
\end{dfn}

Dans l'ennui, je sens de façon très profonde le temps qui passe, même si paradoxalement j’ai l’impression que celui-ci ne passe pas. En effet, dans l'ennui, je sens particulièrement mon existence en tant qu’elle est prisonnière du temps qui s’écoule. Cette sensation advient dans l’ennui car justement lorsque je m’ennuie mes pensées ne sont occupées par rien d’autre et sont donc tout à fait libres de penser au temps qui passe. L’expérience de l’ennui est alors un grossissement de l’angoisse, un temps dans lequel je sens plus que n’importe quel autre moment l’écoulement du temps puisque j’ai justement tout le temps d’y songer.

\subsection{Comment tuer le temps ?}

Les êtres humains ont alors mis en place des stratégies pour ne pas penser au temps, pour ne pas sentir le passage du temps et ainsi être délivrés de l’angoisse. Il s’agit de lutter efficacement contre l’ennui.

En effet, on peut considérer que toutes nos actions tournent autour de la tentative d’échapper à l’ennui, l’ennemi principal de l’existence (Cf. Txt Kant, Réflexions sur l’éducation, dans le Corpus de Textes sur le Travail : l’homme a « besoin d’occupations » pour qu’il ne « se sente plus lui-même » : « Il est tout aussi faux de s’imaginer que si Adam et Ève étaient demeurés au Paradis, ils n’auraient rien fait d’autre que d’être assis ensemble, chanter des chants pastoraux, et contempler la beauté de la nature. L’ennui les eût torturés tout aussi bien que d’autres hommes dans une situation semblable »).

Dans le tome I de L’Alternative, dans « La Culture alternée », Kierkegaard peint l’habitude et l’ennui comme les principes régisseurs du monde. L’ennui y apparaît comme « la source de tout mal » et il est ainsi « naturel » que « cherche à le surmonter ». L’homme, depuis sa création, s’ennuie et le processus même de la création a pour but d’éviter l’ennui, ce qu’explique le personnage de L’Alternative (appelé A) dans cette sorte de « genèse de l’ennui », ou plutôt de « genèse de l’homme » qui se construirait autour de l’ennui : « Les dieux s’ennuyaient : ils créèrent les hommes. Adam s’ennuyait seul : Ève fut créée. Dès ce moment, l’ennui fut dans le monde et il s’étendit dans l’exacte mesure où la population se propagea. Adam s’ennuya d’abord seul, puis en compagnie d’Ève ; puis, Adam, Ève, Caïn et Abel s’ennuyèrent en famille ; puis, la population croissant dans le monde, les peuples s’ennuyèrent en masse ».

Le divertissement en est le principal moyen d’échapper à l’ennui. En effet, se divertir permet de ne pas songer au temps qui passe, de se détourner de cette pensée en occupant notre esprit autrement.

\begin{dfn}[Divertissement]
Moyen trouvé par l’homme pour supporter l’angoisse. C’est le fait de se livrer à des activités multiples pour oublier quelque chose de désagréable (ici le passage du temps), action de se distraire, de se changer les idées par l’amusement.
\end{dfn}

\begin{refe}
	Pascal, \textit{Pensées}, L136
\end{refe}

Le divertissement tel qu’il est théorisé par Pascal consiste à détourner son esprit par des occupations diverses qui remplissent l’esprit. Pascal donne beaucoup d’exemples de divertissement de son temps : le jeu, la chasse, la danse, le théâtre, la poésie, la guerre, l’apprentissage, la conversation, la passion en général... (Les divertissements aujourd'hui se sont considérablement multipliés). Le fait de me divertir me permet de me détourner des pensées angoissantes qui m’assaillent lorsque je n’ai aucune occupation, donc quand je m’ennuie. En effet, le sens premier du divertissement est celui du détournement (« divertissement » vient du latin « divertere » qui signifie « se détourner »). Se divertir est un moyen de tuer le temps puisqu’en s’amusant mon esprit est occupé et ne songe pas au temps qui passe. Ainsi le divertissement m’empêche de me préoccuper de la dimension temporelle de l’existence.

Pascal prend des exemples de toutes les conditions sociales pour montrer qu’aucun homme n’échappe au divertissement (bien qu’il faille avoir les moyens financiers de se divertir... Mais telle n’est pas la préoccupation de Pascal) : même les rois sont touchés par ce phénomène.

Ce que recherchent les hommes dans le divertissement c’est l’agitation, le bruit et le remuement, qui s’opposent à la solitude et au silence qui leur feraient penser à leur condition d’êtres finis : « Aussi les hommes qui sentent naturellement leur condition n’évitent rien tant que le repos ; il n’y a rien qu’ils ne fassent pour chercher le trouble ». Plus leur occupation est « violente et impétueuse », moins ils pensent à leur condition. Mais c’est paradoxalement le repos que les hommes cherchent dans l’agitation du divertissement : « Ainsi s’écoule toute la vie ; on cherche le repos en combattant quelques obstacles et si on les a surmontés le repos devient insupportable par l’ennui qu’il engendre. Il en faut sortir et mendier le tumulte ».

Grâce à son imagination, l’homme s’illusionne lui-même dans le processus du divertissement. L’homme sait que ce n’est pas le gain au jeu ou bien la prise à la chasse qu’il poursuit. Cependant sans se fixer ce but il ne voudrait pas jouer ou chasser, il n’y verrait pas d’intérêt – bien que l’intérêt ne soit pas dans le gain ou la prise mais
dans l’agitation du jeu ou de la chasse. Il faut la poursuite d’un but à l’homme pour qu’il y ait divertissement. Il faut donc que l’homme s’imagine poursuivre le gain ou la prise pour que le divertissement soit effectif : « Il faut qu’il s’y échauffe, et qu’il se pipe [= se trompe] lui-même en s’imaginant qu’il serait heureux de gagner ce qu’il ne voudrait pas qu’on lui donnât à condition de ne point jouer ».

Cependant il faut souligner que Pascal ne prône ni ne blâme le divertissement, il ne fait qu’un constat de son importance au sein de l’existence humaine (même si on peut sentir un certain ton de blâme par moments, et c’est le cas dans d’autres fragments).


Kierkegaard présente lui aussi le divertissement comme un moyen de combattre l’ennemi principal qu’est l’ennui. De la même manière, le processus du divertissement est paradoxal dans le sens où l’homme cherche le repos par le trouble : « L’homme cherche le repos, mais que de changements : le jour et la nuit, l’été et l’hiver, la vie et la mort ; il cherche le repos, mais que de vicissitudes : bonheur et malheur, joie et tristesse se succèdent ; il cherche le repos et la constance, mais quelle variation : à l’ardeur du projet se substitue le dégoût de l’impuissance, aux riantes perspectives de l’attente la gloire désabusée du succès ; il cherche le repos, et jusqu’où ne le cherche-t-il pas : jusque dans les distractions inquiètes ; où ne le cherche-t-il pas en vain, même dans la tombe ! » (Discours édifiant, « L’écharde dans la chair »). L’homme cherche paradoxalement la constance par l’inconstance et le repos dans l’agitation de la distraction.

Baudelaire dans Petits poèmes en prose substitue le divertissement à l’ivresse : « Il faut être toujours ivre. Tout est là : c’est l’unique question. Pour ne pas sentir l’horrible fardeau du Temps qui brise vos épaules et vous penche vers la terre, il faut vous enivrer sans trêve. Mais de quoi ? De vin, de poésie, ou de vertu, à votre guise. Mais enivrez-vous ». Le travail, la lecture, le visionnage d’un film... Tout cela peut être considéré comme un divertissement. En effet, la fiction nous permet d’oublier sa propre dimension temporelle en se plongeant dans une autre. La narration élimine le temps perdu, on entre dans le temps des personnages de l’histoire. Dans une fiction, les personnages perdent rarement leur temps, on les voit toujours vaquer à des occupations (sinon le spectateur s’ennuie). La narration présente des personnages pour qui toute heure compte, qui n’ont pas à « tuer le temps ».

→ Le divertissement est un moyen de faire passer le temps en nous occupant. Ainsi, on ne peut oublier le temps qui passe qu’en faisant quelque chose de son temps.

Transition :

Nous avons dit que ce qui était insupportable dans l’ennui était le fait que quand nous nous ennuyons, nous avons tout le temps de penser au temps, à la dimension temporelle et à la mort qui approche. Pour combattre l’ennui, nous voulons faire passer le temps plus rapidement, pour ne plus sentir le temps qui passe. Or, faire passer le temps plus rapidement, c’est littéralement s’approcher de la mort. Ainsi à travers le divertissement on voudrait paradoxalement éloigner la pensée de la mort en se rapprochant de la mort : la volonté de tuer le temps ne révèle- t-elle pas alors une volonté de mort ?

\section{Nous voulons tuer le temps pour le faire passer plus rapidement et ainsi se rapprocher de la mort}

\subsection{Comment le temps peut-il être accéléré ?}

Vouloir tuer le temps revient également à vouloir se rapprocher de la mort. En effet vouloir qu’un laps de temps passe plus rapidement, c’est vouloir le rendre le plus court possible, et si l’on applique cette logique à tous les laps de temps, à toute la durée de l’existence, cela revient à vouloir faire de son existence l’existence la plus courte possible. Or, réduire son existence au maximum, c’est mourir. Le désir de tuer le temps reviendrait alors à un désir de mort, un désir d’écourter sa propre existence, la mort étant un état de repos absolu. La mort serait-elle le but final du divertissement ?

Mais comment est-il même possible de faire passer le temps plus rapidement ? Cela est-il possible ? Nous ne pouvons concrètement accélérer le temps.

Il faut alors distinguer deux types de temps : le temps objectif et le temps subjectif.

\begin{refe}
	Bergson, \textit{La Pensée et le Mouvant} (Partie I : Introduction).
\end{refe}

Dans cet essai, et dans de nombreux autres livres, Bergson distingue le temps objectif du temps subjectif (c’est un des grands apports de sa philosophie).


Le temps objectif est le « temps des mathématiciens », le temps scientifique, le temps tel qu’il est mesurable objectivement, le « temps des horloges ». C’est la mesure quantitative du temps.


Le temps subjectif quant à lui est le temps tel qu’il est perçu et vécu par chaque individu : cette perception est singulière et subjective. Ce temps est appelé la durée. Notre ressenti du temps est rarement le même que le temps objectif. Ainsi, pour diverses raisons, une minute peut nous paraître « une éternité » (hyperbole pour dire que la minute est ressentie comme plusieurs). À l’inverse, une heure peut passer extrêmement vite, sans que je m’en aperçoive, « en un clin d’œil », là aussi pour plusieurs raisons. En ce sens la durée est une appréciation non plus quantitative mais qualitative du temps.

→ Grâce à cette distinction, on comprend que le temps peut nous paraître long. Bien que le temps objectif lui ne change pas, le temps subjectif, lui, varie continuellement.

Mais pour Bergson seule la durée est la véritable conception du temps. En effet le temps scientifique, lui, ne voit le temps que comme une succession d’instants et non pas comme une réalité fluide et indivisible (ce qu’est pourtant réellement le temps). Voir le temps comme une succession d’instants, c’est le spatialiser, c’est analyser le temps par le prisme de l’espace. Or, le temps est avant tout un passage (son essence est « de passer »). Ainsi seule la durée est « le temps réel » (ainsi quand Bergson parle du « temps » il se réfère principalement à la durée).

\begin{dfn}[Espace]
Milieu homogène dans lequel nous situons tous les objets et leurs déplacements. Domaine dans lequel nous-mêmes nous mouvons, nous vivons, nous agissons.
\end{dfn}

Le temps mathématique peut être représenté par une ligne (allant du présent au futur) avec une multitude de segments qui sont les instants ou les moments, divisibles en secondes, minutes, heures, jours, etc. Représenter le temps c’est bien le spatialiser. Or la durée ne peut être rendue immobile par une représentation spatiale : « La ligne qu’on mesure est immobile, le temps est mobilité. La ligne est du tout fait, le temps est ce qui se fait ». La durée est une dimension propre de la conscience, elle se ressent et il est difficile d’en avoir une conception claire et objective : « on la sent et on la vit ».

\subsection{La mort est l’horizon indépassable de nos pensées}

Toutes nos préoccupations sont faites en fonction de la mort. Nous songeons continuellement et inévitablement à la mort et nous avons beau avoir tous les divertissements possibles, nous ne pouvons échapper à la pensée de la mort.

\begin{refe}
	Épicure, \textit{Lettre à Ménécée}
\end{refe}

\begin{bio}[Épicure]
	Épicure est un philosophe de l’Antiquité grecque (fin du IVe siècle / début du IIIe siècle av. J.-C.). Il est le fondateur du courant antique appelé l’épicurisme. L’épicurisme fait partie, comme de nombreuses autres philosophies antiques (à l’instar du stoïcisme), des philosophies dites « du bonheur », des philosophies pratiques pour enseigner à l’homme comment être heureux. Cependant l’épicurisme repose tout de même sur une théorie (notamment physique : Épicure soutient notamment que la nature est composée d’un ensemble d’atomes).
\end{bio}

Le grand principe de l’épicurisme consiste à associer le bonheur au plaisir. Le plaisir est ainsi le plus grand des biens. À l’inverse, le malheur est associé à la souffrance. La souffrance (physique ou mentale) doit être éviter pour notre bien-être : « Le plaisir est reconnu par nous comme le bien primitif et conforme à notre nature, et c’est de lui que nous partons pour déterminer ce qu’il faut choisir et ce qu’il faut éviter ».

\begin{dfn}[Plaisir]
État affectif agréable que procure la satisfaction d'un besoin ou d’un désir. C’est un sentiment fugace de bien-être. Le plaisir se dit plutôt pour une sensation physique agréable, mais on peut tout aussi parler de plaisir intellectuel.
\end{dfn}

Cependant il ne faut surtout pas voir Épicure et les épicuriens comme des personnes qui recherchent à tout prix le plaisir dans l’excès. Les Épicuriens ont dû en effet faire face à de vives critiques, venant du fait que leur doctrine était mal comprise. En effet, l’épicurisme repose avant tout sur la sagesse et la sobriété : les plaisirs doivent être modérés, et tous les désirs ne sont pas à conserver. Épicure met en garde contre ces mauvaises interprétations de sa philosophie : « Quand donc nous disons que le plaisir est le but de la vie, nous ne parlons pas des plaisirs des voluptueux inquiets, ni de ceux qui consistent dans les jouissances déréglées, ainsi que l’écrivent des gens qui ignorent notre doctrine, ou qui la combattent et la prennent dans un mauvais sens. Le plaisir dont nous parlons est celui qui consiste, pour le corps, à ne pas souffrir [aponie] et, pour l’âme, à être sans trouble [ataraxie]. Car ce n’est pas une suite ininterrompue de jours passés à boire et à manger, ce n’est pas la jouissance des jeunes garçons et des femmes, ce n’est pas la saveur des poissons et des autres mets que porte une table somptueuse, ce n’est pas tout cela qui engendre la vie heureuse, mais c’est le raisonnement vigilant, capable de trouver en toute circonstance les motifs de ce qu’il faut choisir et de ce qu’il faut éviter, et de rejeter les vaines opinions d’où provient le plus grand trouble des âmes ».

Il faut donc distinguer l’épicurisme de l’hédonisme. En effet, l’hédoniste recherche le plaisir pour lui-même, en fait le but de la vie, alors que le but de l’épicurisme est d’être heureux. Tous les plaisirs ne sont pas bons, et par conséquent tous les désirs ne sont pas à satisfaire, et il faut parfois passer par des douleurs pour accéder à un plaisir plus grand : « Mais, précisément parce que le plaisir est le bien primitif et conforme à notre nature, nous ne recherchons pas tout plaisir, et il y a des cas où nous passons par-dessus beaucoup de plaisirs, savoir lorsqu’ils doivent avoir pour suite des peines qui les surpassent ; et, d’autre part, il y a des douleurs que nous estimons valoir mieux que des plaisirs, savoir lorsque, après avoir longtemps supporté les douleurs, il doit résulter de là pour nous un plaisir qui les surpasse ».

Ainsi, Épicure distingue :
les désirs naturels et nécessaires (pour vivre).
les désirs naturels et non nécessaires (variation des plaisirs et recherche de l’agréable).
les désirs non naturels et non nécessaires, appelés aussi désirs vains (les désirs artificiels, par exemple être millionnaire ; et les désirs irréalisables, par exemple être immortel).

Nous ne devons suivre que les désirs naturels, qu’ils soient nécessaires ou non nécessaires, et ce avec modération.

Épicure part du constat que nous vivons sans cesse dans la crainte de la mort et des dieux. Nous avons peur de la mort, elle est notre principale préoccupation et ce mal est « celui des maux qui nous donne le plus d’horreur ». Deux attitudes sont constatées par Épicure : « La multitude tantôt fuit la mort comme le pire des maux, tantôt l’appelle comme le terme des maux de la vie ». Or, que ce soit en voulant « fuir » la mort à tout prix, ou en voulant « l’appeler » on fait de la mort une réalité proche de nous puisqu’elle ne quitte plus nos pensées.

Or, pour Épicure, il n’y a pas de raison d’avoir peur de la mort. En effet, dans l’épicurisme, le malheur est associé à un mal-être physique. Or, la mort ne peut être douloureuse puisque quand nous mourrons nous n’avons plus aucune sensation. Le raisonnement est le suivant : « Prends l’habitude de penser que la mort n’est rien pour nous. Car tout bien et tout mal résident dans la sensation : or la mort est privation de toute sensibilité. Par conséquent, la connaissance de cette vérité que la mort n’est rien pour nous, nous rend capables de jouir de cette vie mortelle ». Et Epicure de conclure : « Ainsi celui de tous les maux qui nous donne le plus d’horreur, la mort, n’est rien pour nous, puisque, tant que nous existons nous-mêmes, la mort n’est pas, et que, quand la mort existe, nous ne sommes plus. Donc la mort n’existe ni pour les vivants ni pour les morts, puisqu’elle n’a rien à faire avec les premiers, et que les seconds ne sont plus ».

Bien sûr, cette philosophie (comme les philosophies pratiques), est difficile à mettre en place ! On peut être conscient de cela sans pour autant pouvoir l’appliquer (ce qui apparaît encore plus dans le stoïcisme).

\subsection{Vouloir supprimer la dimension temporelle de l’existence, c’est vouloir se rapprocher de la mort}

Dans son extrême conséquence, vouloir tuer le temps ne se réduirait-il pas à vouloir mourir ? Vouloir occulter la dimension temporelle de l’existence, c’est sortir de l’existence, c'est-à-dire mourir. En effet vouloir qu’un laps de temps passe plus rapidement, c’est vouloir le rendre le plus court possible, et si l’on applique cette logique à tous les laps de temps, à toute la durée de l’existence, cela revient à vouloir faire de son existence l’existence la plus courte possible. Or, réduire son existence au maximum, c’est mourir : vouloir que son existence soit la plus courte possible, c’est vouloir se rapprocher le plus possible de la mort, c’est vouloir mourir. La conséquence extrême de ne pas vouloir sentir le passage du temps est de vouloir mourir : dans la mort, il est évident que nous ne nous ennuyons plus et que nous ne sentons plus notre existence. Ainsi un désir d’éternité reviendrait à un désir de mort. De la même manière dans la religion chrétienne, l’éternité peut espérer s’atteindre dans « l’Au-delà », qui ne s’acquiert qu’une fois que l’on meurt, qu’une fois que cette vie-ci, si tant est qu’il y en ait une autre, se termine. Vouloir tuer le temps revient ainsi à vouloir mourir.

\begin{refe}
	Nietzsche, \textit{Généalogie de la morale}, Traité II (« La « faute », la « mauvaise conscience » et ce qui leur ressemble »)
\end{refe}

Nietzsche, dans le second traité de la Généalogie de la morale, analyse l’origine de ce ressentiment contre le temps en l’appliquant à la morale chrétienne. Si le chrétien croit en la vie éternelle, croit en l’Au-delà, accessible après la mort, cela vient d’une haine du temps présent, d’un ressentiment contre ce temps-ci (qui est en réalité l’unique temps pour Nietzsche) et il justifie cette haine en dévaluant ce temps, impur par rapport au temps éternel auquel le croyant accède au moment de sa mort. C’est ce que Nietzsche appelle la « morale du ressentiment ». Cette morale se matérialise par de nombreuses restrictions, elle prône des valeurs telles que la pauvreté, la faiblesse, le renoncement, le silence, la patience, la solitude... Ce qui mène à ce que Nietzsche appelle « la vie ascétique » ou « l’idéal ascétique », véhiculée par la figure du prêtre ascétique, appelé l’ « ennemi de la vie » ou encore le « négateur de la vie ».

La religion chrétienne méprise la vie terrestre et aspire à trouver son bonheur uniquement dans l’Au-delà, autrement dit dans la mort. Or la vie terrestre est la seule qui existe pour Nietzsche : mépriser la vie terrestre revient à mépriser tout simplement la vie. Ainsi la morale chrétienne s’oppose à la vie elle-même. Dans un fragment posthume (Fragments 11 [159] (1881)), Nietzsche écrit que la religion « méprise cette vie-ci en tant que
fugitive » et veut nous apprendre à « élever nos regards vers une incertaine, autre vie », c'est-à-dire la vie éternelle (qui pour Nietzsche n’existe pas).

Transition :
Nous avons dit que la volonté de tuer le temps relevait cette fois-ci d’une volonté d’écourter notre existence et donc de se rapprocher de la mort. Apparaît alors un paradoxe : vouloir tuer le temps serait à la fois vouloir occulter la dimension de la mort, et à la fois vouloir s’en rapprocher. Mais n’est-ce pas plus fondamentalement un désir d’échapper à la dimension de l’existence elle-même ? En effet, vouloir occulter la mort revient à vouloir supprimer la pensée de notre condition humaine, qui se termine inévitablement par la mort. Et vouloir se rapprocher de la mort, c’est vouloir également échapper à notre condition humaine. Le paradoxe se résout alors si on comprend que vouloir tuer le temps revient en dernière instance à vouloir échapper à notre condition humaine (finie et absurde).

\section{Mais finalement cela revient à la même volonté : la volonté d’échapper à l’angoisse que procure la condition humaine}

\subsection{La volonté de tuer le temps naît face au constat de la réalité de notre condition humaine}

\begin{refe}
	Nietzsche, \textit{Généalogie de la morale}, Traité III (« Que signifient les idéaux ascétiques ? »), §28
\end{refe}

Nietzsche ne s’arrête pas là dans son raisonnement. La haine de la religion chrétienne contre le temps ne peut se réduire à un désir de mort. Dans ce troisième traité, il souligne en effet le côté paradoxal de ce désir. La vie ascétique apparaît comme une « contradiction en soi ».

L’idéal ascétique est en réalité une ruse pour conserver la vie. Le véritable sens de cet idéal est de lutter contre la mort et de se battre pour la vie (c’est une « ruse de la conservation de la vie »). Le prêtre ascétique n’est qu’un « ennemi apparent de la vie », en réalité il fait partie « des très grandes forces conservatrices et affirmatives de la vie ».

Comment ce paradoxe s’explique-t-il ? Nietzsche explique cela à la toute fin de la Généalogie de la morale, par ce qu’il appelle « la véritable signification de l’idéal ascétique ». L’idéal ascétique donne en réalité un sens à la souffrance de l’homme. Pourquoi l’homme souffre-t-il ? Non seulement car la vie est faite de souffrances, de maladies, de difficultés, etc. mais surtout car cette souffrance n’a pas de sens. Ainsi, plus encore que la
souffrance, « l’absence de sens de la souffrance » fait encore plus souffrir. En définitive, la volonté de tuer le temps relève d’une volonté de trouver un sens à l’existence. L’idéal ascétique, dans sa lutte contre le temps, comble l’« immense vide » ressenti face au non-sens de l’existence. Grâce à cet idéal, la vie acquiert un sens, l’homme est « sauvé ». Même si cela entraîne une souffrance plus grande : c’est cela le paradoxe de l’idéal ascétique : la vie ascétique est tout sauf agréable (elle est faite de pauvreté, de mortifications, de sacrifices, de renoncements au plaisir... ; elle repose sur la conscience douloureuse du péché originel, et ainsi sur la culpabilité...), mais au moins elle donne à la vie un sens.

Tout cela doit vous faire penser à la réflexion de Marx sur la religion ! (Cf. Séquence 2). Mais le point de vue est plus existentiel ici.

Plus encore, Nietzsche explique que l’idéal ascétique entraîne une haine de la vie, une haine de l’humain, une peur du bonheur, une horreur des sens, du devenir, du désir... En bref, l’idéal ascétique veut anéantir tout ce qui fait que nous sommes humains : cela revient à une « volonté de néant » pour Nietzsche (ce qui est différent de la volonté de mort), mais tout cela s’explique en dernière instance par le fait que « l’homme aime mieux vouloir le néant que de ne pas vouloir ».

On voyait déjà des traces de cette pensée chez Pascal (mais chez Pascal la critique ne se porte bien sûr pas sur la morale chrétienne). Si l’homme se divertit, c’est pour ne pas songer à la mort, mais c’est avant tout pour ne pas songer à sa « condition misérable ». Chez Pascal, la condition de l’homme est « misérable » principalement car elle est « vaine ». L’homme est faible, condamné à mourir et condamné à la souffrance.

\begin{dfn}[Vanité]
Au sens premier et étymologique du mot, la vanité est le sentiment qui vient de la prise de conscience de l’insignifiance de la vie. C’est le caractère de ce qui est sans consistance, sans valeur et illusoire. C’est également le sens dit « théologique », utilisé le plus souvent par la tradition chrétienne pour montrer que cette vie-ci n’est pas la véritable vie, qu’elle n’a pas de valeur et qu’elle est futile.
\end{dfn}

En effet dans l’ennui l’homme a tout le temps de songer à sa condition, comme le montre Pascal à travers le fragment L622 des Pensées : « Rien n’est si insupportable à l’homme que d’être dans un plein repos, sans passions, sans affaire, sans divertissement, sans application. Il sent alors son néant, son abandon, son insuffisance, sa dépendance, son impuissance, son vide. Incontinent [= aussitôt] il sortira du fond de son âme l’ennui, la noirceur, la tristesse, le chagrin, le dépit, le désespoir ».

Le divertissement est un moyen de dissimuler le vide de l’existence et de le combler et serait ainsi comme une sorte de leurre qui masquerait le vide. En ce sens, le divertissement est bien efficace en ce qu’il empêche les hommes de penser à leur misère. Cependant cette solution n’est pas durable et ne fait que recouvrir le malheur des hommes.

\subsection{Mais cette volonté est vaine car nous ne pourrons jamais échapper à la dimension de l’existence}

Vouloir tuer le temps pour échapper à l’absurdité de l’existence, pour donner un sens à l’existence ou pour ne pas songer à cette absurdité à travers le divertissement est illusoire. En effet, nous ne pourrons jamais échapper à cette condition.
Tuer le temps est ainsi impossible car nous ne pouvons échapper à qui ce constitue  l’existence elle-même, à part en mourant.

\begin{refe}
	Nietzsche, \textit{Ainsi parlait Zarathoustra}, Partie II, « De la rédemption »
\end{refe}

La volonté est impuissante à tuer le temps, elle est « impuissante à l’égard de tout ce qui a été fait ». La volonté ne peut revenir en arrière, ce qui a été fait ne peut être défait (« il n’y a pas d’action qui puisse être anéantie »), la volonté se retrouve prisonnière du passé. C’est « l’affliction la plus solitaire » de la volonté, qui mène au ressentiment : la volonté « en veut au temps », elle « contemple le passé, pleine de colère », impuissante face à la dimension temporelle de l’existence : « Que le temps ne fasse pas marche arrière, voilà ce qui l’irrite ».

Ce que la volonté invente alors pour tuer le temps, c’est la vengeance. Faute de pouvoir être en acte, car on ne peut modifier ce qui est passé, la vengeance sera spirituelle (naît alors « l’esprit de vengeance ») : on dit que si tout passe, c’est que tout mérite de passer, et que seul a de la valeur ce qui ne passe pas. C’est une forme de rédemption (se racheter d’une faute passée), mais une mauvaise rédemption pour Nietzsche (c’est de la « folie »). Cette mauvaise rédemption est notamment celle du Christianisme et du prêtre ascétique. La véritable rédemption pour Nietzsche consiste à libérer la volonté elle-même, qui est « prisonnière » de cet esprit de ressentiment. 

La véritable rédemption dit que la volonté ne peut vouloir en arrière, mais elle peut vouloir à nouveau. Vouloir qu’une chose soit, ce n’est pas vouloir qu’elle reste sans subir de changement, mais vouloir qu’elle devienne ce qu’elle est et revienne ce qu’elle était, et ceci éternellement (« métamorphoser tout « c’était » en un « je le voulais ainsi ! »). Il faut que la « volonté créatrice » dise du passé « Mais je le veux ainsi ! Je le voudrai ainsi ! ». Vouloir, c’est vouloir le passage du temps, et ce vouloir ne requiert rien d’autre qu’un acte d’affirmation. La volonté est son propre libérateur, elle se « réconcilie » ainsi avec le temps, et se délivre du désir de vengeance.

Ainsi, c’est parce que l’homme, à travers la mauvaise rédemption, veut arrêter ou changer ce qui ne peut l’être qu’il subit le passage du temps.

\subsection{Accepter l’existence, c’est accepter sa dimension temporelle}

Ainsi vivre, c’est accepter d’être soumis au temps, c’est accepter que le temps s’écoule dans que nous puissions le retenir. Accepter la vie, c’est accepter qu’il soit impossible de tuer le temps.
Il y aurait tout de même des moyens d’accéder, du moins psychologiquement, à une forme d’éternité, qui ne tuerait pas le temps mais qui le conserverait, et permettrait de lutter contre la vanité du temps qui passe et qui mène à la mort.

\begin{refe}
	Platon, \textit{Le Banquet}, 208a
\end{refe}

Dans \textit{le Banquet}, Platon, à travers le discours de Diotime, distingue deux moyens de se rapprocher de l’éternité, sans pour autant supprimer toute dimension temporelle : par la procréation et par le travail philosophique. Les hommes ont en eux un désir d’immortalité, désir qu’il est physiquement impossible d’assouvir, mais dont on peut se rapprocher.
Le premier moyen, celui de la procréation, consiste à laisser une trace de nous-mêmes sur terre par notre progéniture. Nous continuons ainsi d’exister à travers notre descendance. Le deuxième moyen, celui du travail philosophique, consiste à engendrer une sorte de progéniture livresque, une descendance philosophique. Certains hommes ne veulent pas procréer mais peuvent tout de même laisser une trace de leur passage sur terre en écrivant ou en véhiculant des pensées. C’est le cas du philosophe.
Nous pourrions ajouter à cette idée de Platon le travail de l’artiste, dont l’œuvre engendrée est comme son enfant, et qui lui permet d’accéder à l’immortalité. Ainsi créer une œuvre serait un moyen de lutter contre le temps et d’entrer dans l’éternité, sans pour autant supprimer la dimension temporelle de l’existence. L’éternité acquise par celui qui laisse une descendance n’est pas une véritable éternité mais une éternité d’ordre psychologique, une éternité qui ne supprime pas la dimension temporelle. Plus encore, c’est l’angoisse même du temps qui motive l’entrée dans cette éternité : l’angoisse est donc un moteur utile et vouloir la supprimer, c’est vouloir mourir.

\chapter{La notion du Devoir}

\begin{rep}[Obligation/Contrainte]
Obligation
Contrainte
On peut se soustraire à l’obligation
On ne peut pas se soustraire à la contrainte.
Vient de soi-même, de l’intérieur
Imposée de l’extérieur
Vient de ma raison (Kant) ou de ma conscience (au sens moral)
\end{rep}


Question de morale et de moralité
Pas de question de morale et de moralité
Inclut la liberté
Pas question du tout de liberté
On obéit => Activité
On subit => Passivité
Acte volontaire
Acte involontaire
Ensemble des obligations constitue la loi morale qui vient soit de ma raison (Kant) soit de ma conscience (Rousseau)


Peut être liée à des instances religieuses, familiales, professionnelles, sociétales




\begin{dfn}[Morale]
Ensemble de règles que doit suivre un individu pour faire le bien. Fixe les fins et les moyens de l’action.
\end{dfn}

\begin{refe}
	Kant, \textit{Fondements de la métaphysique des mœurs} 
\end{refe}

La morale relève de la raison pratique (par opposition avec la raison pure). Elle permet de distinguer le bien du mal pour mieux agir.

À partir de Kant, on passe de morale du bonheur (comment agir pour être heureux, stoïcisme, épicurisme) à la morale du devoir (comment bien agir, parce que bien agir c'est bien).

Pour Kant, bien agir c'est agir avec bonne volonté (les intentions pures et la morale) : c'est ce que l'on nomme le kantisme ou morale kantienne.
Le bien se situe donc dans le fondement de l'action (la volonté). Le bien est donc détectable à la volonté.

Il y a chez Kant une distinction entre deux impératifs :
\begin{dfn}[Impératif]
Commandement, règle de conduite, ce qui nous oblige à agir et relève du devoir
"Formules qui déterminent l'action" selon Kant
\end{dfn}

Impératifs hypothétiques 
Impératifs catégoriques
Peut être autrement
Ne peut être autrement
Hypothèse 
Affirmation, certitude 
Conseil
Ordre
Contestable
Incontestable
Nuancé
Radical
Hésitation
Volonté ferme, détermination 
Moyen pour arriver à une fin :
"Si l'action n'est bonne que comme un moyen pour quelque chose d'autre"
->Fausse morale
-> Conditionnel
L'action est à la fois un moyen et une fin :
"Si l'action est représentée comme bonne en … et comme devant être nécessairement le principe d'une volonté conforme à la raison"
-> Inconditionnel

Pour Kant seul l'impératif catégorique dépend de la morale.

"Je ne sais pas d'avance ce qu'il contiendra jusqu'à ce que la condition me soit donnée" : l'impératif hypothétique est conditionnel.

"Je sais aussitôt ce qu'il contient" : l'impératif catégorique est inconditionnel.

L'impératif hypothétique dépend des sentiments, émotions (variables) alors que l'impératif catégorique ne dépend que de la morale (immuable).

Impératif hypothétique : conseil, intéressé
Impératif catégorique : ordre, désintéressé

Kant rejette fermement tout mensonge, qu'il classe comme impératif hypothétique.

\begin{rep}[Formel/Matériel]
L’impératif catégorique “ne concerne pas la matière l’action et ce qui doit en résulter mais la forme”, alors que l’impératif hypothétique, lui, ne concerne que le contenu, la matière de l’action puisqu’il adapte son contenu à la fin que je vise, qui dépend de mon propre intérêt. L’impératif catégorique est dit formel + les deux définitions ci-dessous
\end{rep}
\begin{dfn}[Formel]
Ce qui détermine la nature et les propriétés d’un objet
\end{dfn}
\begin{dfn}[Matériel]
Ce qui relève de l’élément dont un chose est faite, son contenu, son sens
\end{dfn}


Ainsi, pour Kant, la seule action bonne est bonne dans sa forme et non dans son contenu.
Tous les impératifs catégoriques peuvent être résumés en un seul impératif catégorique :

Impératif de l'universalité de la loi morale :
«Agis uniquement d'après la maxime qui fait que tu peux vouloir en même temps qu'elle devienne une loi universelle.»

\begin{rep}[Universel / Général / Singulier / Particulier]
	\textit{cf} les 4 définitions ci-dessous
\end{rep}
\begin{dfn}[Universel]
Ce qui est valable pour tous, ou ce qui peut l’être.
\end{dfn}
\begin{dfn}[Général]
Ce qui est valable dans la majorité des cas, le plus fréquent, la norme, mais qui n’est pas valable pour tous.
\end{dfn}
\begin{dfn}[Singulier]
Ce qui est valable pour une seule personne, un seul individu. Ce qui est unique à une entité.
\end{dfn}
\begin{dfn}[Particulier]
Exception(s) du général. Ce qui est propre à un individu ou à une partie minoritaire des individus.
\end{dfn}


La raison nous ordonné d'obéir aux actions qui sont universalisables. On ne doit donc jamais mentir pour Kant.

Dans la morale kantienne, la raison pratique est à l'origine de l'action puisque c'est elle qui indique quoi faire et suit ses propres lois. C'est la différence entre l'autonomie et l'hétéronomie.
Le fait de suivre un impératif catégorique conserve notre liberté puisqu'on suit la raison.
L'objectif de la morale kantienne est de fonder la morale sur l'autonomie de la raison : ainsi l'obligation morale n'est pas en contradiction avec la liberté, au contraire.


On ne peut pas aller vérifier dans la tête d'une personne si son intention est bonne ou mauvaise, on ne peut donc jamais savoir si l'action est effectuée par devoir ou par intérêt. On ne peut pas voir les "principes intérieurs" car ils sont invisibles :
"Nous ne pouvons jamais pénétrer les mobiles secrets".

Une intention pure peut-elle réellement exister et ne pas être pervertie (même inconsciemment) ?


\begin{rep}[Absolu / Relatif]

Est dit absolu ce qui est valable sous tous les rapports, ce qui est valable en tout lieu en tout temps, en toutes circonstances. Ce qui est absolu est complet en lui-même et n'a pas besoin d'un autre terme pour exister. Ce qui ne dépend de rien.

Est dit relatif ce qui est valable uniquement sous certains rapports, en un certain lieu, en un certain temps où en certaines circonstances.
\end{rep}



\chapter{La notion du Langage}

\begin{dfn}[Langage]
Système de signaux, il sert à exprimer la pensée et à communiquer
\end{dfn}

\begin{dfn}[Système] 
Ensemble structuré d'éléments qui n'ont de sens qu'en relation les uns avec les autres. Ces éléments sont définis par les relations qu'ils entretiennent entre eux. Les éléments sont dits interdépendants et se soutiennent mutuellement.
\end{dfn}

\begin{dfn}[Signal]
Élément convenu dans le but de transmettre une information ou une instruction afin de déclencher une réaction chez le récepteur du signal.
\end{dfn}

\begin{dfn}[Langue]
Instrument de communication propre à une communauté spécifique. Ensemble de signaux verbaux écrits et/ou verbaux propres à une société.
\end{dfn}

\begin{dfn}[Parole(s)]
Acte individuel par lequel s'exprime le langage. Faculté par laquelle s'actualise le langage dans chaque individu.
Suites de mots qui expriment la parole.
\end{dfn}


\section{Les spécificités du langage humain}

cf. I A Sqce 3 Lettre à Morus du 5 février 1649 Descartes
Le langage est le signe de la conscience humaine.
Les humains ont un langage spécifique.

\subsection{Conditions anatomiques}
Cordes vocales
Position debout de l'Homme montre sa capacité de projection de la voix.
Ces deux éléments permettent à l'Homme d'avoir un langage articulé, i.e. décomposable en éléments dont la composition est réglée (syntaxe, orthographe, grammaire…).

Phonèmes
Monèmes
Unité minimale de son
Unité minimale de langage ayant du sens → souvent liée à l'étymologie 
Voyelles, consonnes et leurs sons
Chaque unité à une signification distincte
Alphabet phonétique
Un mot peut être composé de 1 ou plusieurs monèmes
A|th|é|i|s|me
A|thé|isme


\subsection{Le signe}

La naissance du signe correspond à la naissance de la linguistique.
Il s'agit d'une science créée au XIXe par Ferdinand de Saussure (1857-1913) dans son Cours de linguistique générale (1916, posthume).

\begin{dfn}[Linguistique]
Science qui étudie la langue et essaie de découvrir les lois de la langue
\end{dfn}

\begin{dfn}[Signe]
Qqch qui renvoie à autre chose que lui-même. La présence du signe permet de se tourner vers le signifié. Les mots sont des signes linguistiques, et cet élément est nommé le signifiant.
\end{dfn}



"Le signe linguistique est arbitraire", il unit arbitrairement le signifiant et le signifié (des "éléments psychiques "). Le signifiant ne renvoie pas naturellement au signifié (la preuve, plusieurs langues existent).
Ferdinand de Saussure prend l'exemple du pot "soeur" : «l'idée de soeur n'est liée par aucun rapport intérieur avec la suite de mots s-ö-r, elle pourrait être aussi bien représentée par n'importe quelle autre»

La signification est psychique (l'idée de sœur) alors que la désignation du référent ("ma sœur") est réelle.

Par exemple, les onomatopées sont différentes selon les langues (voir compilations de "Meow/Miaou" en différentes langues).

Pour établir le langage, il faut absolument une convention.

Le signe selon Ferdinand de Saussure (schéma) :


\section{À quoi sert le langage ?}

La communication est la fonction première du langage, mais il y a d'autres fonctions.

Nous allons voir quelles sont les 6 fonctions du langage selon Jakobson dans Essai de linguistique générale :
Référentielle : transmettre une information → information
Phatique : permet de prendre contact avec le récepteur → prise de contact ("Allô", "ça va", "Dites", "Bonjour" et autres salutations…)
Conative : faire agir le récepteur → réaction (impératif, "ô", "debout")
Poétique : ce qui intervient quand la valeur sonore, musicale, rythmique ou visuelle (le signifiant grossièrement) devient aussi importante voire plus importante que le signifié → art, pub dans une moindre mesure (rimes, figures poétiques, jeux de mots, jeux sonores…)
Expressive : permet à l'émetteur d'exprimer une émotion (interjections, plaintes…) → psychanalyse
Métalinguistique : capacité due langage à se questionner (sur) lui-même, ce qui permet aux interlocuteurs de vérifier qu'ils utilisent le même code (définitions, usage des guillemets autour d'un mot…)

\subsection{Le lien entre pensée et langage}

Qui est le premier du langage ou de la pensée ?
Le langage est-il une traduction de la pensée ou bien la pensée peut-elle exister sans langage ?
Émotions inexprimables sans pouvoir mettre des mots ou des signes dessus, aucun mot ne correspond à ce que je ressens → Mais est-il vraiment possible de penser sans langage → cf. Txt 2 du corpus pour la dissert sur le langage (Hegel)

\subsection{Le rôle du langage dans la culture}

Le langage est une activité sociale et possède un caractère historique. On ne peut concevoir de langage sans apprentissage, et donc sans une quelconque éducation : le langage s'inscrit dans un contexte culturel.

On pourrait dire que chaque langue révèle une façon particulière qu'ont les Hommes de concevoir le réel (vision du monde), avec notamment les différentes traductions.

\subsection{Le langage comme un moyen d'expression}
cf. Cours sur l'art (dans un futur proche)

Fonction expressive dissociable de la fonction de communication → Extériorisation des émotions (cf. Cours sur l'inconscient)

L'expression de la langue est indissociable de la parole (sauf exception), mais l'art permet aussi d'exprimer ses émotions.


\chapter{La notion du Travail}

Introduction :

Étymologie :
Le mot "travail" viendrait du latin :
"tripalium" : un instrument de torture qui empêchait les esclaves de fuir et les forçait à travailler (formé de 3 pales de bois autour du cou) ;
et de "tripaliare" : littéralement torturer.
Il faut toutefois ne pas hésiter à nuancer cette étymologie car elle n'est pas complètement attestée et fait un peu "cliché" de l'élève de terminale sur la notion du travail.



\begin{dfn}[Travail]
	Attitude généralement opposée au loisir ayant une valeur économique, impliquant l'idée que l'Homme est producteur et que par cette activité il entretient des rapports sociaux avec autrui
\end{dfn}


\begin{dfn}[Travail]
	Temps dont on dispose pour faire une activité plaisante, s'oppose au temps pendant lequel on travaille.
Les loisirs sont faits pour se ressourcer et se distraire.
\end{dfn}


On distingue deux types de travail :
le travail manuel : travail où des compétences ou capacités physiques sont principalement nécessaires ;
et le travail intellectuel : travail où les compétences intellectuelles sont plus valorisées.

\section{Le travail comme le propre de l'humanité}

\begin{refe}
	Bataille, \textit{Les larmes d'Éros}
\end{refe}

Le travail apparaît ici comme étant le fondement de l'humanité car il s'agit d'un développement pratique des capacités latentes, en germe chez l'Homme qui ne travaille pas.
→ cf. Hegel, Esthétique, Bac Blanc 2

Pour Bataille, le travail serait la condition sine qua non de l'éloignement à l'animalité première de l'Homme. Sans travail, l'Homme resterait un animal et ne pourrait se détacher de sa condition "initiale" ; sans le travail, l'animal qu'est l'Homme ne serait pas humain à proprement parler.
→ cf. conditions physiologiques du langage (cours sur le Langage)
→ et les conditions physiques ?

La bipédie permet d'avoir des mains dégagées, ce qui permet à l'Homme de travailler.

Pour Bataille, le travail est "le fondement de la connaissance et de la raison".
La connaissance et le savoir sont des acquis. Par exemple, construire une cabane permet de comprendre comment faut-il faire pour construire une cabane qui tienne debout : c'est la connaissance pratique. L'apprentissage devrait donc être empirique puisque le travail implique la connaissance.
Le travail oblige l'Homme à utiliser sa raison.

\begin{rep}[En acte/En puissance]
Est dit en acte ce qui est effectivement, ce qui est réel, ce qui existe dans la réalité.
Est dit en puissance ce qui tend vers un acte mais qui n'est pas encore effectif, ce qui n'est donc que virtuel ou possible. La puissance est en ce sens la possibilité ou encore la potentialité.
\end{rep}

Par exemple, Timéo est un pianiste en puissance tandis que je suis un pianiste en acte.

La faculté de raisonner est en puissance chez l'Homme et elle devient en acte grâce au travail.
Une fois réalisé, l'objet technique de l'artisan ou l'œuvre d'art de l'artiste est en acte.
Le passage de la puissance à l'acte est l'actualisation, aussi appelée réalisation.

Cette distinction acte/puissance provient de 
Aristote, Métaphysique
"L'acte est le fait pour une chose d'exister en réalité et non de la façon dont nous disons qu'elle existe en puissance."
"Nous appelons savant en puissance celui qui même ne spécule pas s'il a la faculté de spéculer."

Bataille dit que la fabrication d'armes et d'outils "fut le point de départ des premiers raisonnements qui humanisèrent l'animal que nous étions". En effet, l'outil est le médiat entre le travail effectué et celui qui travaille. C'est ce qui différencie le travail humain du travail animal.

Outil : instrument fabriqué par l'Homme pour modifier la matière

Grâce à ses mains, l'Homme peut créer des outils et se distinguer des animaux.

L'Homme est capable, en assimilant une certaine fin aux objets qu'il construit, de façonner le monde qui l'entoure.

La transformation faite par le travail est double : il modifie l'environnement de l'Homme et l'Homme.


\section{L'aliénation dans le travail}

En quoi un travail peut-il être dit non-libre ou aliéné ?

Caractéristiques de l'aliénation dans le travail :
Pas de liberté de création, pas d'innovation ;
Pas de liberté de mouvement (dirigé par une machine par exemple) ;
Juste pour la survie ;
Pas de décision ;
Mauvaises conditions de travail (va avec le reste) ;
Pas de pause, pas de congé.

Marx a théorisé l'aliénation au travail : pour lui, le critère de l'aliénation est l'extériorité.

Extériorité :
le travail s’impose de l’extérieur et l’objet du travail est extérieur à la finalité du travail ;
l’activité n’est pas organisée par le travailleur mais s’impose à lui de l’extérieur ;
le produit du travail est extérieur au travailleur (pas de sentiment de paternité) et peut même se retourner contre eux (caissiers et machines).

Dans le travail aliéné, l’ouvrier n’est qu’un moyen au service d’une fin. L’ouvrier est dépossédé de ce qui le constitue au profit d’une autre instance qui l’asservit (cf. définition d’aliénation). L’ouvrier devient alors étranger à son travail ;  on peut parler de dépossession psychologique et économique.

Cette description de l’aliénation est effectuée par Marx dans les textes 1 à 3 du corpus.

Ceux qui asservissent les travailleurs aliénés sont les bourgeois et les capitalistes, ce qui leur permet de réaliser du profit (augmentation du capital) et de gagner du temps de loisir.
=> cf. textes 1,6 et 7 du corpus sur Marx

D’après Marx, l’aliénation au travail est propre à la société capitaliste.

Une autre caractéristique de l’aliénation au travail est la division du travail : 

\begin{dfn}[Division du travail]
Organisation des tâches au travail propre au capitalisme. La division du travail assigne à chacun une sphère d’activité propre dont il ne peut sortir. Chaque ouvrier est spécialisé dans un domaine particulier.
\end{dfn}

ATTENTION : Marx critique l'aliénation dans le travail mais pas le travail lui-même. Il ne remet pas en question la valeur même du travail. Il dit même que le travail est le propre de l'Homme. 

\section{La remise en cause de la valeur travail}

Nietzsche a remarqué qu’avec la modernité, le travail est de plus en plus valorisé. A l'inverse, tout ce qui n’est pas un travail est dévalorisé. 

Dans d’autres civilisations, le travail n’était pas toujours valorisé. Chez les Grecs Anciens, il n’y a pas de traduction exacte du mot “travail”, seuls les esclaves travaillaient.

On peut donc remettre en caus la valeur du travail => cf. texte de Sénèque dans le corpus (s’oppose au texte de Kant)

\begin{refe}
	Nietzsche, \textit{Le Gai Savoir}, IV, paragraphe 329
\end{refe}

Il y a eu une inversion des valeurs, une transvaluation entre d’anciennes époques et l’époque de Nietzsche, qu’il juge illégitime. On cherche sans cesse à travailler sans se reposer et on néglige les “otia” (pluriel de “otium”).

“otium” (latin) = “scholê” (grec) = loisir intellectuel

Le fait de chercher à travailler toujours plus fait qu'on n’a plus de force pour les loisirs.

\chapter{La notion du Bonheur}

Introduction :

À l'origine, le bonheur désigne la bonne fortune, le bon augure ou la chance.
Le mot "bonheur" vient des mots "bon" (bon, positif) et "heur" (chance) → Porte-bonheur et "au petit bonheur la chance"



\begin{dfn}[Bonheur]
	État de satisfaction complète chez l'Homme quand ses aspirations sont comblées
On peut graduer le bonheur en trois étapes :
Le plaisir - La jouissance
Le contentement
La béatitude
\end{dfn}



\begin{dfn}[Plaisir]
	État affectif agréable que procure la satisfaction d'un désir ou d'un besoin. L'accomplissement non-entravé d'une activité désirée.
\end{dfn}


\begin{dfn}[Contentement]
	Sentiment d'être satisfait de soi-même, d'avoir accompli son devoir. Se débarrasser de la mauvaise conscience et de la culpabilité. Procure une paix intérieure.
\end{dfn}


\begin{dfn}[Béatitude]
	Sentiment de plénitude, de sérénité parfaite. A souvent une connotation religieuse. Dans la perspective chrétienne, c’est le bonheur éternel atteint par l’Homme au paradis.
\end{dfn}


\section{Qu’est-ce que le bonheur ? Le bonheur est le propre de l’Homme, tout Homme chercher à être heureux}

\begin{refe}
	Aristote, \textit{Éthique à Nicomaque}, Livre I
\end{refe}

Le bonheur ne peut pas être un moyen, il est nécessairement la fin dernière de toute action. Le bonheur est la fin commune à tous les Hommes, c'est le “Bien suprême” ou le “Souverain Bien”.

Le bonheur se suffit à lui-même.

\section{Comment être heureux ? Le bonheur ne peut reposer que sur les choses qui dépendent de nous}

\begin{refe}
	Le stoïcisme / les stoïciens
\end{refe}

Comment faire pour que notre bonheur dépende entièrement de nous ?

Selon les stoïciens, pour être heureux, nous ne pouvons compter que sur les choses qui dépendent de nous, et il faut accepter les choses qui ne dépendent pas de nous de façon stoïque, avec résignation.

Étude du texte de Sénèque :

\begin{refe}
	Sénèque, \textit{De la vie heureuse}
\end{refe}

Paragraphe 1 : Comment être heureux ?
→ La sagesse et la vertu permettent de s’accorder avec la nature
→ Ne pas compter sur la fortune et le hasard
→ Tranquillité
→ Liberté

Paragraphe 2 : Un contre-exemple à ne pas imiter : Comment être malheureux ?
→ Une vie qui n’est pas sage
→ Compter sur la fortune
→ Inquiétude, agitation
→ Esclavage

Il y a donc une opposition entre le bonheur et le plaisir chez les stoïciens.



\begin{refe}
	Sénèque, \textit{De la tranquillité de l’âme}
\end{refe}
Un des personnages se nomme Serenus (à rapprocher du mot “sérénité”).

“Toute existence est esclavage, il faut donc nous faire à notre condition, nous en plaindre le moins possible et profiter de tous les avantages qu’elle peut nous offrir.”

“L’habitude adoucit les maux.”

Il faut se moquer du destin, le rendre risible et ridicule plutôt que de s’en lamenter :
“Faisons en sorte que les vices des Hommes ne nous paraissent pas odieux mais risibles”

\section{Le bonheur est-il atteignable ? L’être humain est condamné à souffrir}

\begin{refe}
	Schopenhauer, \textit{Le monde comme volonté et comme représentation}, Livre 4
\end{refe}



L’Homme souffre donc perpétuellement sauf à l’instant même où il réalise son désir, ce après quoi il souffre de nouveau.

Il faudrait arrêter de désirer pour être heureux, or, cela est impossible puisque c’est dans la nature des Hommes de désirer. Schopenhauer appelle la tendance naturelle des Hommes à désirer la “volonté”.

Il faut donc se détacher de tout et devenir un sage/un ermite afin d'arrêter de désirer ; c’est la seule manière d’atteindre le bonheur.

\begin{refe}
	Rousseau, \textit{Julie ou la Nouvelle Héloïse}, Partie Vi, Lettre VIII (lettre de Madame de Wolmar)
\end{refe}

Rousseau transforme le cercle vicieux de Schopenhauer en cercle vertueux :
“Malheur à qui n’a plus rien à désirer !”

\chapter{Sujets types et comment y répondre} 
\begin{tabular}{|c|c|}
	Mots interrogatifs du sujet&
	Aspects à questionner\\
	Y a-t-il&
	Existence, de fait ou de droit. Ce qui existe mais aussi ce qui est possible.\\
	est-il&
	Essence. Questionner une définition, en montrant limites, autres defs et ce qui peut convenir\\
	est-ce&
	Identité, équivalence, cause et effet\\
	peut-on&
	possibilité et non-contradiction, capacité, légitimité et permission\\
	permet-il&
	pouvoirs et limites\\
	suffit-il&
	conditions nécessaire et suffisante\\
	faut-il … pour ….&
	déterminer si X est une condition suffisante à Y et si X est le moyen le + efficace ou moralement + acceptable pour Y\\
	faut-il&
	obligation\\
	doit-on&
	obligation, nécessité, contrainte et devoir\\
	pourquoi&
	objectif, cause, raison, sens\\
\end{tabular}








\chapter{Méthode}
Voir fiches distribuées sur la dissertation et l’explication de texte


\backmatter
% bibliography, glossary and index would go here.

\end{document}